\documentclass[11pt,a4paper,openany]{book}
\usepackage{geometry}
\geometry{left=20mm,right=20mm,top=20mm,bottom=22mm}
\setlength{\headheight}{22pt}
\setlength{\headsep}{5mm}
\usepackage[utf8]{inputenc}
\usepackage[T2A]{fontenc}
\usepackage[russian]{babel}
\usepackage{dejavu}
\usepackage{amsmath,amssymb,amsthm,mathtools}
\usepackage{microtype}
\usepackage{booktabs,tabularx,longtable,array,multicol,enumitem,needspace,lastpage}
\usepackage{xcolor,graphicx}
\usepackage[most]{tcolorbox}
\usepackage{tikz}
\usetikzlibrary{arrows.meta,positioning,shapes.geometric,calc,decorations.pathreplacing,decorations.pathmorphing}
\definecolor{mbaccent}{HTML}{486780}
\definecolor{mbgold}{HTML}{66798B}
\definecolor{mbpurple}{HTML}{486780}
\definecolor{mbpale}{HTML}{F8FAFB}
\definecolor{mbcream}{HTML}{FFFFFF}
\definecolor{mbink}{HTML}{243447}
\definecolor{mbmuted}{HTML}{66798B}
\definecolor{mbline}{HTML}{BAC9D4}
\usepackage[colorlinks=true,linkcolor=mbaccent,urlcolor=mbaccent,pdfencoding=auto]{hyperref}
\usepackage{bookmark}
\hypersetup{pdftitle={Что такое математическое доказательство? Истинность утверждений},pdfauthor={Лёвин Артём Александрович},pdfsubject={Пособие для самостоятельного изучения олимпиадной математики},pdfkeywords={математическое доказательство, логика, олимпиада, Math Blast}}
\usepackage{fancyhdr}
\makeatletter\renewcommand{\@pnumwidth}{2.5em}\renewcommand{\@tocrmarg}{3.5em}\makeatother
\pagestyle{fancy}\fancyhf{}
\fancyhead[L]{\small\color{mbmuted}Математическая логика}
\fancyhead[R]{\small\color{mbmuted}Доказательства и истинность}
\fancyfoot[C]{\small\color{mbmuted}\thepage\,/\,\pageref{LastPage}}
\renewcommand{\headrulewidth}{.35pt}
\renewcommand{\headrule}{\hbox to\headwidth{\color{mbline}\leaders\hrule height \headrulewidth\hfill}}
\setlength{\parindent}{0pt}\setlength{\parskip}{.43em}
\linespread{1.21}\setlength{\jot}{7pt}
\setcounter{secnumdepth}{2}\setcounter{tocdepth}{2}
\setlist[itemize]{leftmargin=*,topsep=4pt,itemsep=3pt}
\setlist[enumerate]{label=\arabic*),leftmargin=*,topsep=5pt,itemsep=4pt}
\newtcolorbox{keybox}[1][]{enhanced,breakable,colback=white,colframe=white,
  borderline west={.7pt}{0pt}{mbaccent},arc=0pt,boxrule=0pt,left=3.5mm,right=3mm,top=1.5mm,bottom=1.5mm,#1}
\newtcolorbox{learnbox}[1][]{enhanced,breakable,colback=white,colframe=white,
  borderline west={.5pt}{0pt}{mbline},arc=0pt,boxrule=0pt,left=3.5mm,right=3mm,top=1.5mm,bottom=1.5mm,#1}
\newcommand{\mtag}[1]{\textsf{\small\color{mbaccent}\bfseries #1}}
\newcommand{\muid}[1]{\texttt{\footnotesize #1}}
\renewcommand{\arraystretch}{1.38}
\setlength{\emergencystretch}{1.2em}
\newcommand{\mathshowcase}{%
\[\begin{aligned}
(a+b)^2 &= (a+b)(a+b)\\
        &= a^2+ab+ba+b^2\\
        &= a^2+2ab+b^2.
\end{aligned}\]}
\begin{document}
\let\cleardoublepage\clearpage
\thispagestyle{empty}
\begin{center}
\vspace*{23mm}
{\small\color{mbaccent}\bfseries MATH BLAST / УЧЕБНОЕ ПОСОБИЕ}\par
\vspace{21mm}
{\fontsize{26}{34}\selectfont\bfseries Что такое математическое\\[3mm] доказательство?}\par
\vspace{6mm}
{\fontsize{17}{24}\selectfont Истинность утверждений}\par
\vspace{20mm}
\begin{tikzpicture}[x=1cm,y=1cm]
\draw[mbaccent,line width=.65pt] (-5,0) .. controls (-2.8,1.45) and (-1.4,-1.0) .. (0,.1)
  .. controls (1.5,1.2) and (3.1,-.75) .. (5,0);
\end{tikzpicture}\par
\vspace{22mm}
{\small\color{mbmuted}Олимпиадная математика с нуля / 10–11 классы}\par
\vspace{3mm}
{\small Модуль 01 · Неделя 01 · Занятие 01 / 144}\par
\vspace{3mm}
{\small 30 разобранных примеров · 90 задач с полными решениями}\par
\vspace{17mm}
{\small Ученик:\enspace\rule[-.8ex]{85mm}{.4pt}}\par
\vspace{5mm}
{\small 20 ноября 2026 года}\par
\vfill
{\small\color{mbink}Лёвин Артём Александрович эксклюзивно для\enspace\rule[-.8ex]{70mm}{.4pt}}\par
\vspace{10mm}
\end{center}
\clearpage
\chapter*{Как пользоваться книгой}\addcontentsline{toc}{chapter}{Как пользоваться книгой}
Пособие предназначено для ученика, который знаком со школьной арифметикой и простейшей алгеброй, но только начинает олимпиадную подготовку. Урок по плану длится 90 минут; объём книги позволяет вернуться к теме, закрепить язык доказательства и постепенно перейти от чтения к самостоятельному оформлению.

\begin{keybox}
\textbf{План самостоятельной работы}\par
1) Прочитайте первые главы и проговорите главные определения.\par
2) Решите диагностические задания A-001–A-008 без подсказок.\par
3) Разберите 30 примеров, сравнивая свою идею с полным разбором.\par
4) Выполните 90 упражнений по уровням A–D, сначала не раскрывая подсказок.\par
5) Сверьтесь с решениями и таблицей ответов; отметьте задачи для повторения.
\end{keybox}
\textbf{Важное соглашение.} Здесь $\mathbb N=\{1,2,3,\dots\}$, $\mathbb Z$ содержит отрицательные числа и нуль, $\mathbb R$ означает вещественные числа. Слова «доказано» и «опровергнуто» всегда относятся к точно указанной области.\par
\textbf{Проверка самостоятельности.} Если вы открыли одну подсказку, сначала вернитесь к задаче. Если увидели полное решение, закройте книгу и повторите доказательство по памяти на следующий день. Для доказательных задач правильность аргумента важнее краткости ответа.
\tableofcontents
\chapter{Паспорт урока и диагностическая карта}
\begin{keybox}
\textbf{Точное название по учебному плану:} Что такое математическое доказательство? Истинность утверждений.\par
\textbf{Учебный блок:} модуль 1 «Математический язык и первые доказательства», неделя 1, первое из двух занятий.\par
\textbf{Домашняя работа недели:} 6 заданий на истинность утверждений и 4 простых доказательства.\par
\textbf{Следующий урок:} «Контрпримеры, прямые и обратные утверждения».\par
\textbf{Ожидаемый вход:} арифметические операции, числовая прямая, основы алгебраических преобразований и понятие решения уравнения.
\end{keybox}
К окончанию занятия ученик отличает утверждение от открытой формулы; сообщает область переменной; правильно читает «для всех» и «существует»; объясняет ограничения примеров; предъявляет допустимый контрпример; пишет короткое прямое доказательство; указывает логический разрыв в некорректном рассуждении. Каждая цель проверяется самостоятельной задачей с полным образцом ответа.
\section{Диагностика до чтения}
Попробуйте решить задания A-001–A-008 из практикума. Для каждого укажите тип рассуждения: вычисление, свидетель существования, контрпример, общее доказательство или уточнение области. Это входная диагностика; отсутствие правильных ответов здесь нормально и помогает выбрать скорость изучения. Затем прочитайте первые главы и решите эти же восемь заданий вновь.
\section{Карта соответствия: навык — материал — проверка}
\begin{longtable}{p{.27\textwidth}p{.39\textwidth}p{.25\textwidth}}\toprule Навык & Основная опора & Практическая проверка\\\midrule\endhead
Высказывание и область & Главы 2–3, примеры A-001–A-003 & A-001–A-010\\
Кванторы и существование & Глава 4, примеры A-006, A-010 & A-011–A-020\\
Проверка и контрпример & Главы 5–6, примеры B-011, D-030 & A-021–A-030, B-031–B-040\\
Логическая стрелка & Главы 7–9, примеры B-012 и C-024 & B-041–B-060, C-061–C-070\\
Строгое оформление & Главы 10–13, примеры C-021–D-029 & C-071–D-090\\\bottomrule\end{longtable}
\chapter{Теория: от наблюдения к доказательству}

\section{Зачем математике доказательства}\label{theory:1}

\begin{keybox}\textbf{Главная мысль.} Верное предположение — начало работы. Доказательство объясняет, почему вывод обязателен при данных условиях.\end{keybox}

В повседневном разговоре слово «доказать» иногда означает «привести убедительные примеры». В математике требуется более сильное обоснование. Если утверждается, что некоторое свойство выполнено для \emph{всех} объектов определённого типа, то пять, сто или миллион проверок ещё не охватывают бесконечное множество. Доказательство связывает вывод с условиями и уже обоснованными фактами. Задача читателя — восстановить эту связь без догадок.


Учебный пример. Проверим равенство $n+(n+1)=2n+1$ при $n=1,2,3$. Каждый опыт успешен. Затем заменим левую часть: $n+(n+1)=n+n+1=2n+1$. Последняя цепочка работает при каждом целом $n$ благодаря определению сложения. Именно она даёт общее доказательство. Конкретные вычисления полезны для поиска закономерности и поиска ошибки; логический переход к общему выводу обеспечивается цепочкой равенств.


Критерий качества: в любой строке ученик должен ответить на вопрос «почему эту строку можно написать?». Слова «очевидно», «ясно» и «по рисунку видно» могут скрывать пропущенное звено. На первом занятии мы учимся замечать такие места и дописывать обоснование, даже если само утверждение очень простое.


У математического доказательства есть адресат. Для начинающего полезно проговаривать определения и элементарные правила; в олимпиадной работе стандартные факты допускается упоминать кратко. Уровень детализации зависит от читателя, однако логические основания должны сохраняться.


\begin{center}
\begin{tikzpicture}[>=Stealth,font=\small,node distance=1.0cm]
\node[draw=mbaccent,rounded corners=3pt,fill=mbpale,minimum width=3.7cm,minimum height=0.9cm] (a) {Наблюдения: 1, 2, 3, \dots};
\node[draw=mbgold,rounded corners=3pt,fill=mbcream,right=1.3cm of a,minimum width=4.5cm,minimum height=0.9cm] (b) {Общий механизм: $n+n+1=2n+1$};
\draw[->,thick] (a)--node[above,font=\scriptsize]{ищем правило}(b);
\node[below=5mm of a,text width=3.7cm,align=center,font=\scriptsize]{Примеры помогают заметить закономерность};
\node[below=5mm of b,text width=4.5cm,align=center,font=\scriptsize]{Цепочка равенств доказывает её для всех $n$};
\end{tikzpicture}
\end{center}

\begin{learnbox}\textbf{Проверьте понимание.} \muid{MB-W01-L1-TH-01}\par Объясните своими словами, какая часть проверки $1+2=3$, $2+3=5$, $3+4=7$ является наблюдением, а какая формула даёт доказательство для любого $n$.\end{learnbox}

\Needspace{9\baselineskip}
\subsection*{Мини-тренировка: три коротких шага}
\begin{enumerate}[label=\arabic*),leftmargin=*,itemsep=3pt]

\item Назовите тип записи $4+5=9$. \textit{(Истинное числовое высказывание.)}

\item Доказывает ли проверка $n=1,2$ формулу для всех целых $n$? \textit{(Нет, проверены два случая.)}

\item Как обосновать $n+0=n$ для произвольного $n$? \textit{(Свойством нуля при сложении.)}

\end{enumerate}

\medskip

\section{Объект, область и высказывание}\label{theory:2}

\begin{keybox}\textbf{Главная мысль.} Истинность имеет смысл после того, как понятно, о каких объектах идёт речь.\end{keybox}

Высказывание — повествовательное предложение, которому при заданной интерпретации можно приписать значение «истина» или «ложь». Фраза «семь больше пяти» является высказыванием. Запись $x+3=10$ без выбора значения $x$ и без квантора — условие с переменной, то есть \emph{открытая формула}. При $x=7$ она истинна; при $x=4$ ложна. Нельзя объявлять всю открытую формулу истинной или ложной, не уточнив смысл.


\textbf{Область рассмотрения} — набор допустимых значений. Например, фраза «для всякого $x$ имеем $x^2\ge x$» ложна над вещественными числами: $x=\tfrac12$ даёт $\tfrac14<\tfrac12$. Она истинна над целыми числами, потому что целое $x$ либо $x\le0$, либо $x\ge1$, и в обоих случаях $x(x-1)\ge0$. Одинаковые символы здесь описывают разные утверждения.


Обозначения: $\mathbb N=\{1,2,3,\ldots\}$ — положительные целые, $\mathbb Z$ — все целые, $\mathbb R$ — вещественные числа. Нуль в нашем пособии не входит в $\mathbb N$; если требуется множество с нулём, пишем $\mathbb Z_{\ge0}$. Указание области в начале доказательства избавляет от множества скрытых ошибок.


Слово «произвольный» означает: в доказательстве выбран объект без специальных свойств, кроме перечисленных в условии. Пример «возьмём $x=2$» разрешён для доказательства существования, но обычно бесполезен для доказательства утверждения обо всех вещественных $x$.


\begin{center}
\begin{tikzpicture}[font=\small]
\draw[mbaccent,very thick] (0,0) ellipse (4.3cm and 1.5cm);
\draw[mbgold,very thick] (-.8,0) ellipse (2.9cm and 1.15cm);
\draw[mbpurple,very thick] (-1.8,0) ellipse (1.45cm and .72cm);
\node at (-2,0) {$\mathbb N$};\node at (.0,.2) {$\mathbb Z$};\node at (3,.5) {$\mathbb R$};
\node[below] at (0,-1.75) {Одна формула — разные множества допустимых значений};
\end{tikzpicture}
\end{center}

\begin{learnbox}\textbf{Проверьте понимание.} \muid{MB-W01-L1-TH-02}\par Для каждого выражения $x^2=9$, «все квадраты неотрицательны», $2+2=4$ определите, является ли оно высказыванием в записанном виде.\end{learnbox}

\Needspace{9\baselineskip}
\subsection*{Мини-тренировка: три коротких шага}
\begin{enumerate}[label=\arabic*),leftmargin=*,itemsep=3pt]

\item Является ли запись $x+1=3$ замкнутым высказыванием? \textit{(Нет, переменная свободна.)}

\item Что означает $x\in\mathbb Z$? \textit{(Переменная принимает целые значения.)}

\item Верно ли $x^2\ge x$ для всех $x\in\mathbb R$? \textit{(Нет, $x=\tfrac12$ — контрпример.)}

\end{enumerate}

\medskip

\section{Кванторы «для всех» и «существует»}\label{theory:3}

\begin{keybox}\textbf{Главная мысль.} От одного короткого слова иногда зависит истинность целой теоремы.\end{keybox}

Запись $\forall x\in A\ P(x)$ читается: «для каждого $x$ из $A$ выполнено свойство $P(x)$». Запись $\exists x\in A\ P(x)$ читается: «существует хотя бы один $x$ из $A$ со свойством $P(x)$». Ни один из этих символов не сообщает, \emph{почему} утверждение истинно. Он лишь фиксирует силу требования.


Для доказательства существования достаточно предъявить подходящий объект и проверить его свойства. Например, $\exists n\in\mathbb N\colon n^2=49$: берём $n=7$, поскольку $7^2=49$. Для опровержения универсального утверждения достаточно единственного допустимого контрпримера. Обратные правила требуют осторожности: десять примеров не доказывают «для всех», а одно неудачное значение не опровергает «существует».


Порядок кванторов имеет значение. Утверждение $\forall a\in\mathbb Z\ \exists b\in\mathbb Z\colon b>a$ истинно: можно положить $b=a+1$. Утверждение $\exists b\in\mathbb Z\ \forall a\in\mathbb Z\colon b>a$ ложно: никакое фиксированное целое число не превосходит все остальные. Для опровержения второго укажите $a=b+1$. Здесь важно, \emph{можно ли выбирать} $b$ после того, как стало известно $a$.


Часто вместо символов полезно сначала написать предложение по-русски. «У каждого ученика есть любимое число» допускает разное число для разных людей; «существует одно число, любимое всеми учениками» утверждает значительно больше. Перевод естественного языка в математическую запись — самостоятельный навык.


\begin{center}
\begin{tikzpicture}[>=Stealth,node distance=4mm,font=\small]
\node[draw,rounded corners,fill=mbpale,minimum width=3.5cm,minimum height=8mm] (x) {$\forall a\in\mathbb Z$};
\node[below=of x,draw,rounded corners,fill=mbcream,minimum width=3.5cm,minimum height=8mm] (y) {$\exists b=a+1$};
\node[below=of y,draw,rounded corners,fill=mbpale,minimum width=3.5cm,minimum height=8mm] (z) {$b>a$};
\draw[->] (x)--(y);\draw[->] (y)--(z);
\node[anchor=west,text width=6cm] at (2,-.55) {Сначала выбирают произвольный $a$; затем можно построить подходящий $b$.};
\end{tikzpicture}
\end{center}

\begin{learnbox}\textbf{Проверьте понимание.} \muid{MB-W01-L1-TH-03}\par Сравните $\forall n\in\mathbb N\ \exists m\in\mathbb N\colon m>n$ и $\exists m\in\mathbb N\ \forall n\in\mathbb N\colon m>n$.\end{learnbox}

\Needspace{9\baselineskip}
\subsection*{Мини-тренировка: три коротких шага}
\begin{enumerate}[label=\arabic*),leftmargin=*,itemsep=3pt]

\item Укажите натуральный свидетель $n^2=16$. \textit{($n=4$.)}

\item Истинно ли $\forall n\in\mathbb N\colon n+1>n$? \textit{(Да, прибавление $1>0$ увеличивает число.)}

\item Существует ли наибольшее целое число? \textit{(Нет: для любого $m$ число $m+1$ больше.)}

\end{enumerate}

\subsection*{Самопроверка рубежа 1}\mtag{После первых трёх глав}\par Могу отличить закрытое высказывание от открытой формулы, показать роль области и расшифровать $\forall$ и $\exists$.\par \textit{Проба:} \muid{MB-W01-L1-CP-01} Назовите контрпример к «все вещественные $x$ имеют $x^2>0$» и свидетель к «существует целое $x$ с $x^2=9$». 

\medskip

\section{Истина, ложь и недостаток данных}\label{theory:4}

\begin{keybox}\textbf{Главная мысль.} Состояние нашего знания и истинность сформулированного утверждения — разные вопросы.\end{keybox}

У полностью определённого математического высказывания есть значение истинности в выбранной области. Но мы можем пока его не знать. Запись «не знаю, верно ли» описывает состояние решающего, а не третье логическое значение в обычной двузначной математике. Сложная гипотеза не становится ложной только потому, что мы ещё не умеем её доказать.


В учебной работе удобно различать четыре ситуации: (1) верно, есть доказательство; (2) неверно, есть контрпример или вывод несовместим с условиями; (3) открытая формула, требуется уточнить переменные и кванторы; (4) формулировка полна, но обоснование пока отсутствует. Эти ситуации нельзя смешивать в итоговом ответе.


Условие задачи тоже нельзя произвольно расширять. Если дано «$x$ — натуральное», нельзя брать $x=-2$ для опровержения. Если дано «два различных числа», пример с совпадающими числами неприменим. Контрпример обязан удовлетворять \emph{всем} предпосылкам проверяемого утверждения.


Задавайте себе контрольный вопрос: что именно должно быть доказано — равенство, неравенство, существование объекта, невозможность конфигурации или единственность? Для каждой цели годятся разные свидетельства. Единичный объект доказывает существование; доказательство единственности обычно требует сравнить два произвольных предполагаемых решения.


\begin{center}
\begin{tikzpicture}[font=\small]
\foreach \x/\t/\s in {0/ИСТИНА/Есть обоснование,3.2/ЛОЖЬ/Есть контрпример,6.4/НЕ ИЗВЕСТНО/Нужен поиск}{\node[draw,rounded corners,fill=mbpale,minimum width=2.8cm,minimum height=.9cm] at (\x,0) {\t};\node[text width=2.8cm,align=center,font=\scriptsize] at (\x,-.85) {\s};}
\end{tikzpicture}
\end{center}

\begin{learnbox}\textbf{Проверьте понимание.} \muid{MB-W01-L1-TH-04}\par Приведите пример неверного универсального утверждения и объясните, почему ваш объект действительно удовлетворяет всем условиям.\end{learnbox}

\Needspace{9\baselineskip}
\subsection*{Мини-тренировка: три коротких шага}
\begin{enumerate}[label=\arabic*),leftmargin=*,itemsep=3pt]

\item Опровергает ли $x=0$ условие $x^2>0$ на $\mathbb R$? \textit{(Да, $0^2=0$.)}

\item Можно ли проверять натуральное число значением $-1$? \textit{(Нет, $-1\notin\mathbb N$.)}

\item Становится ли гипотеза ложной, если доказательство пока неизвестно? \textit{(Нет.)}

\end{enumerate}

\medskip

\section{Проверка на примерах и пределы эксперимента}\label{theory:5}

\begin{keybox}\textbf{Главная мысль.} Вычисление помогает открыть идею, но не переносит ответ автоматически на все значения.\end{keybox}

Проверяя формулу, начинайте с малых положительных чисел, нуля (если допустим), отрицательных чисел (если допустим) и граничных значений. Такой тест ловит множество опечаток и помогает увидеть роли допущений. Например, «произведение двух чисел больше каждого множителя» ломается уже при множителе $0$ или $\tfrac12$.


Предположим, что несколько испытаний дали одну закономерность. Задача исследователя — найти механизм. Для формулы $(a+b)^2=a^2+2ab+b^2$ полезно выписать значения $a,b$, затем раскрыть скобки с помощью распределительного закона. Только последний шаг объясняет, почему равенство справедливо при любом выборе вещественных чисел.


Контрпример, напротив, заканчивает опровержение универсального утверждения. Если сказано «квадрат любого вещественного числа строго положителен», число $0$ допустимо и даёт $0^2=0$, поэтому формулировка ложна. Исправление «неотрицателен» требует отдельного доказательства: $x^2\ge0$ для любого $x\in\mathbb R$.


Наблюдение за конечным количеством значений может быть доказательством, когда \emph{все возможные значения} образуют маленькое конечное множество и проверено каждое. Если $x\in\{0,1,2\}$, трёх подстановок достаточно. Если $x\in\mathbb N$, те же три подстановки не решают задачу.


\begin{center}
\begin{tikzpicture}[font=\small]
\foreach \i/\x in {0/0,1/1,2/2}{\node[draw,circle,minimum size=9mm,fill=mbpale] (p\i) at (\i*1.4,0) {$\x$};\node[below=3mm of p\i] {$\x^2$};}
\draw[decorate,decoration={brace,amplitude=5pt},thick] (-.5,.65)--(3.3,.65) node[midway,above=7pt]{Проверены все элементы конечной области};
\node[anchor=west,align=left] at (4,0) {Полный перебор: 3 случая\\Бесконечная область: требуется общий довод};
\end{tikzpicture}
\end{center}

\begin{learnbox}\textbf{Проверьте понимание.} \muid{MB-W01-L1-TH-05}\par Как отличить полный перебор от недостаточной проверки отдельных случаев?\end{learnbox}

\Needspace{9\baselineskip}
\subsection*{Мини-тренировка: три коротких шага}
\begin{enumerate}[label=\arabic*),leftmargin=*,itemsep=3pt]

\item Достаточно ли трёх примеров для универсального утверждения о $\mathbb N$? \textit{(Нет.)}

\item Достаточен ли перебор всех $x\in\{0,1\}$? \textit{(Да, если область состоит ровно из двух элементов.)}

\item Истинно ли $x^2>0$ при $x=0$? \textit{(Нет.)}

\end{enumerate}

\subsection*{Самопроверка рубежа 2}\mtag{После глав об истинности}\par Могу отличить наблюдение от доказательства, выполнить полный перебор только в конечной области и проверить допустимость контрпримера.\par \textit{Проба:} \muid{MB-W01-L1-CP-02} Почему проверка $n=1,2,3$ не доказывает свойство для всех натуральных $n$?

\medskip

\section{Каркас строгого доказательства}\label{theory:6}

\begin{keybox}\textbf{Главная мысль.} Правильный ответ состоит из допущений, оправданных переходов и явно сформулированного вывода.\end{keybox}

Полезная форма записи: «Дано …; требуется доказать …; пусть … — произвольный допустимый объект; по определению/формуле …; следовательно …; что и требовалось доказать». Такой каркас дисциплинирует мышление, но его нельзя подставлять вместо содержательных аргументов. Внутри обязательно должен быть шаг, связывающий условие с итогом.


В задаче «если $a=b$, то $a+c=b+c$» дано равенство $a=b$. Прибавление одинакового $c$ к двум частям сохраняет равенство. Отсюда $a+c=b+c$. Слова «при любом $c$» выполнены, потому что никаких ограничений на $c$ мы не вводили. Это миниатюрное прямое доказательство.


Каждое математическое равенство требует собственного основания. При переходе из $x=y$ к $x^2=y^2$ можно возвести обе части в квадрат. Обратный переход невозможен без дополнительных условий: $x^2=y^2$ допускает $x=-y$. Направление стрелки в доказательстве имеет значение. Мы учимся явно различать «следует» и «равносильно».


Финальная строка должна совпадать с тем, что спрашивали. Если требуется найти все решения, одного подходящего числа недостаточно: необходимо доказать, что других нет. Если надо доказать существование, перечисление всех вариантов может оказаться избыточным. Важно завершить мысль, а не оборвать выкладку на промежуточном результате.


\begin{center}
\begin{tikzpicture}[>=Stealth,node distance=6mm,font=\small]
\node[draw,rounded corners,fill=mbpale,minimum width=2.1cm,minimum height=9mm] (p) {Условие $P$};
\node[right=of p,draw,rounded corners,fill=mbcream,minimum width=3cm,minimum height=9mm] (a) {Правило перехода};
\node[right=of a,draw,rounded corners,fill=mbpale,minimum width=2.1cm,minimum height=9mm] (q) {Вывод $Q$};
\draw[->,very thick] (p)--(a);\draw[->,very thick] (a)--(q);
\end{tikzpicture}
\end{center}

\begin{learnbox}\textbf{Проверьте понимание.} \muid{MB-W01-L1-TH-06}\par В каком месте цепочки $a=b\Rightarrow a+c=b+c\Rightarrow 2(a+c)=2(b+c)$ используется правило сохранения равенства?\end{learnbox}

\Needspace{9\baselineskip}
\subsection*{Мини-тренировка: три коротких шага}
\begin{enumerate}[label=\arabic*),leftmargin=*,itemsep=3pt]

\item Что следует из $a=b$ после прибавления $2$? \textit{($a+2=b+2$.)}

\item Верно ли $x^2=9\Rightarrow x=3$ для вещественных $x$? \textit{(Нет, также $x=-3$.)}

\item Что должно завершать математическое доказательство? \textit{(Ровно тот вывод, который требовалось доказать.)}

\end{enumerate}

\medskip

\section{Алгебра как язык доказательства}\label{theory:7}

\begin{keybox}\textbf{Главная мысль.} Тождество — утверждение о равенстве выражений при всех допустимых значениях.\end{keybox}

Вычисление с буквами отличается от угадывания ответа по картинке. Равенство $2(a+b)=2a+2b$ следует из распределительного свойства умножения. Чтобы доказать, что $n+(n+1)+(n+2)=3n+3$, выпишем левую часть, сгруппируем однотипные слагаемые и получим правую. Мы объяснили форму выражения, а значит, доказательство работает для каждого целого $n$.


Отдельная опасность — сокращение, когда множитель может оказаться нулём. Из $ax=ay$ следует $x=y$ лишь при $a\ne0$. Если $a=0$, оба произведения равны при любых $x,y$. Правильное рассуждение сначала устанавливает условие деления, затем делит на величину с гарантированным ненулевым значением.


Для сравнения двух выражений полезно вычислить их разность. Если требуется установить $A\ge B$, можно найти $A-B$ и показать, что эта разность неотрицательна. Например, $x^2+1\ge2x$ равносильно $(x-1)^2\ge0$. Здесь последняя строка верна для всех вещественных $x$. Строка с равносильностью оправдана переносом слагаемых.


Среди допустимых инструментов первого урока — порядок арифметических действий, свойства равенства и неравенства, раскрытие скобок, формулы сокращённого умножения, свойство квадрата. Все они могут быть использованы с точным указанием области и знаков. Более сложные методы доказательства мы будем разбирать в следующих уроках.


\noindent\textbf{Образец записи вычисления по шагам.}
\mathshowcase

\begin{center}
\begin{tikzpicture}[>=Stealth,font=\small]
\node[draw,rounded corners,fill=mbpale] (x) at (0,0) {$x^2+2x+1$};
\node[draw,rounded corners,fill=mbcream] (y) at (3.6,0) {$(x+1)^2$};
\node[draw,rounded corners,fill=mbpale] (z) at (7.4,0) {$\ge0$};
\draw[->,thick] (x)--node[above,font=\scriptsize]{тождество}(y);
\draw[->,thick] (y)--node[above,font=\scriptsize]{квадрат}(z);
\end{tikzpicture}
\end{center}

\begin{learnbox}\textbf{Проверьте понимание.} \muid{MB-W01-L1-TH-07}\par Почему из $a(x-2)=0$ нельзя без уточнения сделать вывод $x=2$?\end{learnbox}

\Needspace{9\baselineskip}
\subsection*{Мини-тренировка: три коротких шага}
\begin{enumerate}[label=\arabic*),leftmargin=*,itemsep=3pt]

\item Раскройте квадрат $(a+1)^2$. \textit{($a^2+2a+1$.)}

\item Докажите $x^2+2x+1\ge0$. \textit{(Это $(x+1)^2\ge0$.)}

\item Можно ли сократить $0x=0y$ на нуль? \textit{(Нет, деление на нуль запрещено.)}

\end{enumerate}

\subsection*{Самопроверка рубежа 3}\mtag{После глав о каркасе}\par Могу записать данные, правило перехода и точный вывод; отличаю следствие от равносильности.\par \textit{Проба:} \muid{MB-W01-L1-CP-03} Почему $x^2=4\Rightarrow x=2$ неверно без дополнительного условия?

\medskip

\section{Прямые следствия и условия}\label{theory:8}

\begin{keybox}\textbf{Главная мысль.} Формулировка «если $P$, то $Q$» требует доказать переход от условия к заключению.\end{keybox}

Пусть $P$ обозначает «$x>5$», а $Q$ — «$x>3$». Когда $P$ выполняется, транзитивность порядка гарантирует $Q$. Случаи, когда $x\le5$, не нарушают импликацию: они просто не подпадают под условие. Так устроен математический смысл конструкции «если …, то …».


Обратное утверждение может быть ложным: из $x>3$ нельзя заключать $x>5$. Число $4$ доказывает это. Поэтому связь $P\Rightarrow Q$ не даёт автоматически $Q\Rightarrow P$. Эквивалентность $P\Leftrightarrow Q$ требует обоснования обоих направлений.


Правило «из ложного следует что угодно» в таблице истинности импликации иногда кажется непривычным. В задачах проще придерживаться деятельной интерпретации: надо проверить все случаи, где условие действительно выполнено; именно там нельзя получить ложное заключение. Эта договорённость связывает математическую логику с проверкой доказательств.


Представьте, что ученик пишет: «Если число кратно шести, то кратно трём. Число девять кратно трём, поэтому кратно шести». Последний шаг незаконен: использовано обратное направление. Рецензирование таких доказательств начинается с подписывания стрелок и проверки каждого основания.


\begin{center}
\begin{tikzpicture}[>=Stealth,font=\small]
\node[draw,rounded corners,fill=mbpale] (a) at (0,0) {$x>5$};
\node[draw,rounded corners,fill=mbpale] (b) at (4,0) {$x>3$};
\draw[->,very thick,mbaccent] (a)--node[above]{верно}(b);
\draw[->,dashed,mbgold] (b) to[bend left=35] node[below]{обратно — ложь при $x=4$}(a);
\end{tikzpicture}
\end{center}

\begin{learnbox}\textbf{Проверьте понимание.} \muid{MB-W01-L1-TH-08}\par Запишите утверждение, обратное к «если $x>7$, то $x>0$», и проверьте его на числе $x=1$.\end{learnbox}

\Needspace{9\baselineskip}
\subsection*{Мини-тренировка: три коротких шага}
\begin{enumerate}[label=\arabic*),leftmargin=*,itemsep=3pt]

\item Следует ли $x>5\Rightarrow x>3$? \textit{(Да, по транзитивности.)}

\item Верно ли обратное $x>3\Rightarrow x>5$? \textit{(Нет, $x=4$.)}

\item Если $a=b$, обязательно ли $a^2=b^2$? \textit{(Да, возведение в квадрат сохраняет равенство.)}

\end{enumerate}

\medskip

\section{Числовые свойства и элементарные случаи}\label{theory:9}

\begin{keybox}\textbf{Главная мысль.} Определение свойства должно быть предъявлено в виде, пригодном для вычисления.\end{keybox}

Число $m$ делится на положительное целое $d$, если существует целое $k$ такое, что $m=dk$. В первом занятии мы используем это как пример формы доказательства: чтобы обосновать кратность, надо предъявить целый множитель. Так, $12n+18=6(2n+3)$ делится на $6$ при любом $n\in\mathbb Z$, поскольку $2n+3$ — целое.


Чётность также можно выразить определением: чётное целое число имеет вид $2k$, где $k\in\mathbb Z$; нечётное — вид $2k+1$. При этом полноценная теория чётности входит в занятие №3. Здесь она служит лишь удобным материалом для объяснения, как работает квантор и почему переменная множителя должна оставаться целой.


Разбор случаев допустим, когда список исчерпывающий. Для целого $n$ ситуации $n\ge0$ и $n<0$ покрывают все возможности. Для произвольного вещественного числа при исследовании знака возможны $x<0$, $x=0$, $x>0$. Пропущенный нуль часто разрушает рассуждение, особенно когда используется строгий знак $>$.


Пусть утверждается $n^2+n\ge0$ для любого целого $n$. При $n\ge0$ оба множителя $n$ и $n+1$ неотрицательны. При $n\le-1$ оба неположительны. Следовательно, произведение $n(n+1)$ неотрицательно в каждом допустимом случае. Над вещественными это неверно при $n=-\tfrac12$; область является существенной частью теоремы.


\begin{center}
\begin{tikzpicture}[>=Stealth,font=\small]
\node[draw,rounded corners] (n) at (0,0) {$n\in\mathbb Z$};
\node[draw,rounded corners,fill=mbpale] (p) at (-2.6,-1.5) {$n\ge0$};
\node[draw,rounded corners,fill=mbcream] (m) at (2.6,-1.5) {$n\le-1$};
\node[draw,rounded corners] (pp) at (-2.6,-3) {$n(n+1)\ge0$};
\node[draw,rounded corners] (mm) at (2.6,-3) {$n(n+1)\ge0$};
\draw[->] (n)--(p);\draw[->] (n)--(m);\draw[->] (p)--(pp);\draw[->] (m)--(mm);
\end{tikzpicture}
\end{center}

\begin{learnbox}\textbf{Проверьте понимание.} \muid{MB-W01-L1-TH-09}\par Объясните, в чём состоит разница между доказательством кратности для всех $n$ и проверкой только $n=1$.\end{learnbox}

\Needspace{9\baselineskip}
\subsection*{Мини-тренировка: три коротких шага}
\begin{enumerate}[label=\arabic*),leftmargin=*,itemsep=3pt]

\item Чётно ли $n(n+1)$ для любого целого $n$? \textit{(Да, один из множителей чётный.)}

\item Вычислите $(-1)\cdot0$. \textit{($0$.)}

\item Истинно ли $x^2\ge0$ при всех $x\in\mathbb R$? \textit{(Да, квадрат неотрицателен.)}

\end{enumerate}

\subsection*{Самопроверка рубежа 4}\mtag{После алгебры и условий}\par Могу доказать простое тождество и импликацию, контролируя деление на переменные и знаки.\par \textit{Проба:} \muid{MB-W01-L1-CP-04} При каких условиях из $ax=ay$ следует $x=y$?

\medskip

\section{Рисунок, геометрия и логика}\label{theory:10}

\begin{keybox}\textbf{Главная мысль.} Чертёж подсказывает взаимосвязи, но изображённые пропорции не являются доказательством.\end{keybox}

Если изображена произвольная треугольная конфигурация, на глаз трудно удостовериться в равенстве углов или расстояний. В геометрии нужно отдельно указать, что задано, а что следует из теоремы или определения. Например, если $AB=AC$, то треугольник $ABC$ равнобедренный по определению. При дополнительных данных можно говорить о равенстве углов при основании, но их нельзя обосновывать тем, что на рисунке они выглядят одинаковыми.


Чтобы опровергнуть фразу «любой четырёхугольник — прямоугольник», достаточно предъявить допустимый четырёхугольник, например обычную трапецию с одной парой параллельных сторон и непрямым углом. На схеме зафиксируйте существенные признаки. Если важна длина, обозначьте её явно; не требуйте от читателя угадывать геометрическую истину по масштабу.


Геометрические доказательства строятся по той же схеме, что и арифметические: фиксируем условия, применяем определение или уже известную теорему, получаем вывод. Если в задаче перечислены точки, отдельно проверьте, все ли они различны и лежат ли на нужных прямых или окружностях.


В иллюстрациях пособия штриховка, подписи и графические стрелки используются для понимания логики. Рисунки с координатами показывают конкретную модель, а доказательство либо охватывает все допустимые конфигурации, либо честно заявляет, что пример служит контрпримером.


\begin{center}
\begin{tikzpicture}[scale=.85,font=\small]
\coordinate (A) at (-2,0);\coordinate (B) at (2,0);\coordinate (C) at (1,2);\coordinate (D) at (-1,2);
\draw[very thick,mbaccent] (A)--(B)--(C)--(D)--cycle;
\draw[dashed,mbgold] (A)--(C) (B)--(D);
\foreach \p/\pos in {A/below left,B/below right,C/above right,D/above left}{\fill (\p) circle(2pt) node[\pos]{$\p$};}
\node[anchor=west,text width=5cm] at (3,1) {Равные диагонали\\$AC=BD=\sqrt{13}$\\Трапеция всё же не прямоугольник};
\end{tikzpicture}
\end{center}

\begin{learnbox}\textbf{Проверьте понимание.} \muid{MB-W01-L1-TH-10}\par Почему рисунок с двумя похожими углами сам по себе не доказывает их равенство?\end{learnbox}

\Needspace{9\baselineskip}
\subsection*{Мини-тренировка: три коротких шага}
\begin{enumerate}[label=\arabic*),leftmargin=*,itemsep=3pt]

\item Равные диагонали всегда означают прямоугольник? \textit{(Нет, например, у равнобедренной трапеции они равны.)}

\item Можно ли обосновать параллельность словами «видно на рисунке»? \textit{(Нет, нужно условие или доказательство.)}

\item Может ли рисунок служить подсказкой? \textit{(Да, если вывод затем обоснован.)}

\end{enumerate}

\medskip

\section{Редактор доказательства: найдите разрыв}\label{theory:11}

\begin{keybox}\textbf{Главная мысль.} Умение распознавать ошибку ценнее быстрого согласия с правдоподобной выкладкой.\end{keybox}

Рассмотрим ошибочную запись: $a=b$, значит $a^2=ab$. Вычитая $b^2$, получаем $(a-b)(a+b)=b(a-b)$. «Сократим на $a-b$» и заключим $a+b=b$, то есть $2=1$ после подстановки $a=b=1$. Опасный шаг — сокращение на нуль, ведь $a-b=0$. До него преобразования были допустимы.


Другая ошибка: «При $n=1,2,3$ величина $n^2+n$ чётна, значит она чётна всегда». Вывод может оказаться верным, однако предложенная аргументация не доказывает его для всех натуральных $n$. Чтобы исправить запись, нужно выделить общий множитель или разобрать определение чётности, опираясь на свойства соседних целых чисел.


Третья ошибка — замена требуемого тезиса более слабым. Ученик должен доказать «$x=3$ — единственное решение», но показывает лишь $3^2=9$. Подстановка подтверждает существование одного решения уравнения $x^2=9$. Второе решение $x=-3$ опровергает слово «единственное». При редактировании полезно подчёркивать кванторы и слова «все», «только», «ровно».


Наконец, нельзя использовать то, что требуется доказать, как необоснованный промежуточный факт. Если в доказательстве встречается фраза «мы знаем, что сумма равна …», нужно выяснить, откуда именно это известно. Если источник — сама формулировка задачи, рассуждение становится круговым.


\begin{center}
\begin{tikzpicture}[>=Stealth,font=\small]
\node[draw,fill=mbpale,rounded corners,minimum width=3cm] (a) {$a=b$};
\node[draw,fill=mbcream,rounded corners,right=1cm of a,minimum width=4.4cm] (b) {$a-b=0$: делить нельзя};
\draw[->,very thick,mbgold] (a)--node[above,font=\scriptsize]{проверка}(b);
\end{tikzpicture}
\end{center}

\begin{learnbox}\textbf{Проверьте понимание.} \muid{MB-W01-L1-TH-11}\par Найдите незаконный переход в записи $(x-1)(x+2)=0\Rightarrow x-1=0$ и объясните, какой случай потерян.\end{learnbox}

\Needspace{9\baselineskip}
\subsection*{Мини-тренировка: три коротких шага}
\begin{enumerate}[label=\arabic*),leftmargin=*,itemsep=3pt]

\item Почему нельзя делить на произвольный $a$ без проверки? \textit{($a$ может быть нулём.)}

\item Сколько корней у $x^2=4$ над $\mathbb R$? \textit{(Два: $2$ и $-2$.)}

\item Что проверяют после нахождения кандидатов? \textit{(Допустимость и подстановку каждого ответа.)}

\end{enumerate}

\medskip

\section{Как оформлять олимпиадное решение}\label{theory:12}

\begin{keybox}\textbf{Главная мысль.} Письменная работа должна позволять другому человеку проверить идею и каждый существенный переход.\end{keybox}

Перед оформлением запишите цель в одной строке. Для универсального утверждения начните словами «пусть $x$ — произвольный …». Для существования укажите конструкцию объекта. Для опровержения назовите конкретный контрпример и выполните проверку условия. Для невозможности сформулируйте, какой конфликт с предпосылками получился.


Полезно отделять поиск идеи от окончательного текста. Черновик может содержать примеры и ошибочные попытки; чистовое решение должно раскрывать короткую, воспроизводимую цепь. Иногда краткая таблица исчерпывающих случаев гораздо убедительнее длинного пересказа числовых экспериментов.


В задачах на доказательство нет преимущества у огромного числа вычислений, если они не приводят к общему выводу. Проверяйте крайние случаи, область, неопределённые выражения, деление на переменные и слово «ровно». На раннем этапе достаточно честного доказательства из трёх-четырёх причинно связанных шагов.


В этом пособии каждая самостоятельная задача имеет две дозированные подсказки, полное решение и строку в таблице ответов. Используйте материалы именно в таком порядке: первая попытка без подсказок, затем одна идея, потом повторная самостоятельная работа, и лишь после этого сверка с эталоном. Если открыли решение сразу, решите аналог заново через день.


\begin{center}
\begin{tikzpicture}[font=\small,node distance=2mm]
\foreach \i/\s in {0/Область,1/Кванторы,2/Допустимость операций,3/Полнота случаев,4/Точный итог}{\node[draw,rounded corners,fill=mbpale,minimum width=5.8cm,anchor=north] (c\i) at (0,-\i*.65) {\(\square\)\quad \s};}
\end{tikzpicture}
\end{center}

\begin{learnbox}\textbf{Проверьте понимание.} \muid{MB-W01-L1-TH-12}\par Составьте мини-чек-лист из пяти вопросов к чужому доказательству; объясните выбор каждого.\end{learnbox}

\Needspace{9\baselineskip}
\subsection*{Мини-тренировка: три коротких шага}
\begin{enumerate}[label=\arabic*),leftmargin=*,itemsep=3pt]

\item Что указывают перед универсальным доказательством? \textit{(Произвольный объект и его область.)}

\item Нужно ли объяснять неочевидный переход? \textit{(Да, указать применённое правило.)}

\item Что писать в заключении? \textit{(Точную требуемую формулировку.)}

\end{enumerate}

\medskip

\section{Дальше по маршруту: мост к следующему занятию}\label{theory:13}

\begin{keybox}\textbf{Главная мысль.} Первый урок даёт язык, на котором будут доказаны более сильные утверждения.\end{keybox}

Урок №2 посвящён контрпримерам, прямым и обратным утверждениям. Сегодня уже возникли слова «для всех», «существует», «если», «обратно», «допустимый контрпример». На следующем занятии мы научимся систематически строить опровержения и уточнять условные высказывания.


Урок №3 углубит пример о чётных числах. Урок №4 объяснит доказательство от противного. В шестом уроке будет принцип Дирихле, а в седьмом — индукция. На текущем занятии эти методы нужны только как ориентиры; самостоятельные задания решаются доступными средствами: определениями, алгеброй, корректными выводами и исчерпывающими конечными случаями.


Финишная компетенция первого урока: ученик может точно прочитать высказывание, назвать область, выбрать между универсальным доказательством и контрпримером, написать короткое прямое доказательство, а также показать место логического разрыва в чужом тексте. Сложность задачи определяется глубиной рассуждения, а не размером чисел.


Сформулируйте личный критерий готовности: я умею объяснить другу, почему проверка десяти примеров не доказывает бесконечное утверждение; могу назвать все условия задачи и дать хотя бы два строгих обоснованных вывода. В начале следующей недели вернитесь к задачам, вызвавшим затруднение.


\begin{center}
\begin{tikzpicture}[>=Stealth,font=\scriptsize,node distance=4mm]
\foreach \i/\s in {0/Формулировка,1/Выбор довода,2/Строгая запись,3/Самопроверка}{\node[draw,rounded corners,fill=mbpale,minimum width=2.3cm,minimum height=.8cm] (t\i) at (\i*2.65,0) {\s};}
\foreach \i/\j in {0/1,1/2,2/3}{\draw[->,thick] (t\i)--(t\j);}
\end{tikzpicture}
\end{center}

\begin{learnbox}\textbf{Проверьте понимание.} \muid{MB-W01-L1-TH-13}\par Запишите два навыка, которые уже получаются, и два, которые требуют повторения.\end{learnbox}

\Needspace{9\baselineskip}
\subsection*{Мини-тренировка: три коротких шага}
\begin{enumerate}[label=\arabic*),leftmargin=*,itemsep=3pt]

\item Что полезно сделать после прочитанного решения? \textit{(Воспроизвести ключевое рассуждение самостоятельно.)}

\item В каком случае достаточно одного контрпримера? \textit{(Для опровержения утверждения «для всех».)}

\item Какая тема следует за этим уроком? \textit{(Контрпримеры, прямые и обратные утверждения.)}

\end{enumerate}

\subsection*{Самопроверка рубежа 5}\mtag{Перед итоговым туром}\par Могу самостоятельно написать короткое доказательство и рецензию на ложное, сохраняя область и кванторы.\par \textit{Проба:} \muid{MB-W01-L1-CP-05} Проверьте, что полное решение содержит и проверку, и объяснение отсутствия других вариантов.

\medskip

\chapter{Тридцать разобранных примеров}
\begin{keybox}
Перед чтением решения попытайтесь сделать один самостоятельный шаг. Для каждого примера в книге приведены: формулировка, план объяснения, полное решение, ловушка и вопрос на перенос. Примеры идут от A к D; обозначение \texttt{EX} сохранится в веб-версии.
\end{keybox}

\subsection*{\mtag{Разбор 01 · уровень A}\quad Истинность}\hypertarget{MB-W01-L1-EX-A-001}{}

\muid{MB-W01-L1-EX-A-001}\par

\begin{keybox}\textbf{Условие.} Верно ли числовое высказывание $7+5=12$?\end{keybox}

\textbf{Пошаговое решение.}\par
\begin{enumerate}[label=\arabic*),leftmargin=*,itemsep=3pt,topsep=4pt]
\item Сначала определяем область: это конкретные натуральные числа, свободных переменных нет.
\item Вычисляем $7+5=12$.
\item Полученное число совпадает с правой частью, значит высказывание истинно. Одна проверка достаточна здесь, потому что утверждение относится к единственному вычислению.
\end{enumerate}


\begin{learnbox}\textbf{Частая ловушка.} Нельзя переносить успешность этого вычисления на незаявленную формулу со всеми числами.\par \textbf{Перенос метода.} \muid{MB-W01-L1-TR-01} Составьте аналогичное ложное числовое высказывание.\end{learnbox}

\subsection*{\mtag{Разбор 02 · уровень A}\quad Открытая формула}\hypertarget{MB-W01-L1-EX-A-002}{}

\muid{MB-W01-L1-EX-A-002}\par

\begin{keybox}\textbf{Условие.} Можно ли объявить истинным выражение $2x+3=7$ без пояснений?\end{keybox}

\textbf{Пошаговое решение.}\par
\begin{enumerate}[label=\arabic*),leftmargin=*,itemsep=3pt,topsep=4pt]
\item Буква $x$ свободна: область и значение не заданы.
\item При $x=2$ равенство истинно, поскольку $2\cdot2+3=7$.
\item При $x=0$ получается $3\ne7$.
\item Значит, перед нами открытая формула; высказыванием она станет после задания значения или квантора.
\end{enumerate}


\begin{learnbox}\textbf{Частая ловушка.} Истинность при одном значении не равна истинности для всех.\par \textbf{Перенос метода.} \muid{MB-W01-L1-TR-02} Добавьте квантор, чтобы получить истинное утверждение.\end{learnbox}

\subsection*{\mtag{Разбор 03 · уровень A}\quad Натуральная область}\hypertarget{MB-W01-L1-EX-A-003}{}

\muid{MB-W01-L1-EX-A-003}\par

\begin{keybox}\textbf{Условие.} Докажите $3n+2>0$ для каждого $n\in\mathbb N$.\end{keybox}

\textbf{Пошаговое решение.}\par
\begin{enumerate}[label=\arabic*),leftmargin=*,itemsep=3pt,topsep=4pt]
\item По принятому определению $n\ge1$.
\item Умножение на положительное $3$ сохраняет порядок: $3n\ge3$.
\item Прибавим $2$: $3n+2\ge5>0$.
\item Так как $n$ произвольное натуральное, утверждение доказано для всех элементов области.
\end{enumerate}


\begin{learnbox}\textbf{Частая ловушка.} Без области $n\in\mathbb N$ утверждение неверно для больших по модулю отрицательных чисел.\par \textbf{Перенос метода.} \muid{MB-W01-L1-TR-03} Проверьте, остаётся ли оно истинным при $n\in\mathbb Z$.\end{learnbox}

\subsection*{\mtag{Разбор 04 · уровень A}\quad Разбор двух случаев}\hypertarget{MB-W01-L1-EX-A-004}{}

\muid{MB-W01-L1-EX-A-004}\par

\begin{keybox}\textbf{Условие.} Почему произведение двух последовательных целых чисел чётно?\end{keybox}

\textbf{Пошаговое решение.}\par
\begin{enumerate}[label=\arabic*),leftmargin=*,itemsep=3pt,topsep=4pt]
\item Пусть первое число $n\in\mathbb Z$.
\item Если $n=2k$, произведение $n(n+1)=2k(n+1)$ кратно двум.
\item Если $n=2k+1$, то $n+1=2(k+1)$ и произведение снова кратно двум.
\item Любое целое число либо чётно, либо нечётно, поэтому случаи исчерпаны.
\end{enumerate}


\begin{learnbox}\textbf{Частая ловушка.} Проверка $2\cdot3$ и $3\cdot4$ ещё не доказывает утверждение.\par \textbf{Перенос метода.} \muid{MB-W01-L1-TR-04} Где в доказательстве используется целочисленность $k$?\end{learnbox}

\subsection*{\mtag{Разбор 05 · уровень A}\quad Неотрицательность}\hypertarget{MB-W01-L1-EX-A-005}{}

\muid{MB-W01-L1-EX-A-005}\par

\begin{keybox}\textbf{Условие.} Докажите, что $x^2\ge0$ для всех $x\in\mathbb R$.\end{keybox}

\textbf{Пошаговое решение.}\par
\begin{enumerate}[label=\arabic*),leftmargin=*,itemsep=3pt,topsep=4pt]
\item Если $x\ge0$, то произведение $x\cdot x$ неотрицательно.
\item Если $x<0$, то произведение двух отрицательных чисел положительно.
\item Оба случая дают $x^2\ge0$.
\item Случай $x=0$ также включён в первый и даёт равенство.
\end{enumerate}


\begin{learnbox}\textbf{Частая ловушка.} Не заменяйте $\ge0$ на $>0$: при нуле строгое неравенство ложно.\par \textbf{Перенос метода.} \muid{MB-W01-L1-TR-05} Почему важно назвать случай равенства?\end{learnbox}

\subsection*{\mtag{Разбор 06 · уровень A}\quad Существование}\hypertarget{MB-W01-L1-EX-A-006}{}

\muid{MB-W01-L1-EX-A-006}\par

\begin{keybox}\textbf{Условие.} Докажите, что существует натуральное $n$ с $n^2=49$.\end{keybox}

\textbf{Пошаговое решение.}\par
\begin{enumerate}[label=\arabic*),leftmargin=*,itemsep=3pt,topsep=4pt]
\item Для существования достаточно одного свидетеля.
\item Выбираем $n=7$, это натуральное число.
\item Подстановка даёт $7^2=49$.
\item Следовательно, требуемое натуральное $n$ существует.
\end{enumerate}


\begin{learnbox}\textbf{Частая ловушка.} Слова «возможно существует» без предъявленного объекта слабее доказательства.\par \textbf{Перенос метода.} \muid{MB-W01-L1-TR-06} Что изменится, если требуется $n\in\mathbb Z$?\end{learnbox}

\subsection*{\mtag{Разбор 07 · уровень A}\quad Порядок}\hypertarget{MB-W01-L1-EX-A-007}{}

\muid{MB-W01-L1-EX-A-007}\par

\begin{keybox}\textbf{Условие.} Докажите $n+1>n$ при любом $n\in\mathbb Z$.\end{keybox}

\textbf{Пошаговое решение.}\par
\begin{enumerate}[label=\arabic*),leftmargin=*,itemsep=3pt,topsep=4pt]
\item Пусть $n$ — произвольное целое.
\item Для любого вещественного числа $1>0$.
\item Прибавляя $n$ к обеим частям, получаем $n+1>n$.
\item Мы нигде не выбирали специальное значение $n$, поэтому доказано универсальное высказывание.
\end{enumerate}


\begin{learnbox}\textbf{Частая ловушка.} Отрицательные значения допустимы и не требуют отдельного доказательства.\par \textbf{Перенос метода.} \muid{MB-W01-L1-TR-07} Запишите аналог для $n-1<n$.\end{learnbox}

\subsection*{\mtag{Разбор 08 · уровень A}\quad Решения уравнения}\hypertarget{MB-W01-L1-EX-A-008}{}

\muid{MB-W01-L1-EX-A-008}\par

\begin{keybox}\textbf{Условие.} Найдите все вещественные числа $x$, удовлетворяющие $x^2=4$.\end{keybox}

\textbf{Пошаговое решение.}\par
\begin{enumerate}[label=\arabic*),leftmargin=*,itemsep=3pt,topsep=4pt]
\item Переносим правую часть: $x^2-4=0$.
\item Раскладываем: $(x-2)(x+2)=0$.
\item По свойству нулевого произведения $x=2$ или $x=-2$.
\item Подставляем оба числа: квадраты равны $4$.
\item Других решений нет, поскольку произведение равно нулю только в перечисленных случаях.
\end{enumerate}


\begin{learnbox}\textbf{Частая ловушка.} Если извлечь корень и написать лишь $x=2$, теряется второе решение.\par \textbf{Перенос метода.} \muid{MB-W01-L1-TR-08} Каков ответ при ограничении $x\in\mathbb N$?\end{learnbox}

\subsection*{\mtag{Разбор 09 · уровень A}\quad Сравнение}\hypertarget{MB-W01-L1-EX-A-009}{}

\muid{MB-W01-L1-EX-A-009}\par

\begin{keybox}\textbf{Условие.} Докажите $n^2\ge n$ для любого $n\in\mathbb N$.\end{keybox}

\textbf{Пошаговое решение.}\par
\begin{enumerate}[label=\arabic*),leftmargin=*,itemsep=3pt,topsep=4pt]
\item Натуральное $n\ge1$.
\item Поэтому $n-1\ge0$.
\item Умножение на $n>0$ сохраняет знак: $n(n-1)\ge0$.
\item Раскрыв скобки, получаем $n^2-n\ge0$, то есть $n^2\ge n$.
\end{enumerate}


\begin{learnbox}\textbf{Частая ловушка.} Для вещественного $n=\tfrac12$ утверждение неверно.\par \textbf{Перенос метода.} \muid{MB-W01-L1-TR-09} Укажите все случаи равенства внутри $\mathbb N$.\end{learnbox}

\subsection*{\mtag{Разбор 10 · уровень A}\quad Несуществование}\hypertarget{MB-W01-L1-EX-A-010}{}

\muid{MB-W01-L1-EX-A-010}\par

\begin{keybox}\textbf{Условие.} Существует ли целое $x$, для которого $2x=7$?\end{keybox}

\textbf{Пошаговое решение.}\par
\begin{enumerate}[label=\arabic*),leftmargin=*,itemsep=3pt,topsep=4pt]
\item Из уравнения следует $x=\tfrac72$.
\item Число $\tfrac72$ не является целым.
\item Значит, в области $\mathbb Z$ решения не существует.
\item При смене области на $\mathbb R$ решение существует, поэтому вывод нельзя отрывать от домена.
\end{enumerate}


\begin{learnbox}\textbf{Частая ловушка.} Не объявляйте уравнение «неразрешимым вообще»: оно имеет вещественное решение.\par \textbf{Перенос метода.} \muid{MB-W01-L1-TR-10} Как изменится ответ для $2x=8$?\end{learnbox}

\subsection*{\mtag{Разбор 11 · уровень B}\quad Контрпример}\hypertarget{MB-W01-L1-EX-B-011}{}

\muid{MB-W01-L1-EX-B-011}\par

\begin{keybox}\textbf{Условие.} Опровергните фразу «всякое простое число нечётно».\end{keybox}

\textbf{Пошаговое решение.}\par
\begin{enumerate}[label=\arabic*),leftmargin=*,itemsep=3pt,topsep=4pt]
\item Фраза имеет универсальный квантор над простыми числами.
\item Рассмотрим $2$: оно простое, так как имеет ровно два положительных делителя — $1$ и $2$.
\item Число $2$ чётное.
\item Допустимый пример нарушает заключение, поэтому универсальное высказывание ложно.
\end{enumerate}


\begin{learnbox}\textbf{Частая ловушка.} Составное чётное число, например $4$, не подходит: оно не удовлетворяет условию «простое».\par \textbf{Перенос метода.} \muid{MB-W01-L1-TR-11} Переформулируйте утверждение, исключив найденное исключение.\end{learnbox}

\subsection*{\mtag{Разбор 12 · уровень B}\quad Импликация}\hypertarget{MB-W01-L1-EX-B-012}{}

\muid{MB-W01-L1-EX-B-012}\par

\begin{keybox}\textbf{Условие.} Докажите: если $x>3$, то $x>1$.\end{keybox}

\textbf{Пошаговое решение.}\par
\begin{enumerate}[label=\arabic*),leftmargin=*,itemsep=3pt,topsep=4pt]
\item Предположим $x>3$.
\item Поскольку $3>1$, из транзитивности порядка следует $x>1$.
\item Вывод получен для любого допустимого $x$ при выполненном условии.
\item Случаи $x\le3$ не требуются для доказательства условного утверждения.
\end{enumerate}


\begin{learnbox}\textbf{Частая ловушка.} Обратное неверно: $x=2$ удовлетворяет $x>1$, но нарушает $x>3$.\par \textbf{Перенос метода.} \muid{MB-W01-L1-TR-12} Придумайте другую верную цепочку из двух неравенств.\end{learnbox}

\subsection*{\mtag{Разбор 13 · уровень B}\quad Распределительность}\hypertarget{MB-W01-L1-EX-B-013}{}

\muid{MB-W01-L1-EX-B-013}\par

\begin{keybox}\textbf{Условие.} Докажите тождество $2(a+b)=2a+2b$.\end{keybox}

\textbf{Пошаговое решение.}\par
\begin{enumerate}[label=\arabic*),leftmargin=*,itemsep=3pt,topsep=4pt]
\item Пусть $a,b\in\mathbb R$ произвольны.
\item По распределительному закону умножения числа на сумму $2(a+b)=2\cdot a+2\cdot b$.
\item Правая часть совпадает с требуемым выражением.
\item Никаких дополнительных условий на $a,b$ не вводилось, поэтому тождество справедливо всегда.
\end{enumerate}


\begin{learnbox}\textbf{Частая ловушка.} Не стоит «доказывать» тождество проверкой $a=b=1$.\par \textbf{Перенос метода.} \muid{MB-W01-L1-TR-13} Как выглядит аналогичное равенство для множителя $-3$?\end{learnbox}

\subsection*{\mtag{Разбор 14 · уровень B}\quad Сохранение равенства}\hypertarget{MB-W01-L1-EX-B-014}{}

\muid{MB-W01-L1-EX-B-014}\par

\begin{keybox}\textbf{Условие.} Докажите: из $a=b$ следует $a+c=b+c$.\end{keybox}

\textbf{Пошаговое решение.}\par
\begin{enumerate}[label=\arabic*),leftmargin=*,itemsep=3pt,topsep=4pt]
\item Возьмём произвольные $a,b,c\in\mathbb R$ и допустим $a=b$.
\item По свойству равенства к обеим частям можно прибавить одно и то же число $c$.
\item Получаем $a+c=b+c$, что и требовалось.
\item Заметьте, что на $c$ нет ограничения.
\end{enumerate}


\begin{learnbox}\textbf{Частая ловушка.} Прибавление разных чисел к разным частям не сохраняет равенство.\par \textbf{Перенос метода.} \muid{MB-W01-L1-TR-14} Сформулируйте похожее следствие при умножении на $c$.\end{learnbox}

\subsection*{\mtag{Разбор 15 · уровень B}\quad Соседние числа}\hypertarget{MB-W01-L1-EX-B-015}{}

\muid{MB-W01-L1-EX-B-015}\par

\begin{keybox}\textbf{Условие.} Докажите $n(n+1)>n^2$ для $n\in\mathbb N$.\end{keybox}

\textbf{Пошаговое решение.}\par
\begin{enumerate}[label=\arabic*),leftmargin=*,itemsep=3pt,topsep=4pt]
\item Разность левой и правой частей равна $n(n+1)-n^2=n$.
\item Для натурального $n$ имеем $n>0$.
\item Поэтому разность положительна, следовательно $n(n+1)>n^2$.
\item Область существенна: при $n=0$ получилось бы равенство.
\end{enumerate}


\begin{learnbox}\textbf{Частая ловушка.} Свойство не распространяется на все целые числа.\par \textbf{Перенос метода.} \muid{MB-W01-L1-TR-15} Проверьте случай $n=-2$.\end{learnbox}

\subsection*{\mtag{Разбор 16 · уровень B}\quad Целый множитель}\hypertarget{MB-W01-L1-EX-B-016}{}

\muid{MB-W01-L1-EX-B-016}\par

\begin{keybox}\textbf{Условие.} Докажите, что $12n+18$ кратно $6$ для любого целого $n$.\end{keybox}

\textbf{Пошаговое решение.}\par
\begin{enumerate}[label=\arabic*),leftmargin=*,itemsep=3pt,topsep=4pt]
\item Выделим общий множитель: $12n+18=6(2n+3)$.
\item Число $2n+3$ целое, поскольку $n$ целое.
\item По определению делимости исходное выражение делится на $6$.
\item Это универсальный вывод, так как $n$ оставалось произвольным.
\end{enumerate}


\begin{learnbox}\textbf{Частая ловушка.} Одной записи множителя $6$ недостаточно, если другой множитель не доказан целым.\par \textbf{Перенос метода.} \muid{MB-W01-L1-TR-16} Почему при $n=\tfrac14$ такой вывод потребовал бы иной проверки?\end{learnbox}

\subsection*{\mtag{Разбор 17 · уровень B}\quad Одинаковое слагаемое}\hypertarget{MB-W01-L1-EX-B-017}{}

\muid{MB-W01-L1-EX-B-017}\par

\begin{keybox}\textbf{Условие.} Докажите: из $a>b$ следует $a+c>b+c$.\end{keybox}

\textbf{Пошаговое решение.}\par
\begin{enumerate}[label=\arabic*),leftmargin=*,itemsep=3pt,topsep=4pt]
\item Пусть $a,b,c\in\mathbb R$ и $a>b$.
\item Прибавление одинакового числа к двум сторонам строгого неравенства сохраняет порядок.
\item Следовательно, $a+c>b+c$.
\item Знак $c$ роли не играет: мы прибавляем, а не умножаем на $c$.
\end{enumerate}


\begin{learnbox}\textbf{Частая ловушка.} При умножении на отрицательное число направление знака меняется; не путайте операции.\par \textbf{Перенос метода.} \muid{MB-W01-L1-TR-17} Поясните отдельно случай $c<0$.\end{learnbox}

\subsection*{\mtag{Разбор 18 · уровень B}\quad Неотрицательные слагаемые}\hypertarget{MB-W01-L1-EX-B-018}{}

\muid{MB-W01-L1-EX-B-018}\par

\begin{keybox}\textbf{Условие.} Докажите: если $a,b\ge0$, то $a+b\ge a$.\end{keybox}

\textbf{Пошаговое решение.}\par
\begin{enumerate}[label=\arabic*),leftmargin=*,itemsep=3pt,topsep=4pt]
\item Из условия $b\ge0$.
\item Прибавляем $a$ к обеим частям: $a+b\ge a+0$.
\item Поэтому $a+b\ge a$.
\item Равенство возможно при $b=0$, это не нарушает нестрогий знак.
\end{enumerate}


\begin{learnbox}\textbf{Частая ловушка.} Без требования $b\ge0$ вывод может стать ложным.\par \textbf{Перенос метода.} \muid{MB-W01-L1-TR-18} Найдите допустимые $a,b$, для которых выполняется равенство.\end{learnbox}

\subsection*{\mtag{Разбор 19 · уровень B}\quad Квадрат нуля}\hypertarget{MB-W01-L1-EX-B-019}{}

\muid{MB-W01-L1-EX-B-019}\par

\begin{keybox}\textbf{Условие.} Докажите: если $x^2=0$, то $x=0$.\end{keybox}

\textbf{Пошаговое решение.}\par
\begin{enumerate}[label=\arabic*),leftmargin=*,itemsep=3pt,topsep=4pt]
\item Для любого вещественного $x$ число $x^2$ неотрицательно.
\item Если $x\ne0$, произведение $x\cdot x$ строго положительно.
\item Поскольку по условию квадрат равен нулю, случай $x\ne0$ исключён.
\item Следовательно, $x=0$, и подстановка подтверждает вывод.
\end{enumerate}


\begin{learnbox}\textbf{Частая ловушка.} Не используйте некорректное сокращение на $x$ без проверки $x\ne0$.\par \textbf{Перенос метода.} \muid{MB-W01-L1-TR-19} Верно ли обратное утверждение?\end{learnbox}

\subsection*{\mtag{Разбор 20 · уровень B}\quad Выражение с параметром}\hypertarget{MB-W01-L1-EX-B-020}{}

\muid{MB-W01-L1-EX-B-020}\par

\begin{keybox}\textbf{Условие.} Докажите, что сумма $n+(n+1)+(n+2)$ равна $3n+3$.\end{keybox}

\textbf{Пошаговое решение.}\par
\begin{enumerate}[label=\arabic*),leftmargin=*,itemsep=3pt,topsep=4pt]
\item Запишем сумму трёх последовательных целых чисел.
\item Раскрытие скобок даёт $n+n+1+n+2$.
\item Объединим три одинаковых слагаемых и числа: $3n+3$.
\item Тождество справедливо при любом $n\in\mathbb Z$, так как преобразование использует только законы сложения.
\end{enumerate}


\begin{learnbox}\textbf{Частая ловушка.} Нельзя выводить универсальный факт из одного примера $n=1$.\par \textbf{Перенос метода.} \muid{MB-W01-L1-TR-20} Выведите аналогичную формулу для четырёх последовательных чисел.\end{learnbox}

\subsection*{\mtag{Разбор 21 · уровень C}\quad Целые и вещественные}\hypertarget{MB-W01-L1-EX-C-021}{}

\muid{MB-W01-L1-EX-C-021}\par

\begin{keybox}\textbf{Условие.} Докажите, что $n(n+1)\ge0$ для всех $n\in\mathbb Z$, и опровергните вариант для $n\in\mathbb R$.\end{keybox}

\textbf{Пошаговое решение.}\par
\begin{enumerate}[label=\arabic*),leftmargin=*,itemsep=3pt,topsep=4pt]
\item Если целое $n\ge0$, оба множителя неотрицательны.
\item Если $n<0$, целочисленность даёт $n\le-1$, поэтому оба множителя $n$ и $n+1$ неположительны.
\item В обоих случаях произведение неотрицательно.
\item Для вещественного $n=-\tfrac12$ произведение равно $-\tfrac14<0$.
\item Область — содержательная часть утверждения.
\end{enumerate}


\begin{learnbox}\textbf{Частая ловушка.} Для отрицательного вещественного $n$ не обязательно выполняется $n\le-1$.\par \textbf{Перенос метода.} \muid{MB-W01-L1-TR-21} Сравните области, в которых утверждение верно.\end{learnbox}

\subsection*{\mtag{Разбор 22 · уровень C}\quad Двойное ограничение}\hypertarget{MB-W01-L1-EX-C-022}{}

\muid{MB-W01-L1-EX-C-022}\par

\begin{keybox}\textbf{Условие.} Докажите: из $x^2<1$ следует $-1<x<1$.\end{keybox}

\textbf{Пошаговое решение.}\par
\begin{enumerate}[label=\arabic*),leftmargin=*,itemsep=3pt,topsep=4pt]
\item Перенесём единицу: $x^2-1<0$.
\item Разложим: $(x-1)(x+1)<0$.
\item Произведение двух линейных множителей отрицательно только между корнями $-1$ и $1$.
\item Значит, $-1<x<1$.
\item Граничные значения дают квадрат $1$ и исключены строгим неравенством.
\end{enumerate}


\begin{learnbox}\textbf{Частая ловушка.} Не теряйте левую границу: запись $x<1$ слабее требования.\par \textbf{Перенос метода.} \muid{MB-W01-L1-TR-22} Докажите также обратное направление.\end{learnbox}

\subsection*{\mtag{Разбор 23 · уровень C}\quad Разность квадратов}\hypertarget{MB-W01-L1-EX-C-023}{}

\muid{MB-W01-L1-EX-C-023}\par

\begin{keybox}\textbf{Условие.} Если $a>b\ge0$, докажите $a^2>b^2$.\end{keybox}

\textbf{Пошаговое решение.}\par
\begin{enumerate}[label=\arabic*),leftmargin=*,itemsep=3pt,topsep=4pt]
\item Вычислим разность: $a^2-b^2=(a-b)(a+b)$.
\item Из $a>b$ следует $a-b>0$.
\item Из $a>b\ge0$ следует $a+b>0$.
\item Произведение положительных множителей положительно.
\item Следовательно, $a^2-b^2>0$, то есть $a^2>b^2$.
\end{enumerate}


\begin{learnbox}\textbf{Частая ловушка.} При отрицательных числах утверждение «большее число имеет больший квадрат» вообще неверно.\par \textbf{Перенос метода.} \muid{MB-W01-L1-TR-23} Испытайте $a=-1$, $b=-2$.\end{learnbox}

\subsection*{\mtag{Разбор 24 · уровень C}\quad Транзитивность}\hypertarget{MB-W01-L1-EX-C-024}{}

\muid{MB-W01-L1-EX-C-024}\par

\begin{keybox}\textbf{Условие.} Докажите: если $a\ge b$ и $b\ge c$, то $a\ge c$.\end{keybox}

\textbf{Пошаговое решение.}\par
\begin{enumerate}[label=\arabic*),leftmargin=*,itemsep=3pt,topsep=4pt]
\item Пусть предпосылки выполняются одновременно.
\item Вычтем $b$ из первого: $a-b\ge0$.
\item Вычтем $c$ из второго: $b-c\ge0$.
\item Складываем: $(a-b)+(b-c)=a-c\ge0$.
\item Значит, $a\ge c$. Это свойство транзитивности отношения порядка.
\end{enumerate}


\begin{learnbox}\textbf{Частая ловушка.} Подставлять специальные числа нельзя: доказательство относится к любым трём числам.\par \textbf{Перенос метода.} \muid{MB-W01-L1-TR-24} Какая строка меняется для строгих знаков?\end{learnbox}

\subsection*{\mtag{Разбор 25 · уровень C}\quad Квадрат разности}\hypertarget{MB-W01-L1-EX-C-025}{}

\muid{MB-W01-L1-EX-C-025}\par

\begin{keybox}\textbf{Условие.} Докажите $a+b\ge2\sqrt{ab}$ для неотрицательных $a,b$.\end{keybox}

\textbf{Пошаговое решение.}\par
\begin{enumerate}[label=\arabic*),leftmargin=*,itemsep=3pt,topsep=4pt]
\item Положим $u=\sqrt a\ge0$ и $v=\sqrt b\ge0$.
\item Квадрат $(u-v)^2\ge0$.
\item После раскрытия скобок: $a-2\sqrt{ab}+b\ge0$.
\item Переносим член и получаем $a+b\ge2\sqrt{ab}$.
\item Равенство достигается тогда и только тогда, когда $u=v$, то есть $a=b$.
\end{enumerate}


\begin{learnbox}\textbf{Частая ловушка.} Корни $\sqrt a,\sqrt b$ вещественны именно благодаря неотрицательности аргументов.\par \textbf{Перенос метода.} \muid{MB-W01-L1-TR-25} Почему нельзя применить тот же ход к $a=-1$?\end{learnbox}

\subsection*{\mtag{Разбор 26 · уровень C}\quad Остатки квадратов}\hypertarget{MB-W01-L1-EX-C-026}{}

\muid{MB-W01-L1-EX-C-026}\par

\begin{keybox}\textbf{Условие.} Докажите: квадрат целого числа при делении на $4$ имеет остаток $0$ или $1$.\end{keybox}

\textbf{Пошаговое решение.}\par
\begin{enumerate}[label=\arabic*),leftmargin=*,itemsep=3pt,topsep=4pt]
\item Любое целое $n$ имеет форму $2k$ либо $2k+1$.
\item В первом случае $n^2=4k^2$, остаток равен $0$.
\item Во втором $n^2=4k(k+1)+1$, остаток равен $1$.
\item Эти формы исчерпывают все целые числа.
\item Следовательно, иных остатков у квадратов целых нет.
\end{enumerate}


\begin{learnbox}\textbf{Частая ловушка.} Не утверждайте, что любой остаток по модулю $4$ достигается квадратом.\par \textbf{Перенос метода.} \muid{MB-W01-L1-TR-26} Как изменится список для кубов?\end{learnbox}

\subsection*{\mtag{Разбор 27 · уровень C}\quad Нулевое произведение}\hypertarget{MB-W01-L1-EX-C-027}{}

\muid{MB-W01-L1-EX-C-027}\par

\begin{keybox}\textbf{Условие.} Докажите: если $xy=0$ для вещественных $x,y$, то $x=0$ или $y=0$.\end{keybox}

\textbf{Пошаговое решение.}\par
\begin{enumerate}[label=\arabic*),leftmargin=*,itemsep=3pt,topsep=4pt]
\item Допустим, что $x\ne0$.
\item Тогда обе части $xy=0$ можно разделить на ненулевое $x$.
\item Получим $y=0$.
\item Итак, при $x=0$ уже выполнено первое условие, а при $x\ne0$ обязательно второе.
\item Случаи исчерпаны; доказано утверждение с «или».
\end{enumerate}


\begin{learnbox}\textbf{Частая ловушка.} «Или» в математике допускает, что оба множителя равны нулю.\par \textbf{Перенос метода.} \muid{MB-W01-L1-TR-27} Где точно применено требование $x\ne0$?\end{learnbox}

\subsection*{\mtag{Разбор 28 · уровень D}\quad Равносильность и стрелки}\hypertarget{MB-W01-L1-EX-D-028}{}

\muid{MB-W01-L1-EX-D-028}\par

\begin{keybox}\textbf{Условие.} Для высказываний $P,Q$ проверьте равносильность $P\Rightarrow Q$ и $\neg Q\Rightarrow\neg P$.\end{keybox}

\textbf{Пошаговое решение.}\par
\begin{enumerate}[label=\arabic*),leftmargin=*,itemsep=3pt,topsep=4pt]
\item Составим таблицу четырёх вариантов значений $(P,Q)$.
\item Импликация $P\Rightarrow Q$ ложна только при истинном $P$ и ложном $Q$.
\item Импликация $\neg Q\Rightarrow\neg P$ ложна в том же единственном случае.
\item Следовательно, столбцы истинности совпадают во всех четырёх строках.
\item Это равносильность контрапозитивных формулировок; подробнее метод будет изучен позднее.
\end{enumerate}


\begin{learnbox}\textbf{Частая ловушка.} Не спутайте контрапозицию $\neg Q\Rightarrow\neg P$ с обратной стрелкой $Q\Rightarrow P$.\par \textbf{Перенос метода.} \muid{MB-W01-L1-TR-28} Составьте собственный пример на неравенствах.\end{learnbox}

\subsection*{\mtag{Разбор 29 · уровень D}\quad Произведение трёх чисел}\hypertarget{MB-W01-L1-EX-D-029}{}

\muid{MB-W01-L1-EX-D-029}\par

\begin{keybox}\textbf{Условие.} Докажите, что $n(n+1)(n+2)$ делится на $6$ при любом $n\in\mathbb Z$.\end{keybox}

\textbf{Пошаговое решение.}\par
\begin{enumerate}[label=\arabic*),leftmargin=*,itemsep=3pt,topsep=4pt]
\item Среди трёх последовательных целых чисел хотя бы одно чётное, значит произведение кратно $2$.
\item При делении трёх последовательных целых чисел на $3$ остатки последовательно принимают все значения $0,1,2$ в некотором порядке, поэтому один из множителей кратен $3$.
\item Произведение кратно одновременно $2$ и $3$.
\item Поскольку эти числа взаимно просты, произведение кратно $6$.
\item Результат корректен также для отрицательных $n$ и нулевого множителя.
\end{enumerate}


\begin{learnbox}\textbf{Частая ловушка.} Необходимо объяснить кратность $2$ и $3$ раздельно; отдельные примеры не охватывают все $n$.\par \textbf{Перенос метода.} \muid{MB-W01-L1-TR-29} Какова самая короткая форма записи этого доказательства?\end{learnbox}

\subsection*{\mtag{Разбор 30 · уровень D}\quad Опасное обобщение}\hypertarget{MB-W01-L1-EX-D-030}{}

\muid{MB-W01-L1-EX-D-030}\par

\begin{keybox}\textbf{Условие.} Опровергните «$n^2+n+41$ — простое число для каждого $n\in\mathbb N$».\end{keybox}

\textbf{Пошаговое решение.}\par
\begin{enumerate}[label=\arabic*),leftmargin=*,itemsep=3pt,topsep=4pt]
\item Для опровержения достаточно одного натурального $n$.
\item Выберем $n=41$.
\item Тогда $n^2+n+41=41^2+41+41=41(41+2)=41\cdot43$.
\item Оба множителя больше единицы, поэтому значение составное.
\item Значит, заявленное универсальное утверждение неверно, независимо от результатов проверок для малых $n$.
\end{enumerate}


\begin{learnbox}\textbf{Частая ловушка.} Найденный контрпример обязан быть натуральным; здесь $41$ подходит.\par \textbf{Перенос метода.} \muid{MB-W01-L1-TR-30} Сколько первых подстановок может скрывать ошибочное обобщение?\end{learnbox}

\chapter{Практикум: 90 самостоятельных задач}
\begin{keybox}
\textbf{Четыре ступени.} A (30 задач): распознавание и первые доказательства. B (30): перенос и логические связи. C (20): синтез и критика доказательств. D (10): олимпиадное рассуждение.\par
Для каждого задания доступны два шага подсказки в отдельной главе. Полное решение приводится под тем же ID в другой главе. Норматив: на A — 3–7 минут, на B — 7–12, на C — 12–20, на D — 20–30; эти интервалы служат ориентиром, а не ограничением времени.
\end{keybox}

\clearpage\phantomsection\addcontentsline{toc}{section}{Алгоритм 1: Высказывание и область}

\section*{Алгоритм 1. Высказывание и область}

\vspace{18mm}\begin{center}

\begin{tikzpicture}[>=Stealth,
      flowstart/.style={ellipse,draw=mbaccent,line width=.6pt,align=center,text width=47mm,minimum height=12mm},
      flowstep/.style={rounded corners=3pt,rectangle,draw=mbline,line width=.6pt,align=center,text width=58mm,minimum height=14mm},
      flowdecision/.style={diamond,aspect=2.1,draw=mbaccent,line width=.7pt,align=center,text width=33mm,inner sep=1mm},
      flowbranch/.style={rectangle,rounded corners=3pt,draw=mbline,line width=.6pt,align=center,text width=49mm,minimum height=20mm},
      flowarrow/.style={-{Stealth[length=2.7mm]},draw=mbaccent,line width=.65pt}]

\node[flowstart] (start) at (0,0) {Проверка истинности};

\node[flowstep] (act) at (0,-2) {Определите область и смысл записи};

\node[flowdecision] (gate) at (0,-4.4) {Есть свободная переменная?};

\node[flowbranch] (yes) at (-4.0,-7.4) {Уточните значение или квантор};

\node[flowbranch] (no) at (4.0,-7.4) {Вычислите и обоснуйте результат};

\node[flowstart] (finish) at (0,-10.1) {Сформулируйте точный статус};

\draw[flowarrow] (start)--(act);
\draw[flowarrow] (act)--(gate);
\draw[flowarrow] (gate.south west)--node[above left,sloped,font=\footnotesize]{да}(yes.north);
\draw[flowarrow] (gate.south east)--node[above right,sloped,font=\footnotesize]{нет}(no.north);
\draw[flowarrow] (yes.south) |- ([xshift=-12mm]finish.north);
\draw[flowarrow] (no.south) |- ([xshift=12mm]finish.north);
\end{tikzpicture}

\end{center}\par\vspace{7mm}\noindent\textit{Что заметить:} Определите область и смысл записи. Сформулируйте точный статус.

\clearpage\section{Тренировочный блок 1 · Высказывание и область}

\textit{Десять задач: от первого применения метода к самостоятельной проверке.}\par

\Needspace{4\baselineskip}\noindent\hypertarget{MB-W01-L1-PR-A-001}{}\muid{MB-W01-L1-PR-A-001}\quad Числовая истина\par

Определите истинность высказывания $14-6=8$ и объясните ответ.\par\medskip

\Needspace{4\baselineskip}\noindent\hypertarget{MB-W01-L1-PR-A-002}{}\muid{MB-W01-L1-PR-A-002}\quad Числовая ложь\par

Верно ли $8\cdot7=54$? Укажите точную причину.\par\medskip

\Needspace{4\baselineskip}\noindent\hypertarget{MB-W01-L1-PR-A-003}{}\muid{MB-W01-L1-PR-A-003}\quad Открытая формула\par

Является ли запись $x-4=1$ высказыванием без дополнительных условий?\par\medskip

\Needspace{4\baselineskip}\noindent\hypertarget{MB-W01-L1-PR-A-004}{}\muid{MB-W01-L1-PR-A-004}\quad Область\par

Истинно ли $\forall x\in\mathbb R\colon x+2>x$?\par\medskip

\Needspace{4\baselineskip}\noindent\hypertarget{MB-W01-L1-PR-A-005}{}\muid{MB-W01-L1-PR-A-005}\quad Свидетель\par

Докажите: $\exists n\in\mathbb N\colon 3n=15$.\par\medskip

\Needspace{4\baselineskip}\noindent\hypertarget{MB-W01-L1-PR-A-006}{}\muid{MB-W01-L1-PR-A-006}\quad Несуществование\par

Существует ли натуральное $n$ с $n+2=1$?\par\medskip

\Needspace{4\baselineskip}\noindent\hypertarget{MB-W01-L1-PR-A-007}{}\muid{MB-W01-L1-PR-A-007}\quad Допустимость примера\par

Можно ли использовать $x=-3$ как контрпример к утверждению «для всех натуральных $x$ выполняется $x>0$»?\par\medskip

\Needspace{4\baselineskip}\noindent\hypertarget{MB-W01-L1-PR-A-008}{}\muid{MB-W01-L1-PR-A-008}\quad Квадраты\par

Истинно ли «каждый квадрат вещественного числа неотрицателен»?\par\medskip

\Needspace{4\baselineskip}\noindent\hypertarget{MB-W01-L1-PR-A-009}{}\muid{MB-W01-L1-PR-A-009}\quad Строгость\par

Опровергните утверждение $\forall x\in\mathbb R\colon x^2>0$.\par\medskip

\Needspace{4\baselineskip}\noindent\hypertarget{MB-W01-L1-PR-A-010}{}\muid{MB-W01-L1-PR-A-010}\quad Отрицательные числа\par

Верно ли $\forall z\in\mathbb Z\colon z+1>z$?\par\medskip

\clearpage\section*{Ответы к тренировочному блоку 1}

\addcontentsline{toc}{subsection}{Ответы к тренировочному блоку 1}

\begingroup\footnotesize\renewcommand{\arraystretch}{1.45}

\begin{tabularx}{\textwidth}{@{}p{10mm}X p{10mm}X@{}}\toprule № & Ответ / идея & № & Ответ / идея\\\midrule

1 & Истина: $14-6=8$. & 6 & Нет; единственное числовое решение $n=-1$.\\

2 & Ложь: $8\cdot7=56\ne54$. & 7 & Нельзя: $-3\notin\mathbb N$.\\

3 & Открытая формула; нужен выбор $x$ или квантор. & 8 & Истина: $x^2\ge0$ при всех $x\in\mathbb R$.\\

4 & Истина для каждого $x\in\mathbb R$. & 9 & Ложь; контрпример $x=0$.\\

5 & Свидетель $n=5$. & 10 & Истина; разность равна $1>0$.\\

\bottomrule\end{tabularx}\endgroup

\clearpage\phantomsection\addcontentsline{toc}{section}{Алгоритм 2: Существование и конечная проверка}

\section*{Алгоритм 2. Существование и конечная проверка}

\vspace{18mm}\begin{center}

\begin{tikzpicture}[>=Stealth,
      flowstart/.style={ellipse,draw=mbaccent,line width=.6pt,align=center,text width=47mm,minimum height=12mm},
      flowstep/.style={rounded corners=3pt,rectangle,draw=mbline,line width=.6pt,align=center,text width=58mm,minimum height=14mm},
      flowdecision/.style={diamond,aspect=2.1,draw=mbaccent,line width=.7pt,align=center,text width=33mm,inner sep=1mm},
      flowbranch/.style={rectangle,rounded corners=3pt,draw=mbline,line width=.6pt,align=center,text width=49mm,minimum height=20mm},
      flowarrow/.style={-{Stealth[length=2.7mm]},draw=mbaccent,line width=.65pt}]

\node[flowstart] (start) at (0,0) {Поиск свидетельства};

\node[flowstep] (act) at (0,-2) {Назовите требуемый объект};

\node[flowdecision] (gate) at (0,-4.4) {Подходит конкретный свидетель?};

\node[flowbranch] (yes) at (-4.0,-7.4) {Проверьте все условия};

\node[flowbranch] (no) at (4.0,-7.4) {Исследуйте допустимую область};

\node[flowstart] (finish) at (0,-10.1) {Докажите существование или невозможность};

\draw[flowarrow] (start)--(act);
\draw[flowarrow] (act)--(gate);
\draw[flowarrow] (gate.south west)--node[above left,sloped,font=\footnotesize]{да}(yes.north);
\draw[flowarrow] (gate.south east)--node[above right,sloped,font=\footnotesize]{нет}(no.north);
\draw[flowarrow] (yes.south) |- ([xshift=-12mm]finish.north);
\draw[flowarrow] (no.south) |- ([xshift=12mm]finish.north);
\end{tikzpicture}

\end{center}\par\vspace{7mm}\noindent\textit{Что заметить:} Назовите требуемый объект. Докажите существование или невозможность.

\clearpage\section{Тренировочный блок 2 · Существование и конечная проверка}

\textit{Десять задач: от первого применения метода к самостоятельной проверке.}\par

\Needspace{4\baselineskip}\noindent\hypertarget{MB-W01-L1-PR-A-011}{}\muid{MB-W01-L1-PR-A-011}\quad Квантор существования\par

Верно ли $\exists x\in\mathbb Z\colon x^2=16$?\par\medskip

\Needspace{4\baselineskip}\noindent\hypertarget{MB-W01-L1-PR-A-012}{}\muid{MB-W01-L1-PR-A-012}\quad Конечная область\par

Проверьте $x^2\le4$ для всех $x\in\{-2,-1,0,1,2\}$.\par\medskip

\Needspace{4\baselineskip}\noindent\hypertarget{MB-W01-L1-PR-A-013}{}\muid{MB-W01-L1-PR-A-013}\quad Обратное\par

Сформулируйте обратное к «если $x>9$, то $x>0$» и оцените его истинность.\par\medskip

\Needspace{4\baselineskip}\noindent\hypertarget{MB-W01-L1-PR-A-014}{}\muid{MB-W01-L1-PR-A-014}\quad Переменная в задаче\par

Является ли предложение «найдите $x$, если $x+7=10$» математическим высказыванием?\par\medskip

\Needspace{4\baselineskip}\noindent\hypertarget{MB-W01-L1-PR-A-015}{}\muid{MB-W01-L1-PR-A-015}\quad Неравенство\par

Докажите $2a+5>2a+4$ для любого вещественного $a$.\par\medskip

\Needspace{4\baselineskip}\noindent\hypertarget{MB-W01-L1-PR-A-016}{}\muid{MB-W01-L1-PR-A-016}\quad Вычисление и вывод\par

Оцените фразу: «$1^2+1=2$ и $2^2+2=6$, значит $n^2+n$ чётно при любом $n$».\par\medskip

\Needspace{4\baselineskip}\noindent\hypertarget{MB-W01-L1-PR-A-017}{}\muid{MB-W01-L1-PR-A-017}\quad Утверждение о нуле\par

Верно ли $\exists x\in\mathbb R\colon x^2=0$?\par\medskip

\Needspace{4\baselineskip}\noindent\hypertarget{MB-W01-L1-PR-A-018}{}\muid{MB-W01-L1-PR-A-018}\quad Положительный корень\par

Существует ли натуральное $n$, для которого $n^2=25$?\par\medskip

\Needspace{4\baselineskip}\noindent\hypertarget{MB-W01-L1-PR-A-019}{}\muid{MB-W01-L1-PR-A-019}\quad Равенство\par

Докажите тождество $a+0=a$ для всех вещественных $a$.\par\medskip

\Needspace{4\baselineskip}\noindent\hypertarget{MB-W01-L1-PR-A-020}{}\muid{MB-W01-L1-PR-A-020}\quad Контрпример по области\par

Опровергните утверждение «любой целый $n$ положителен».\par\medskip

\clearpage\section*{Ответы к тренировочному блоку 2}

\addcontentsline{toc}{subsection}{Ответы к тренировочному блоку 2}

\begingroup\footnotesize\renewcommand{\arraystretch}{1.45}

\begin{tabularx}{\textwidth}{@{}p{10mm}X p{10mm}X@{}}\toprule № & Ответ / идея & № & Ответ / идея\\\midrule

11 & Истина; например $x=4$. & 16 & Аргумент недостаточен: проверены два значения.\\

12 & Истина на указанном пятиэлементном множестве. & 17 & Истина; $x=0$.\\

13 & Обратное ложно; $x=1$. & 18 & Да; $n=5$.\\

14 & Нет; это инструкция к решению. & 19 & Истина по свойству нуля при сложении.\\

15 & Истина: разность сторон равна $1$. & 20 & Ложь; $n=0$.\\

\bottomrule\end{tabularx}\endgroup

\clearpage\phantomsection\addcontentsline{toc}{section}{Алгоритм 3: Прямое доказательство и контрпример}

\section*{Алгоритм 3. Прямое доказательство и контрпример}

\vspace{18mm}\begin{center}

\begin{tikzpicture}[>=Stealth,
      flowstart/.style={ellipse,draw=mbaccent,line width=.6pt,align=center,text width=47mm,minimum height=12mm},
      flowstep/.style={rounded corners=3pt,rectangle,draw=mbline,line width=.6pt,align=center,text width=58mm,minimum height=14mm},
      flowdecision/.style={diamond,aspect=2.1,draw=mbaccent,line width=.7pt,align=center,text width=33mm,inner sep=1mm},
      flowbranch/.style={rectangle,rounded corners=3pt,draw=mbline,line width=.6pt,align=center,text width=49mm,minimum height=20mm},
      flowarrow/.style={-{Stealth[length=2.7mm]},draw=mbaccent,line width=.65pt}]

\node[flowstart] (start) at (0,0) {Выбор способа};

\node[flowstep] (act) at (0,-2) {Выделите квантор и заключение};

\node[flowdecision] (gate) at (0,-4.4) {Требуется «для всех»?};

\node[flowbranch] (yes) at (-4.0,-7.4) {Возьмите произвольный объект};

\node[flowbranch] (no) at (4.0,-7.4) {Поищите допустимый контрпример};

\node[flowstart] (finish) at (0,-10.1) {Запишите общий аргумент или опровержение};

\draw[flowarrow] (start)--(act);
\draw[flowarrow] (act)--(gate);
\draw[flowarrow] (gate.south west)--node[above left,sloped,font=\footnotesize]{да}(yes.north);
\draw[flowarrow] (gate.south east)--node[above right,sloped,font=\footnotesize]{нет}(no.north);
\draw[flowarrow] (yes.south) |- ([xshift=-12mm]finish.north);
\draw[flowarrow] (no.south) |- ([xshift=12mm]finish.north);
\end{tikzpicture}

\end{center}\par\vspace{7mm}\noindent\textit{Что заметить:} Выделите квантор и заключение. Запишите общий аргумент или опровержение.

\clearpage\section{Тренировочный блок 3 · Прямое доказательство и контрпример}

\textit{Десять задач: от первого применения метода к самостоятельной проверке.}\par

\Needspace{4\baselineskip}\noindent\hypertarget{MB-W01-L1-PR-A-021}{}\muid{MB-W01-L1-PR-A-021}\quad Сумма\par

Покажите, что $(a+1)+(b+1)=a+b+2$ при любых вещественных $a,b$.\par\medskip

\Needspace{4\baselineskip}\noindent\hypertarget{MB-W01-L1-PR-A-022}{}\muid{MB-W01-L1-PR-A-022}\quad Неверная формула\par

Опровергните формулу $(a+b)^2=a^2+b^2$ для всех вещественных $a,b$.\par\medskip

\Needspace{4\baselineskip}\noindent\hypertarget{MB-W01-L1-PR-A-023}{}\muid{MB-W01-L1-PR-A-023}\quad Части доказательства\par

В фразе «пусть $n\in\mathbb N$; тогда $n>0$; следовательно, $n+2>2$» назовите данное, основание и вывод.\par\medskip

\Needspace{4\baselineskip}\noindent\hypertarget{MB-W01-L1-PR-A-024}{}\muid{MB-W01-L1-PR-A-024}\quad Строгое и нестрогое\par

Определите истинность $\forall x\in\mathbb R\colon x^2+1\ge1$.\par\medskip

\Needspace{4\baselineskip}\noindent\hypertarget{MB-W01-L1-PR-A-025}{}\muid{MB-W01-L1-PR-A-025}\quad Составное число\par

Опровергните «каждое натуральное число больше единицы является простым».\par\medskip

\Needspace{4\baselineskip}\noindent\hypertarget{MB-W01-L1-PR-A-026}{}\muid{MB-W01-L1-PR-A-026}\quad Свидетель и проверка\par

Докажите $\exists k\in\mathbb Z\colon k+5=0$.\par\medskip

\Needspace{4\baselineskip}\noindent\hypertarget{MB-W01-L1-PR-A-027}{}\muid{MB-W01-L1-PR-A-027}\quad Порядок величин\par

Верно ли «если $x\ge4$, то $x\ge1$»?\par\medskip

\Needspace{4\baselineskip}\noindent\hypertarget{MB-W01-L1-PR-A-028}{}\muid{MB-W01-L1-PR-A-028}\quad Ложная импликация\par

Опровергните «если целое $n$ кратно $4$, то оно кратно $8$».\par\medskip

\Needspace{4\baselineskip}\noindent\hypertarget{MB-W01-L1-PR-A-029}{}\muid{MB-W01-L1-PR-A-029}\quad Уточнение формулировки\par

Исправьте «каждое вещественное число имеет положительный квадрат», изменив одно слово.\par\medskip

\Needspace{4\baselineskip}\noindent\hypertarget{MB-W01-L1-PR-A-030}{}\muid{MB-W01-L1-PR-A-030}\quad Множество решений\par

Найдите все целые $x$, удовлетворяющие $x^2=1$.\par\medskip

\clearpage\section*{Ответы к тренировочному блоку 3}

\addcontentsline{toc}{subsection}{Ответы к тренировочному блоку 3}

\begingroup\footnotesize\renewcommand{\arraystretch}{1.45}

\begin{tabularx}{\textwidth}{@{}p{10mm}X p{10mm}X@{}}\toprule № & Ответ / идея & № & Ответ / идея\\\midrule

21 & Истина: обе части тождественны. & 26 & Истина; $k=-5$.\\

22 & Ложь; при $(1,1)$ имеем $4\ne2$. & 27 & Истина; $x\ge4\ge1$.\\

23 & Дано $n\in\mathbb N$; основание $n>0$; вывод $n+2>2$. & 28 & Ложь; $n=4$.\\

24 & Истина; равенство при $x=0$. & 29 & «Каждое вещественное число имеет неотрицательный квадрат».\\

25 & Ложь; контрпример $4$. & 30 & $x\in\{-1,1\}$.\\

\bottomrule\end{tabularx}\endgroup

\clearpage\phantomsection\addcontentsline{toc}{section}{Алгоритм 4: Алгебра и ограничения}

\section*{Алгоритм 4. Алгебра и ограничения}

\vspace{18mm}\begin{center}

\begin{tikzpicture}[>=Stealth,
      flowstart/.style={ellipse,draw=mbaccent,line width=.6pt,align=center,text width=47mm,minimum height=12mm},
      flowstep/.style={rounded corners=3pt,rectangle,draw=mbline,line width=.6pt,align=center,text width=58mm,minimum height=14mm},
      flowdecision/.style={diamond,aspect=2.1,draw=mbaccent,line width=.7pt,align=center,text width=33mm,inner sep=1mm},
      flowbranch/.style={rectangle,rounded corners=3pt,draw=mbline,line width=.6pt,align=center,text width=49mm,minimum height=20mm},
      flowarrow/.style={-{Stealth[length=2.7mm]},draw=mbaccent,line width=.65pt}]

\node[flowstart] (start) at (0,0) {Проверка переходов};

\node[flowstep] (act) at (0,-2) {Запишите предпосылки и равенства};

\node[flowdecision] (gate) at (0,-4.4) {Все операции допустимы?};

\node[flowbranch] (yes) at (-4.0,-7.4) {Выполните преобразования};

\node[flowbranch] (no) at (4.0,-7.4) {Уточните область и делитель};

\node[flowstart] (finish) at (0,-10.1) {Проверьте каждую математическую стрелку};

\draw[flowarrow] (start)--(act);
\draw[flowarrow] (act)--(gate);
\draw[flowarrow] (gate.south west)--node[above left,sloped,font=\footnotesize]{да}(yes.north);
\draw[flowarrow] (gate.south east)--node[above right,sloped,font=\footnotesize]{нет}(no.north);
\draw[flowarrow] (yes.south) |- ([xshift=-12mm]finish.north);
\draw[flowarrow] (no.south) |- ([xshift=12mm]finish.north);
\end{tikzpicture}

\end{center}\par\vspace{7mm}\noindent\textit{Что заметить:} Запишите предпосылки и равенства. Проверьте каждую математическую стрелку.

\clearpage\section{Тренировочный блок 4 · Алгебра и ограничения}

\textit{Десять задач: от первого применения метода к самостоятельной проверке.}\par

\Needspace{4\baselineskip}\noindent\hypertarget{MB-W01-L1-PR-B-031}{}\muid{MB-W01-L1-PR-B-031}\quad Прямое доказательство\par

Докажите $5(n+2)-5n=10$ для любого целого $n$.\par\medskip

\Needspace{4\baselineskip}\noindent\hypertarget{MB-W01-L1-PR-B-032}{}\muid{MB-W01-L1-PR-B-032}\quad Область истинности\par

На какой из областей $\mathbb N$, $\mathbb Z$, $\mathbb R$ истинно $\forall x\colon x^2\ge x$?\par\medskip

\Needspace{4\baselineskip}\noindent\hypertarget{MB-W01-L1-PR-B-033}{}\muid{MB-W01-L1-PR-B-033}\quad Слабый вывод\par

Если $x>8$, докажите $2x>12$.\par\medskip

\Needspace{4\baselineskip}\noindent\hypertarget{MB-W01-L1-PR-B-034}{}\muid{MB-W01-L1-PR-B-034}\quad Граница условия\par

Опровергните «если $x^2\ge4$, то $x\ge2$» над вещественными числами.\par\medskip

\Needspace{4\baselineskip}\noindent\hypertarget{MB-W01-L1-PR-B-035}{}\muid{MB-W01-L1-PR-B-035}\quad Целый множитель\par

Докажите, что $15k+25$ кратно $5$ при любом $k\in\mathbb Z$.\par\medskip

\Needspace{4\baselineskip}\noindent\hypertarget{MB-W01-L1-PR-B-036}{}\muid{MB-W01-L1-PR-B-036}\quad Неверное сокращение\par

Найдите ошибку в рассуждении: из $0\cdot a=0\cdot b$ следует $a=b$.\par\medskip

\Needspace{4\baselineskip}\noindent\hypertarget{MB-W01-L1-PR-B-037}{}\muid{MB-W01-L1-PR-B-037}\quad Идентичность\par

Докажите $(x+3)^2-(x-3)^2=12x$ для всех вещественных $x$.\par\medskip

\Needspace{4\baselineskip}\noindent\hypertarget{MB-W01-L1-PR-B-038}{}\muid{MB-W01-L1-PR-B-038}\quad Условная кратность\par

Докажите: если целое $n$ кратно $7$, то $3n$ кратно $7$.\par\medskip

\Needspace{4\baselineskip}\noindent\hypertarget{MB-W01-L1-PR-B-039}{}\muid{MB-W01-L1-PR-B-039}\quad Контрпример к квадрату\par

Опровергните «если $a^2>b^2$, то $a>b$» для вещественных $a,b$.\par\medskip

\Needspace{4\baselineskip}\noindent\hypertarget{MB-W01-L1-PR-B-040}{}\muid{MB-W01-L1-PR-B-040}\quad Существование целого\par

Найдите целое число $m$, для которого $m^2-5m=0$, и докажите существование.\par\medskip

\clearpage\section*{Ответы к тренировочному блоку 4}

\addcontentsline{toc}{subsection}{Ответы к тренировочному блоку 4}

\begingroup\footnotesize\renewcommand{\arraystretch}{1.45}

\begin{tabularx}{\textwidth}{@{}p{10mm}X p{10mm}X@{}}\toprule № & Ответ / идея & № & Ответ / идея\\\midrule

31 & Тождество; результат всегда $10$. & 36 & Ошибка: сокращение на нуль.\\

32 & Истинно в $\mathbb N,\mathbb Z$; ложно в $\mathbb R$. & 37 & Тождество: $12x$.\\

33 & Верно; $2x>16>12$. & 38 & Истина; $3n=7(3k)$.\\

34 & Ложь; $x=-2$. & 39 & Ложь; $(a,b)=(-3,2)$.\\

35 & Истина: $15k+25=5(3k+5)$. & 40 & Существуют $m=0$ и $m=5$.\\

\bottomrule\end{tabularx}\endgroup

\clearpage\phantomsection\addcontentsline{toc}{section}{Алгоритм 5: Следствие, обратное, кванторы}

\section*{Алгоритм 5. Следствие, обратное, кванторы}

\vspace{18mm}\begin{center}

\begin{tikzpicture}[>=Stealth,
      flowstart/.style={ellipse,draw=mbaccent,line width=.6pt,align=center,text width=47mm,minimum height=12mm},
      flowstep/.style={rounded corners=3pt,rectangle,draw=mbline,line width=.6pt,align=center,text width=58mm,minimum height=14mm},
      flowdecision/.style={diamond,aspect=2.1,draw=mbaccent,line width=.7pt,align=center,text width=33mm,inner sep=1mm},
      flowbranch/.style={rectangle,rounded corners=3pt,draw=mbline,line width=.6pt,align=center,text width=49mm,minimum height=20mm},
      flowarrow/.style={-{Stealth[length=2.7mm]},draw=mbaccent,line width=.65pt}]

\node[flowstart] (start) at (0,0) {Анализ логики};

\node[flowstep] (act) at (0,-2) {Разделите предпосылку и вывод};

\node[flowdecision] (gate) at (0,-4.4) {Нужно обратное утверждение?};

\node[flowbranch] (yes) at (-4.0,-7.4) {Проверьте обратную стрелку};

\node[flowbranch] (no) at (4.0,-7.4) {Докажите только прямой вывод};

\node[flowstart] (finish) at (0,-10.1) {Различайте «следует» и «равносильно»};

\draw[flowarrow] (start)--(act);
\draw[flowarrow] (act)--(gate);
\draw[flowarrow] (gate.south west)--node[above left,sloped,font=\footnotesize]{да}(yes.north);
\draw[flowarrow] (gate.south east)--node[above right,sloped,font=\footnotesize]{нет}(no.north);
\draw[flowarrow] (yes.south) |- ([xshift=-12mm]finish.north);
\draw[flowarrow] (no.south) |- ([xshift=12mm]finish.north);
\end{tikzpicture}

\end{center}\par\vspace{7mm}\noindent\textit{Что заметить:} Разделите предпосылку и вывод. Различайте «следует» и «равносильно».

\clearpage\section{Тренировочный блок 5 · Следствие, обратное, кванторы}

\textit{Десять задач: от первого применения метода к самостоятельной проверке.}\par

\Needspace{4\baselineskip}\noindent\hypertarget{MB-W01-L1-PR-B-041}{}\muid{MB-W01-L1-PR-B-041}\quad Усиление тезиса\par

Докажите, что $x^2+4\ge4$ для любого вещественного $x$, и укажите равенство.\par\medskip

\Needspace{4\baselineskip}\noindent\hypertarget{MB-W01-L1-PR-B-042}{}\muid{MB-W01-L1-PR-B-042}\quad Сумма последовательных\par

Докажите, что $k+(k+1)+(k+2)+(k+3)=4k+6$ для всех целых $k$.\par\medskip

\Needspace{4\baselineskip}\noindent\hypertarget{MB-W01-L1-PR-B-043}{}\muid{MB-W01-L1-PR-B-043}\quad Неравенство с двумя условиями\par

Если $a\ge2$ и $b\ge3$, докажите $a+b\ge5$.\par\medskip

\Needspace{4\baselineskip}\noindent\hypertarget{MB-W01-L1-PR-B-044}{}\muid{MB-W01-L1-PR-B-044}\quad Неполное доказательство\par

Оцените аргумент «$3\cdot4$ кратно $3$, значит произведение любых двух натуральных чисел кратно $3$».\par\medskip

\Needspace{4\baselineskip}\noindent\hypertarget{MB-W01-L1-PR-B-045}{}\muid{MB-W01-L1-PR-B-045}\quad Проверка обратного\par

Из $a=b$ следует $a^2=b^2$. Верно ли обратное для вещественных $a,b$?\par\medskip

\Needspace{4\baselineskip}\noindent\hypertarget{MB-W01-L1-PR-B-046}{}\muid{MB-W01-L1-PR-B-046}\quad Отрицательный множитель\par

Если $a>b$, докажите $-2a<-2b$.\par\medskip

\Needspace{4\baselineskip}\noindent\hypertarget{MB-W01-L1-PR-B-047}{}\muid{MB-W01-L1-PR-B-047}\quad Уравнение с двумя решениями\par

Найдите все вещественные $x$ с $(x-4)(x+1)=0$.\par\medskip

\Needspace{4\baselineskip}\noindent\hypertarget{MB-W01-L1-PR-B-048}{}\muid{MB-W01-L1-PR-B-048}\quad Двойная квантификация\par

Докажите $\forall n\in\mathbb N\ \exists m\in\mathbb N\colon m=n+2$.\par\medskip

\Needspace{4\baselineskip}\noindent\hypertarget{MB-W01-L1-PR-B-049}{}\muid{MB-W01-L1-PR-B-049}\quad Кванторы в обратном порядке\par

Опровергните $\exists m\in\mathbb N\ \forall n\in\mathbb N\colon m>n$.\par\medskip

\Needspace{4\baselineskip}\noindent\hypertarget{MB-W01-L1-PR-B-050}{}\muid{MB-W01-L1-PR-B-050}\quad Логика «и»\par

Опровергните «любое натуральное число чётно и кратно $3$».\par\medskip

\clearpage\section*{Ответы к тренировочному блоку 5}

\addcontentsline{toc}{subsection}{Ответы к тренировочному блоку 5}

\begingroup\footnotesize\renewcommand{\arraystretch}{1.45}

\begin{tabularx}{\textwidth}{@{}p{10mm}X p{10mm}X@{}}\toprule № & Ответ / идея & № & Ответ / идея\\\midrule

41 & Истина; равенство только при $x=0$. & 46 & Истина: $-2a<-2b$.\\

42 & Тождество; $4k+6$. & 47 & $x\in\{-1,4\}$.\\

43 & Истина; $a+b\ge5$. & 48 & Истина; $m=n+2$.\\

44 & Ложь; контрпример $2\cdot2=4$. & 49 & Ложь; для любого $m$ подходит $n=m+1$.\\

45 & Обратное ложно; $(1,-1)$. & 50 & Ложь; $n=1$.\\

\bottomrule\end{tabularx}\endgroup

\clearpage\phantomsection\addcontentsline{toc}{section}{Алгоритм 6: Полнота и случаи}

\section*{Алгоритм 6. Полнота и случаи}

\vspace{18mm}\begin{center}

\begin{tikzpicture}[>=Stealth,
      flowstart/.style={ellipse,draw=mbaccent,line width=.6pt,align=center,text width=47mm,minimum height=12mm},
      flowstep/.style={rounded corners=3pt,rectangle,draw=mbline,line width=.6pt,align=center,text width=58mm,minimum height=14mm},
      flowdecision/.style={diamond,aspect=2.1,draw=mbaccent,line width=.7pt,align=center,text width=33mm,inner sep=1mm},
      flowbranch/.style={rectangle,rounded corners=3pt,draw=mbline,line width=.6pt,align=center,text width=49mm,minimum height=20mm},
      flowarrow/.style={-{Stealth[length=2.7mm]},draw=mbaccent,line width=.65pt}]

\node[flowstart] (start) at (0,0) {Контроль доказательства};

\node[flowstep] (act) at (0,-2) {Укажите область и варианты};

\node[flowdecision] (gate) at (0,-4.4) {Все случаи перечислены?};

\node[flowbranch] (yes) at (-4.0,-7.4) {Проверьте каждый случай};

\node[flowbranch] (no) at (4.0,-7.4) {Найдите пропущенную границу};

\node[flowstart] (finish) at (0,-10.1) {Сделайте вывод для всей области};

\draw[flowarrow] (start)--(act);
\draw[flowarrow] (act)--(gate);
\draw[flowarrow] (gate.south west)--node[above left,sloped,font=\footnotesize]{да}(yes.north);
\draw[flowarrow] (gate.south east)--node[above right,sloped,font=\footnotesize]{нет}(no.north);
\draw[flowarrow] (yes.south) |- ([xshift=-12mm]finish.north);
\draw[flowarrow] (no.south) |- ([xshift=12mm]finish.north);
\end{tikzpicture}

\end{center}\par\vspace{7mm}\noindent\textit{Что заметить:} Укажите область и варианты. Сделайте вывод для всей области.

\clearpage\section{Тренировочный блок 6 · Полнота и случаи}

\textit{Десять задач: от первого применения метода к самостоятельной проверке.}\par

\Needspace{4\baselineskip}\noindent\hypertarget{MB-W01-L1-PR-B-051}{}\muid{MB-W01-L1-PR-B-051}\quad Логика «или»\par

Докажите: для любого целого $n$ верно $n\ge0$ или $n<0$.\par\medskip

\Needspace{4\baselineskip}\noindent\hypertarget{MB-W01-L1-PR-B-052}{}\muid{MB-W01-L1-PR-B-052}\quad Деление и область\par

Можно ли из $ax=ay$ заключить $x=y$, если известно $a\ne0$?\par\medskip

\Needspace{4\baselineskip}\noindent\hypertarget{MB-W01-L1-PR-B-053}{}\muid{MB-W01-L1-PR-B-053}\quad Квадрат разности\par

Докажите $x^2-2x+1\ge0$ при любом $x\in\mathbb R$.\par\medskip

\Needspace{4\baselineskip}\noindent\hypertarget{MB-W01-L1-PR-B-054}{}\muid{MB-W01-L1-PR-B-054}\quad Неверная универсальность\par

Опровергните «сумма двух положительных чисел меньше их произведения».\par\medskip

\Needspace{4\baselineskip}\noindent\hypertarget{MB-W01-L1-PR-B-055}{}\muid{MB-W01-L1-PR-B-055}\quad Связанные неравенства\par

Если $a\ge b$ и $b\ge0$, докажите $a\ge0$.\par\medskip

\Needspace{4\baselineskip}\noindent\hypertarget{MB-W01-L1-PR-B-056}{}\muid{MB-W01-L1-PR-B-056}\quad Чётный показатель\par

Докажите $x^4\ge0$ при любом вещественном $x$.\par\medskip

\Needspace{4\baselineskip}\noindent\hypertarget{MB-W01-L1-PR-B-057}{}\muid{MB-W01-L1-PR-B-057}\quad Неравенство с параметром\par

Докажите: если $t\ge1$, то $t^2\ge t$.\par\medskip

\Needspace{4\baselineskip}\noindent\hypertarget{MB-W01-L1-PR-B-058}{}\muid{MB-W01-L1-PR-B-058}\quad Доказательство единственности\par

Докажите, что уравнение $x+7=3$ имеет ровно одно вещественное решение.\par\medskip

\Needspace{4\baselineskip}\noindent\hypertarget{MB-W01-L1-PR-B-059}{}\muid{MB-W01-L1-PR-B-059}\quad Круговое рассуждение\par

Почему «$a+b=b+a$, потому что сумма не меняется при перестановке слагаемых» допустимо, если правило уже изучено, но недостаточно для доказательства самого правила?\par\medskip

\Needspace{4\baselineskip}\noindent\hypertarget{MB-W01-L1-PR-B-060}{}\muid{MB-W01-L1-PR-B-060}\quad Проверка краёв\par

Покажите, почему утверждение «если $x>0$, то $x^2>0$» истинно, а замена условия на $x\ge0$ делает его ложным.\par\medskip

\clearpage\section*{Ответы к тренировочному блоку 6}

\addcontentsline{toc}{subsection}{Ответы к тренировочному блоку 6}

\begingroup\footnotesize\renewcommand{\arraystretch}{1.45}

\begin{tabularx}{\textwidth}{@{}p{10mm}X p{10mm}X@{}}\toprule № & Ответ / идея & № & Ответ / идея\\\midrule

51 & Истина; два случая исчерпывают $\mathbb Z$. & 56 & Истина; $x^4=(x^2)^2\ge0$.\\

52 & Да; делим на ненулевое $a$. & 57 & Истина при $t\ge1$.\\

53 & Истина; $(x-1)^2\ge0$. & 58 & Единственное решение $x=-4$.\\

54 & Ложь; $1+1=2>1\cdot1$. & 59 & Допустимо как ссылка на известный закон; кругово как его доказательство.\\

55 & Истина; $a\ge b\ge0$. & 60 & Первая импликация истинна; вторая ложна при $x=0$.\\

\bottomrule\end{tabularx}\endgroup

\clearpage\phantomsection\addcontentsline{toc}{section}{Алгоритм 7: Кванторы и синтез}

\section*{Алгоритм 7. Кванторы и синтез}

\vspace{18mm}\begin{center}

\begin{tikzpicture}[>=Stealth,
      flowstart/.style={ellipse,draw=mbaccent,line width=.6pt,align=center,text width=47mm,minimum height=12mm},
      flowstep/.style={rounded corners=3pt,rectangle,draw=mbline,line width=.6pt,align=center,text width=58mm,minimum height=14mm},
      flowdecision/.style={diamond,aspect=2.1,draw=mbaccent,line width=.7pt,align=center,text width=33mm,inner sep=1mm},
      flowbranch/.style={rectangle,rounded corners=3pt,draw=mbline,line width=.6pt,align=center,text width=49mm,minimum height=20mm},
      flowarrow/.style={-{Stealth[length=2.7mm]},draw=mbaccent,line width=.65pt}]

\node[flowstart] (start) at (0,0) {Управление зависимостями};

\node[flowstep] (act) at (0,-2) {Запишите кванторы в нужном порядке};

\node[flowdecision] (gate) at (0,-4.4) {Свидетель зависит от переменной?};

\node[flowbranch] (yes) at (-4.0,-7.4) {Постройте свидетель после выбора};

\node[flowbranch] (no) at (4.0,-7.4) {Проверьте постоянный кандидат};

\node[flowstart] (finish) at (0,-10.1) {Обоснуйте итог с правильными кванторами};

\draw[flowarrow] (start)--(act);
\draw[flowarrow] (act)--(gate);
\draw[flowarrow] (gate.south west)--node[above left,sloped,font=\footnotesize]{да}(yes.north);
\draw[flowarrow] (gate.south east)--node[above right,sloped,font=\footnotesize]{нет}(no.north);
\draw[flowarrow] (yes.south) |- ([xshift=-12mm]finish.north);
\draw[flowarrow] (no.south) |- ([xshift=12mm]finish.north);
\end{tikzpicture}

\end{center}\par\vspace{7mm}\noindent\textit{Что заметить:} Запишите кванторы в нужном порядке. Обоснуйте итог с правильными кванторами.

\clearpage\section{Тренировочный блок 7 · Кванторы и синтез}

\textit{Десять задач: от первого применения метода к самостоятельной проверке.}\par

\Needspace{4\baselineskip}\noindent\hypertarget{MB-W01-L1-PR-C-061}{}\muid{MB-W01-L1-PR-C-061}\quad Смена области\par

Докажите $n(n-2)\ge0$ для целых $n\le0$ или $n\ge2$; найдите целое исключение для всех $\mathbb Z$.\par\medskip

\Needspace{4\baselineskip}\noindent\hypertarget{MB-W01-L1-PR-C-062}{}\muid{MB-W01-L1-PR-C-062}\quad Квадратная оценка\par

Докажите $x^2+6x+10\ge1$ для каждого вещественного $x$ и найдите случай равенства.\par\medskip

\Needspace{4\baselineskip}\noindent\hypertarget{MB-W01-L1-PR-C-063}{}\muid{MB-W01-L1-PR-C-063}\quad Ошибка в направлении\par

Рецензируйте «если $n$ кратно $12$, то кратно $4$; $20$ кратно $4$; значит, $20$ кратно $12$».\par\medskip

\Needspace{4\baselineskip}\noindent\hypertarget{MB-W01-L1-PR-C-064}{}\muid{MB-W01-L1-PR-C-064}\quad Три части множества\par

Для $x\in\mathbb R$ докажите, что ровно одна ситуация $x<-2$, $x=-2$, $x>-2$ реализуется.\par\medskip

\Needspace{4\baselineskip}\noindent\hypertarget{MB-W01-L1-PR-C-065}{}\muid{MB-W01-L1-PR-C-065}\quad Равенство двух квадратов\par

Найдите все пары вещественных $a,b$, для которых $a^2=b^2$, в виде условия на $a,b$.\par\medskip

\Needspace{4\baselineskip}\noindent\hypertarget{MB-W01-L1-PR-C-066}{}\muid{MB-W01-L1-PR-C-066}\quad Двойной контрпример\par

Опровергните «если два вещественных числа имеют одинаковые модули, то они равны»; запишите верное исправление.\par\medskip

\Needspace{4\baselineskip}\noindent\hypertarget{MB-W01-L1-PR-C-067}{}\muid{MB-W01-L1-PR-C-067}\quad Квантор и зависимость\par

Докажите $\forall a\in\mathbb Z\ \exists b\in\mathbb Z\colon a+b=0$ и объясните, почему свидетель зависит от $a$.\par\medskip

\Needspace{4\baselineskip}\noindent\hypertarget{MB-W01-L1-PR-C-068}{}\muid{MB-W01-L1-PR-C-068}\quad Ложная перестановка\par

Опровергните $\exists b\in\mathbb Z\ \forall a\in\mathbb Z\colon a+b=0$.\par\medskip

\Needspace{4\baselineskip}\noindent\hypertarget{MB-W01-L1-PR-C-069}{}\muid{MB-W01-L1-PR-C-069}\quad Неравенство по условию\par

Если $u>v\ge0$, докажите $u^2>v^2$ и объясните, почему $u\ge0$ следует из условия.\par\medskip

\Needspace{4\baselineskip}\noindent\hypertarget{MB-W01-L1-PR-C-070}{}\muid{MB-W01-L1-PR-C-070}\quad Эквивалентность границ\par

Докажите: для вещественного $x$ условие $|x|\le3$ равносильно $-3\le x\le3$.\par\medskip

\clearpage\section*{Ответы к тренировочному блоку 7}

\addcontentsline{toc}{subsection}{Ответы к тренировочному блоку 7}

\begingroup\footnotesize\renewcommand{\arraystretch}{1.45}

\begin{tabularx}{\textwidth}{@{}p{10mm}X p{10mm}X@{}}\toprule № & Ответ / идея & № & Ответ / идея\\\midrule

61 & Верно на $n\le0$ или $n\ge2$; исключение $n=1$. & 66 & Контрпример $(3,-3)$; исправление $a=\pm b$.\\

62 & Минимум $1$ при $x=-3$. & 67 & Истина; для данного $a$ берём $b=-a$.\\

63 & Ошибка: незаконная обратная импликация. & 68 & Ложь: $a=0$ требует $b=0$, а $a=1$ — $b=-1$.\\

64 & Три случая образуют полное разбиение. & 69 & Истина; $(u-v)(u+v)>0$.\\

65 & $a=b$ или $a=-b$. & 70 & $|x|\le3\Leftrightarrow-3\le x\le3$.\\

\bottomrule\end{tabularx}\endgroup

\clearpage\phantomsection\addcontentsline{toc}{section}{Алгоритм 8: Границы и единственность}

\section*{Алгоритм 8. Границы и единственность}

\vspace{18mm}\begin{center}

\begin{tikzpicture}[>=Stealth,
      flowstart/.style={ellipse,draw=mbaccent,line width=.6pt,align=center,text width=47mm,minimum height=12mm},
      flowstep/.style={rounded corners=3pt,rectangle,draw=mbline,line width=.6pt,align=center,text width=58mm,minimum height=14mm},
      flowdecision/.style={diamond,aspect=2.1,draw=mbaccent,line width=.7pt,align=center,text width=33mm,inner sep=1mm},
      flowbranch/.style={rectangle,rounded corners=3pt,draw=mbline,line width=.6pt,align=center,text width=49mm,minimum height=20mm},
      flowarrow/.style={-{Stealth[length=2.7mm]},draw=mbaccent,line width=.65pt}]

\node[flowstart] (start) at (0,0) {Полное решение};

\node[flowstep] (act) at (0,-2) {Получите список кандидатов};

\node[flowdecision] (gate) at (0,-4.4) {Все кандидаты допустимы?};

\node[flowbranch] (yes) at (-4.0,-7.4) {Проверьте подстановкой};

\node[flowbranch] (no) at (4.0,-7.4) {Отбросьте нарушающие условие};

\node[flowstart] (finish) at (0,-10.1) {Докажите отсутствие остальных решений};

\draw[flowarrow] (start)--(act);
\draw[flowarrow] (act)--(gate);
\draw[flowarrow] (gate.south west)--node[above left,sloped,font=\footnotesize]{да}(yes.north);
\draw[flowarrow] (gate.south east)--node[above right,sloped,font=\footnotesize]{нет}(no.north);
\draw[flowarrow] (yes.south) |- ([xshift=-12mm]finish.north);
\draw[flowarrow] (no.south) |- ([xshift=12mm]finish.north);
\end{tikzpicture}

\end{center}\par\vspace{7mm}\noindent\textit{Что заметить:} Получите список кандидатов. Докажите отсутствие остальных решений.

\clearpage\section{Тренировочный блок 8 · Границы и единственность}

\textit{Десять задач: от первого применения метода к самостоятельной проверке.}\par

\Needspace{4\baselineskip}\noindent\hypertarget{MB-W01-L1-PR-C-071}{}\muid{MB-W01-L1-PR-C-071}\quad Разбор ошибки с корнем\par

Найдите недостающее условие в выводе «из $x^2=9$ следует $x=3$».\par\medskip

\Needspace{4\baselineskip}\noindent\hypertarget{MB-W01-L1-PR-C-072}{}\muid{MB-W01-L1-PR-C-072}\quad Усреднение\par

Докажите: если $a\ge b$, то $a\ge\frac{a+b}{2}\ge b$ для вещественных $a,b$.\par\medskip

\Needspace{4\baselineskip}\noindent\hypertarget{MB-W01-L1-PR-C-073}{}\muid{MB-W01-L1-PR-C-073}\quad Единственность результата\par

Докажите, что уравнение $3x-5=7$ имеет единственное вещественное решение.\par\medskip

\Needspace{4\baselineskip}\noindent\hypertarget{MB-W01-L1-PR-C-074}{}\muid{MB-W01-L1-PR-C-074}\quad Невозможность\par

Докажите, что не существуют вещественные $x,y$ с $x^2+y^2=-1$.\par\medskip

\Needspace{4\baselineskip}\noindent\hypertarget{MB-W01-L1-PR-C-075}{}\muid{MB-W01-L1-PR-C-075}\quad Контрпример к нестрогому\par

Опровергните «если $a>b$, то $a^2\ge b^2$» без ограничений на знаки.\par\medskip

\Needspace{4\baselineskip}\noindent\hypertarget{MB-W01-L1-PR-C-076}{}\muid{MB-W01-L1-PR-C-076}\quad Разбиение условия\par

Найдите все целые $n$, при которых $n^2\le4$, и докажите полноту списка.\par\medskip

\Needspace{4\baselineskip}\noindent\hypertarget{MB-W01-L1-PR-C-077}{}\muid{MB-W01-L1-PR-C-077}\quad Сравнение долей\par

Докажите: если $a>b>0$, то $\frac1a<\frac1b$.\par\medskip

\Needspace{4\baselineskip}\noindent\hypertarget{MB-W01-L1-PR-C-078}{}\muid{MB-W01-L1-PR-C-078}\quad Строгость доказательства\par

Прочитайте «если $a=b$, то $a^3=b^3$» и напишите обоснование, применимое ко всем вещественным $a,b$.\par\medskip

\Needspace{4\baselineskip}\noindent\hypertarget{MB-W01-L1-PR-C-079}{}\muid{MB-W01-L1-PR-C-079}\quad Анализ ложного доказательства\par

Ученик пишет: «Для любого $x$ имеем $x^2+2x+1>0$, потому что это квадрат». Найдите ошибку и исправьте тезис.\par\medskip

\Needspace{4\baselineskip}\noindent\hypertarget{MB-W01-L1-PR-C-080}{}\muid{MB-W01-L1-PR-C-080}\quad Оценка по двум условиям\par

Если $0\le a\le2$ и $0\le b\le3$, докажите $0\le ab\le6$.\par\medskip

\clearpage\section*{Ответы к тренировочному блоку 8}

\addcontentsline{toc}{subsection}{Ответы к тренировочному блоку 8}

\begingroup\footnotesize\renewcommand{\arraystretch}{1.45}

\begin{tabularx}{\textwidth}{@{}p{10mm}X p{10mm}X@{}}\toprule № & Ответ / идея & № & Ответ / идея\\\midrule

71 & Нужно $x\ge0$; без него $x=\pm3$. & 76 & $n\in\{-2,-1,0,1,2\}$.\\

72 & $a\ge(a+b)/2\ge b$. & 77 & Истина: $1/a<1/b$.\\

73 & $x=4$, единственное решение. & 78 & Истина: подстановка равных величин сохраняет куб.\\

74 & Невозможно: $x^2+y^2\ge0$. & 79 & Ошибка в строгости: верно $\ge0$, равенство при $x=-1$.\\

75 & Ложь; $(-1,-2)$. & 80 & $0\le ab\le6$.\\

\bottomrule\end{tabularx}\endgroup

\clearpage\phantomsection\addcontentsline{toc}{section}{Алгоритм 9: Олимпиадное доказательство}

\section*{Алгоритм 9. Олимпиадное доказательство}

\vspace{18mm}\begin{center}

\begin{tikzpicture}[>=Stealth,
      flowstart/.style={ellipse,draw=mbaccent,line width=.6pt,align=center,text width=47mm,minimum height=12mm},
      flowstep/.style={rounded corners=3pt,rectangle,draw=mbline,line width=.6pt,align=center,text width=58mm,minimum height=14mm},
      flowdecision/.style={diamond,aspect=2.1,draw=mbaccent,line width=.7pt,align=center,text width=33mm,inner sep=1mm},
      flowbranch/.style={rectangle,rounded corners=3pt,draw=mbline,line width=.6pt,align=center,text width=49mm,minimum height=20mm},
      flowarrow/.style={-{Stealth[length=2.7mm]},draw=mbaccent,line width=.65pt}]

\node[flowstart] (start) at (0,0) {Построение идеи};

\node[flowstep] (act) at (0,-2) {Разделите цель на условия};

\node[flowdecision] (gate) at (0,-4.4) {Есть универсальный механизм?};

\node[flowbranch] (yes) at (-4.0,-7.4) {Используйте тождество или оценку};

\node[flowbranch] (no) at (4.0,-7.4) {Ищите крайний допустимый случай};

\node[flowstart] (finish) at (0,-10.1) {Сверьте необходимость, достаточность и итог};

\draw[flowarrow] (start)--(act);
\draw[flowarrow] (act)--(gate);
\draw[flowarrow] (gate.south west)--node[above left,sloped,font=\footnotesize]{да}(yes.north);
\draw[flowarrow] (gate.south east)--node[above right,sloped,font=\footnotesize]{нет}(no.north);
\draw[flowarrow] (yes.south) |- ([xshift=-12mm]finish.north);
\draw[flowarrow] (no.south) |- ([xshift=12mm]finish.north);
\end{tikzpicture}

\end{center}\par\vspace{7mm}\noindent\textit{Что заметить:} Разделите цель на условия. Сверьте необходимость, достаточность и итог.

\clearpage\section{Тренировочный блок 9 · Олимпиадное доказательство}

\textit{Десять задач: от первого применения метода к самостоятельной проверке.}\par

\Needspace{4\baselineskip}\noindent\hypertarget{MB-W01-L1-PR-D-081}{}\muid{MB-W01-L1-PR-D-081}\quad Две границы\par

Для $x\in\mathbb R$ докажите эквивалентность: $x^2-4x+3\le0$ тогда и только тогда, когда $1\le x\le3$.\par\medskip

\Needspace{4\baselineskip}\noindent\hypertarget{MB-W01-L1-PR-D-082}{}\muid{MB-W01-L1-PR-D-082}\quad Зависимость кванторов\par

Сравните истинность $\forall x\in\mathbb R\ \exists y\in\mathbb R\colon x+y=1$ и $\exists y\in\mathbb R\ \forall x\in\mathbb R\colon x+y=1$.\par\medskip

\Needspace{4\baselineskip}\noindent\hypertarget{MB-W01-L1-PR-D-083}{}\muid{MB-W01-L1-PR-D-083}\quad Несколько условий\par

Если $a,b,c\ge0$ и $a+b+c=0$, докажите $a=b=c=0$.\par\medskip

\Needspace{4\baselineskip}\noindent\hypertarget{MB-W01-L1-PR-D-084}{}\muid{MB-W01-L1-PR-D-084}\quad Разбор ложного обобщения\par

Числа $1^2+1+17$, $2^2+2+17$ и $3^2+3+17$ просты. Следует ли, что $n^2+n+17$ всегда простое при $n\in\mathbb N$?\par\medskip

\Needspace{4\baselineskip}\noindent\hypertarget{MB-W01-L1-PR-D-085}{}\muid{MB-W01-L1-PR-D-085}\quad Геометрический контрпример\par

Опровергните «если у четырёхугольника диагонали равны, то он прямоугольник» на примере равнобедренной трапеции.\par\medskip

\Needspace{4\baselineskip}\noindent\hypertarget{MB-W01-L1-PR-D-086}{}\muid{MB-W01-L1-PR-D-086}\quad Невозможность с условиями\par

Докажите, что не существует вещественного $x$, для которого одновременно $x>2$ и $x^2<4$.\par\medskip

\Needspace{4\baselineskip}\noindent\hypertarget{MB-W01-L1-PR-D-087}{}\muid{MB-W01-L1-PR-D-087}\quad Уточнение области\par

Исследуйте истинность $\forall x\in A\colon x^2\ge2x$ для $A=\mathbb N$ и для $A=\{2,3,4,\ldots\}$; приведите верную максимальную целую область.\par\medskip

\Needspace{4\baselineskip}\noindent\hypertarget{MB-W01-L1-PR-D-088}{}\muid{MB-W01-L1-PR-D-088}\quad Правка доказательства\par

Ученик получил из $ab=ac$ равенство $b=c$. Уточните необходимую предпосылку и обоснуйте её необходимость примером.\par\medskip

\Needspace{4\baselineskip}\noindent\hypertarget{MB-W01-L1-PR-D-089}{}\muid{MB-W01-L1-PR-D-089}\quad Синтез ограничений\par

Докажите, что при $a,b\in\mathbb R$ и $a+b=6$ выполняется $ab\le9$, и найдите условие равенства.\par\medskip

\Needspace{4\baselineskip}\noindent\hypertarget{MB-W01-L1-PR-D-090}{}\muid{MB-W01-L1-PR-D-090}\quad Полное доказательство\par

Найдите все вещественные решения системы $x+y=4$, $x^2+y^2=8$ и строго докажите полноту ответа.\par\medskip

\clearpage\section*{Ответы к тренировочному блоку 9}

\addcontentsline{toc}{subsection}{Ответы к тренировочному блоку 9}

\begingroup\footnotesize\renewcommand{\arraystretch}{1.45}

\begin{tabularx}{\textwidth}{@{}p{10mm}X p{10mm}X@{}}\toprule № & Ответ / идея & № & Ответ / идея\\\midrule

81 & Равносильно $x\in[1,3]$. & 86 & Невозможно: из $x>2$ следует $x^2>4$.\\

82 & Первое истинно ($y=1-x$); второе ложно. & 87 & Ложно на $\mathbb N$; истинно при $x\in\mathbb Z$, $x\le0$ или $x\ge2$.\\

83 & $a=b=c=0$. & 88 & Нужно $a\ne0$; при $a=0$, $b=1$, $c=2$ вывод неверен.\\

84 & Нет; при $n=17$ значение $17\cdot19$ составное. & 89 & $ab\le9$; равенство при $a=b=3$.\\

85 & Ложь; симметричная трапеция с диагоналями $\sqrt{13}$. & 90 & Единственная пара $(x,y)=(2,2)$.\\

\bottomrule\end{tabularx}\endgroup

\chapter{Подсказки: открывайте последовательно}
\begin{learnbox}Если после первой подсказки появилась идея, закройте раздел и решите задачу. Вторая подсказка направляет метод, но оставляет проверку условий и оформление вам. Каждый ID совпадает с условием из практикума.\end{learnbox}

\section{Подсказки уровня A}

\noindent\muid{MB-W01-L1-PR-A-001}\par 
\textbf{Шаг 1:} Здесь нет свободной переменной.\par \textbf{Шаг 2:} Вычислите левую часть и сравните с правой.\smallskip

\noindent\muid{MB-W01-L1-PR-A-002}\par 
\textbf{Шаг 1:} Вспомните таблицу умножения.\par \textbf{Шаг 2:} Действительное произведение равно $56$.\smallskip

\noindent\muid{MB-W01-L1-PR-A-003}\par 
\textbf{Шаг 1:} Есть ли в записи свободная буква?\par \textbf{Шаг 2:} Сравните подстановки $x=5$ и $x=2$.\smallskip

\noindent\muid{MB-W01-L1-PR-A-004}\par 
\textbf{Шаг 1:} Рассмотрите разность левой и правой частей.\par \textbf{Шаг 2:} Разность равна положительному числу $2$.\smallskip

\noindent\muid{MB-W01-L1-PR-A-005}\par 
\textbf{Шаг 1:} Для существования нужен один подходящий $n$.\par \textbf{Шаг 2:} Попробуйте $n=5$.\smallskip

\noindent\muid{MB-W01-L1-PR-A-006}\par 
\textbf{Шаг 1:} Решите равенство относительно $n$.\par \textbf{Шаг 2:} Полученное значение сравните с $\mathbb N$.\smallskip

\noindent\muid{MB-W01-L1-PR-A-007}\par 
\textbf{Шаг 1:} Проверьте область, указанную в условии.\par \textbf{Шаг 2:} Принадлежит ли $-3$ натуральным?\smallskip

\noindent\muid{MB-W01-L1-PR-A-008}\par 
\textbf{Шаг 1:} Разберите положительный, нулевой и отрицательный аргумент.\par \textbf{Шаг 2:} В каждом случае проверьте знак $x\cdot x$.\smallskip

\noindent\muid{MB-W01-L1-PR-A-009}\par 
\textbf{Шаг 1:} Ищите граничный случай.\par \textbf{Шаг 2:} Подставьте $x=0$.\smallskip

\noindent\muid{MB-W01-L1-PR-A-010}\par 
\textbf{Шаг 1:} Посмотрите на разность.\par \textbf{Шаг 2:} Она не зависит от $z$.\smallskip

\noindent\muid{MB-W01-L1-PR-A-011}\par 
\textbf{Шаг 1:} Подставьте любое целое, квадрат которого известен.\par \textbf{Шаг 2:} Возможны $4$ и $-4$.\smallskip

\noindent\muid{MB-W01-L1-PR-A-012}\par 
\textbf{Шаг 1:} В области ровно пять чисел.\par \textbf{Шаг 2:} Найдите их квадраты.\smallskip

\noindent\muid{MB-W01-L1-PR-A-013}\par 
\textbf{Шаг 1:} Поменяйте местами условие и заключение.\par \textbf{Шаг 2:} Для обратной стрелки рассмотрите $x=1$.\smallskip

\noindent\muid{MB-W01-L1-PR-A-014}\par 
\textbf{Шаг 1:} Это инструкция или повествование?\par \textbf{Шаг 2:} Отделите вопрос от равенства.\smallskip

\noindent\muid{MB-W01-L1-PR-A-015}\par 
\textbf{Шаг 1:} Вычтите одинаковое выражение $2a$.\par \textbf{Шаг 2:} Останется сравнить $5$ и $4$.\smallskip

\noindent\muid{MB-W01-L1-PR-A-016}\par 
\textbf{Шаг 1:} Сколько значений $n$ проверено?\par \textbf{Шаг 2:} Отделите истинность вывода от достаточности доказательства.\smallskip

\noindent\muid{MB-W01-L1-PR-A-017}\par 
\textbf{Шаг 1:} Найдите самый простой квадрат.\par \textbf{Шаг 2:} Проверьте $x=0$.\smallskip

\noindent\muid{MB-W01-L1-PR-A-018}\par 
\textbf{Шаг 1:} Назовите положительный квадратный корень.\par \textbf{Шаг 2:} Проверьте принадлежность области.\smallskip

\noindent\muid{MB-W01-L1-PR-A-019}\par 
\textbf{Шаг 1:} Используйте нейтральный элемент сложения.\par \textbf{Шаг 2:} Он входит в определение сложения.\smallskip

\noindent\muid{MB-W01-L1-PR-A-020}\par 
\textbf{Шаг 1:} Целые включают нуль и отрицательные.\par \textbf{Шаг 2:} Укажите одно подходящее значение.\smallskip

\noindent\muid{MB-W01-L1-PR-A-021}\par 
\textbf{Шаг 1:} Раскройте скобки.\par \textbf{Шаг 2:} Перегруппируйте две единицы.\smallskip

\noindent\muid{MB-W01-L1-PR-A-022}\par 
\textbf{Шаг 1:} Попробуйте небольшие положительные значения.\par \textbf{Шаг 2:} Положите $a=b=1$.\smallskip

\noindent\muid{MB-W01-L1-PR-A-023}\par 
\textbf{Шаг 1:} Как определено $\mathbb N$?\par \textbf{Шаг 2:} Прибавьте число $2$ к строгому неравенству.\smallskip

\noindent\muid{MB-W01-L1-PR-A-024}\par 
\textbf{Шаг 1:} Известно ли что-то о $x^2$?\par \textbf{Шаг 2:} Прибавьте $1$ к неотрицательному числу.\smallskip

\noindent\muid{MB-W01-L1-PR-A-025}\par 
\textbf{Шаг 1:} Ищите натуральное с нетривиальным делителем.\par \textbf{Шаг 2:} Например, $4=2\cdot2$.\smallskip

\noindent\muid{MB-W01-L1-PR-A-026}\par 
\textbf{Шаг 1:} Решите уравнение.\par \textbf{Шаг 2:} Получится отрицательное целое число.\smallskip

\noindent\muid{MB-W01-L1-PR-A-027}\par 
\textbf{Шаг 1:} Сравните $4$ и $1$.\par \textbf{Шаг 2:} Используйте транзитивность.\smallskip

\noindent\muid{MB-W01-L1-PR-A-028}\par 
\textbf{Шаг 1:} Найдите число, кратное $4$.\par \textbf{Шаг 2:} Оно может быть равно самому $4$.\smallskip

\noindent\muid{MB-W01-L1-PR-A-029}\par 
\textbf{Шаг 1:} Вспомните случай $x=0$.\par \textbf{Шаг 2:} Замените строгую характеристику на слабую.\smallskip

\noindent\muid{MB-W01-L1-PR-A-030}\par 
\textbf{Шаг 1:} Разложите разность квадратов.\par \textbf{Шаг 2:} $x^2-1=(x-1)(x+1)$.\smallskip

\section{Подсказки уровня B}

\noindent\muid{MB-W01-L1-PR-B-031}\par 
\textbf{Шаг 1:} Раскройте первые скобки.\par \textbf{Шаг 2:} После вычитания $5n$ останется число.\smallskip

\noindent\muid{MB-W01-L1-PR-B-032}\par 
\textbf{Шаг 1:} Исследуйте $x(x-1)$.\par \textbf{Шаг 2:} При вещественном $x=\tfrac12$ произведение отрицательно.\smallskip

\noindent\muid{MB-W01-L1-PR-B-033}\par 
\textbf{Шаг 1:} Умножьте условие на положительное число.\par \textbf{Шаг 2:} Получите $2x>16$, затем сравните с $12$.\smallskip

\noindent\muid{MB-W01-L1-PR-B-034}\par 
\textbf{Шаг 1:} Квадрат не сохраняет знак числа.\par \textbf{Шаг 2:} Возьмите $x=-2$.\smallskip

\noindent\muid{MB-W01-L1-PR-B-035}\par 
\textbf{Шаг 1:} Вынесите общий множитель.\par \textbf{Шаг 2:} Проверьте целочисленность второго множителя.\smallskip

\noindent\muid{MB-W01-L1-PR-B-036}\par 
\textbf{Шаг 1:} На что пытаются сократить?\par \textbf{Шаг 2:} Нуль не имеет обратного при умножении.\smallskip

\noindent\muid{MB-W01-L1-PR-B-037}\par 
\textbf{Шаг 1:} Раскройте оба квадрата.\par \textbf{Шаг 2:} Члены $x^2$ и $9$ сократятся.\smallskip

\noindent\muid{MB-W01-L1-PR-B-038}\par 
\textbf{Шаг 1:} Запишите $n=7k$.\par \textbf{Шаг 2:} Множитель $3k$ должен быть целым.\smallskip

\noindent\muid{MB-W01-L1-PR-B-039}\par 
\textbf{Шаг 1:} Сравните отрицательное и положительное числа.\par \textbf{Шаг 2:} Подойдут $a=-3$, $b=2$.\smallskip

\noindent\muid{MB-W01-L1-PR-B-040}\par 
\textbf{Шаг 1:} Разложите выражение на множители.\par \textbf{Шаг 2:} Получите $m(m-5)=0$.\smallskip

\noindent\muid{MB-W01-L1-PR-B-041}\par 
\textbf{Шаг 1:} Используйте свойство квадрата.\par \textbf{Шаг 2:} Когда $x^2=0$?\smallskip

\noindent\muid{MB-W01-L1-PR-B-042}\par 
\textbf{Шаг 1:} Группируйте буквенные и постоянные слагаемые.\par \textbf{Шаг 2:} Сумма смещений равна $0+1+2+3$.\smallskip

\noindent\muid{MB-W01-L1-PR-B-043}\par 
\textbf{Шаг 1:} Сложите данные неравенства.\par \textbf{Шаг 2:} Числовая граница равна $2+3$.\smallskip

\noindent\muid{MB-W01-L1-PR-B-044}\par 
\textbf{Шаг 1:} Один пример не охватывает всех пар.\par \textbf{Шаг 2:} Проверьте $2\cdot2$.\smallskip

\noindent\muid{MB-W01-L1-PR-B-045}\par 
\textbf{Шаг 1:} Вспомните знак числа после возведения в квадрат.\par \textbf{Шаг 2:} Рассмотрите пару $(1,-1)$.\smallskip

\noindent\muid{MB-W01-L1-PR-B-046}\par 
\textbf{Шаг 1:} Умножение на отрицательное меняет знак.\par \textbf{Шаг 2:} Следите за направлением строгого неравенства.\smallskip

\noindent\muid{MB-W01-L1-PR-B-047}\par 
\textbf{Шаг 1:} Используйте нулевое произведение.\par \textbf{Шаг 2:} Каждый множитель может быть нулём.\smallskip

\noindent\muid{MB-W01-L1-PR-B-048}\par 
\textbf{Шаг 1:} Свидетель может зависеть от $n$.\par \textbf{Шаг 2:} Выберите $m$ по явной формуле.\smallskip

\noindent\muid{MB-W01-L1-PR-B-049}\par 
\textbf{Шаг 1:} После фиксирования $m$ найдите большее натуральное.\par \textbf{Шаг 2:} Возьмите $n=m+1$.\smallskip

\noindent\muid{MB-W01-L1-PR-B-050}\par 
\textbf{Шаг 1:} Достаточно нарушить хотя бы одну часть.\par \textbf{Шаг 2:} Выберите натуральное число $1$.\smallskip

\noindent\muid{MB-W01-L1-PR-B-051}\par 
\textbf{Шаг 1:} Используйте линейный порядок чисел.\par \textbf{Шаг 2:} Иных возможностей для сравнения с нулём нет.\smallskip

\noindent\muid{MB-W01-L1-PR-B-052}\par 
\textbf{Шаг 1:} Поделите обе части на $a$.\par \textbf{Шаг 2:} Условие $a\ne0$ делает операцию допустимой.\smallskip

\noindent\muid{MB-W01-L1-PR-B-053}\par 
\textbf{Шаг 1:} Узнайте полный квадрат.\par \textbf{Шаг 2:} $x^2-2x+1=(x-1)^2$.\smallskip

\noindent\muid{MB-W01-L1-PR-B-054}\par 
\textbf{Шаг 1:} Испытайте небольшие числа.\par \textbf{Шаг 2:} Подставьте $a=b=1$.\smallskip

\noindent\muid{MB-W01-L1-PR-B-055}\par 
\textbf{Шаг 1:} Используйте транзитивность порядка.\par \textbf{Шаг 2:} Цепочка $a\ge b\ge0$.\smallskip

\noindent\muid{MB-W01-L1-PR-B-056}\par 
\textbf{Шаг 1:} Представьте четвёртую степень квадратом.\par \textbf{Шаг 2:} $x^4=(x^2)^2$.\smallskip

\noindent\muid{MB-W01-L1-PR-B-057}\par 
\textbf{Шаг 1:} Вычислите $t^2-t$.\par \textbf{Шаг 2:} Оба множителя $t$ и $t-1$ неотрицательны.\smallskip

\noindent\muid{MB-W01-L1-PR-B-058}\par 
\textbf{Шаг 1:} Вычтите $7$ из обеих частей.\par \textbf{Шаг 2:} Равносильность преобразования важна для единственности.\smallskip

\noindent\muid{MB-W01-L1-PR-B-059}\par 
\textbf{Шаг 1:} Откуда взято основание?\par \textbf{Шаг 2:} Нельзя использовать доказываемый факт как единственный аргумент.\smallskip

\noindent\muid{MB-W01-L1-PR-B-060}\par 
\textbf{Шаг 1:} Для $x>0$ произведение двух положительных положительно.\par \textbf{Шаг 2:} В изменённом условии разрешён нуль.\smallskip

\section{Подсказки уровня C}

\noindent\muid{MB-W01-L1-PR-C-061}\par 
\textbf{Шаг 1:} Рассмотрите знак двух множителей на указанной области.\par \textbf{Шаг 2:} Между $0$ и $2$ лежит одно целое число.\smallskip

\noindent\muid{MB-W01-L1-PR-C-062}\par 
\textbf{Шаг 1:} Выделите квадрат.\par \textbf{Шаг 2:} $x^2+6x+10=(x+3)^2+1$.\smallskip

\noindent\muid{MB-W01-L1-PR-C-063}\par 
\textbf{Шаг 1:} Выпишите стрелку между свойствами.\par \textbf{Шаг 2:} Проверьте обратную импликацию на $20$.\smallskip

\noindent\muid{MB-W01-L1-PR-C-064}\par 
\textbf{Шаг 1:} Сравнение вещественных чисел полно и однозначно.\par \textbf{Шаг 2:} Взаимное исключение должно быть указано.\smallskip

\noindent\muid{MB-W01-L1-PR-C-065}\par 
\textbf{Шаг 1:} Перенесите всё в одну сторону.\par \textbf{Шаг 2:} Разложите $a^2-b^2$.\smallskip

\noindent\muid{MB-W01-L1-PR-C-066}\par 
\textbf{Шаг 1:} Модуль стирает информацию о знаке.\par \textbf{Шаг 2:} Возьмите числа $3$ и $-3$.\smallskip

\noindent\muid{MB-W01-L1-PR-C-067}\par 
\textbf{Шаг 1:} Какое число противоположно $a$?\par \textbf{Шаг 2:} Положите $b=-a$.\smallskip

\noindent\muid{MB-W01-L1-PR-C-068}\par 
\textbf{Шаг 1:} Как изменится равенство при $a=0$ и $a=1$?\par \textbf{Шаг 2:} Один и тот же $b$ должен удовлетворять обоим.\smallskip

\noindent\muid{MB-W01-L1-PR-C-069}\par 
\textbf{Шаг 1:} Разложите разность квадратов.\par \textbf{Шаг 2:} Знаки $u-v$ и $u+v$ положительны.\smallskip

\noindent\muid{MB-W01-L1-PR-C-070}\par 
\textbf{Шаг 1:} Используйте определение модуля по знаку.\par \textbf{Шаг 2:} Проверьте оба направления.\smallskip

\noindent\muid{MB-W01-L1-PR-C-071}\par 
\textbf{Шаг 1:} В квадрате исчезает знак.\par \textbf{Шаг 2:} Разберите корни $3$ и $-3$.\smallskip

\noindent\muid{MB-W01-L1-PR-C-072}\par 
\textbf{Шаг 1:} Сравните каждую разность с нулём.\par \textbf{Шаг 2:} Обе разности равны $(a-b)/2$.\smallskip

\noindent\muid{MB-W01-L1-PR-C-073}\par 
\textbf{Шаг 1:} Используйте обратимые операции.\par \textbf{Шаг 2:} Прибавьте $5$, затем поделите на $3$.\smallskip

\noindent\muid{MB-W01-L1-PR-C-074}\par 
\textbf{Шаг 1:} Каков знак каждого квадрата?\par \textbf{Шаг 2:} Сумма двух неотрицательных чисел неотрицательна.\smallskip

\noindent\muid{MB-W01-L1-PR-C-075}\par 
\textbf{Шаг 1:} Рассмотрите два отрицательных числа.\par \textbf{Шаг 2:} Например, $a=-1$, $b=-2$.\smallskip

\noindent\muid{MB-W01-L1-PR-C-076}\par 
\textbf{Шаг 1:} Свяжите квадрат с модулем.\par \textbf{Шаг 2:} Из $n^2\le4$ следует $-2\le n\le2$.\smallskip

\noindent\muid{MB-W01-L1-PR-C-077}\par 
\textbf{Шаг 1:} Умножайте только на положительные величины.\par \textbf{Шаг 2:} Вычислите $\frac1b-\frac1a$.\smallskip

\noindent\muid{MB-W01-L1-PR-C-078}\par 
\textbf{Шаг 1:} Равенство сохраняется при умножении одинаковых величин.\par \textbf{Шаг 2:} Запишите $a^3=a\cdot a\cdot a$.\smallskip

\noindent\muid{MB-W01-L1-PR-C-079}\par 
\textbf{Шаг 1:} Квадрат может быть равен нулю.\par \textbf{Шаг 2:} Подставьте $x=-1$.\smallskip

\noindent\muid{MB-W01-L1-PR-C-080}\par 
\textbf{Шаг 1:} Произведение неотрицательных неотрицательно.\par \textbf{Шаг 2:} Сравните $ab$ с $2b$ и $6$.\smallskip

\section{Подсказки уровня D}

\noindent\muid{MB-W01-L1-PR-D-081}\par 
\textbf{Шаг 1:} Разложите квадратный трёхчлен.\par \textbf{Шаг 2:} Исследуйте знак $(x-1)(x-3)$.\smallskip

\noindent\muid{MB-W01-L1-PR-D-082}\par 
\textbf{Шаг 1:} Для первого выбирайте $y$ после $x$.\par \textbf{Шаг 2:} Во втором одно $y$ должно работать при $x=0$ и $x=1$.\smallskip

\noindent\muid{MB-W01-L1-PR-D-083}\par 
\textbf{Шаг 1:} Сумма неотрицательных слагаемых равна нулю.\par \textbf{Шаг 2:} Предположите, что одно из слагаемых положительно.\smallskip

\noindent\muid{MB-W01-L1-PR-D-084}\par 
\textbf{Шаг 1:} Найдите параметр, при котором выражение имеет множитель $17$.\par \textbf{Шаг 2:} Испытайте $n=17$.\smallskip

\noindent\muid{MB-W01-L1-PR-D-085}\par 
\textbf{Шаг 1:} Подберите координаты симметричной трапеции.\par \textbf{Шаг 2:} Можно взять $A=(-2,0)$, $B=(2,0)$, $C=(1,2)$, $D=(-1,2)$.\smallskip

\noindent\muid{MB-W01-L1-PR-D-086}\par 
\textbf{Шаг 1:} Положительное $x>2$ можно возвести в квадрат.\par \textbf{Шаг 2:} Выведите $x^2>4$.\smallskip

\noindent\muid{MB-W01-L1-PR-D-087}\par 
\textbf{Шаг 1:} Разложите $x(x-2)$.\par \textbf{Шаг 2:} Проверьте $x=1$ и отрицательные целые.\smallskip

\noindent\muid{MB-W01-L1-PR-D-088}\par 
\textbf{Шаг 1:} Какой множитель сократили?\par \textbf{Шаг 2:} Изучите случаи $a=0$ и $a\ne0$.\smallskip

\noindent\muid{MB-W01-L1-PR-D-089}\par 
\textbf{Шаг 1:} Возведите в квадрат разность $a-b$.\par \textbf{Шаг 2:} Используйте $(a-b)^2=(a+b)^2-4ab$.\smallskip

\noindent\muid{MB-W01-L1-PR-D-090}\par 
\textbf{Шаг 1:} Выведите произведение $xy$ из квадрата суммы.\par \textbf{Шаг 2:} Уравнения заставляют $(x-y)^2$ обратиться в нуль.\smallskip

\chapter{Алгоритмы, эвристики и методы проверки}
Ниже перечислены сигналы в условии и действия, полезные на первом занятии. Каждому приёму сопутствует ограничение: универсальных «волшебных правил» для олимпиад не существует. Используйте памятки во время поиска, затем объясняйте каждую применённую операцию в чистовом доказательстве.

\section{Приём 1: Слово «любой»}

\begin{keybox}\textbf{Сигнал в условии.} Слово «любой»\par \textbf{Что пробовать.} Назовите произвольный допустимый объект и выводите формулу без специальных чисел.\end{keybox}

\textbf{Когда опасно.} Опасно ограничиться проверкой нескольких значений.\par \textbf{Мини-проба.} \muid{MB-W01-L1-HR-01} Докажите $(n+1)-n=1$ для каждого целого $n$.


\section{Приём 2: Слово «существует»}

\begin{keybox}\textbf{Сигнал в условии.} Слово «существует»\par \textbf{Что пробовать.} Ищите конструкцию или конкретного свидетеля и проверяйте все условия.\end{keybox}

\textbf{Когда опасно.} Опасно предъявить объект вне заявленной области.\par \textbf{Мини-проба.} \muid{MB-W01-L1-HR-02} Докажите существование целого корня $x^2=36$.


\section{Приём 3: Хочется опровергнуть «для всех»}

\begin{keybox}\textbf{Сигнал в условии.} Хочется опровергнуть «для всех»\par \textbf{Что пробовать.} Проверьте нуль, единицу, отрицательные значения и границы области.\end{keybox}

\textbf{Когда опасно.} Опасно брать запрещённое значение переменной.\par \textbf{Мини-проба.} \muid{MB-W01-L1-HR-03} Опровергните «каждое вещественное число имеет ненулевой квадрат».


\section{Приём 4: Заявлено «ровно одно»}

\begin{keybox}\textbf{Сигнал в условии.} Заявлено «ровно одно»\par \textbf{Что пробовать.} Разделите доказательство на существование и единственность.\end{keybox}

\textbf{Когда опасно.} Опасно показать лишь одну подходящую подстановку.\par \textbf{Мини-проба.} \muid{MB-W01-L1-HR-04} Найдите все решения $x+4=9$.


\section{Приём 5: Утверждение с «если»}

\begin{keybox}\textbf{Сигнал в условии.} Утверждение с «если»\par \textbf{Что пробовать.} Подпишите предпосылку $P$ и заключение $Q$; доказывайте $P\Rightarrow Q$.\end{keybox}

\textbf{Когда опасно.} Опасна самовольная замена стрелки на обратную.\par \textbf{Мини-проба.} \muid{MB-W01-L1-HR-05} Истинно ли «если $x>6$, то $x>2$»?


\section{Приём 6: Перед вами алгебраическое тождество}

\begin{keybox}\textbf{Сигнал в условии.} Перед вами алгебраическое тождество\par \textbf{Что пробовать.} Преобразуйте одну сторону через известные законы и сопоставьте с другой.\end{keybox}

\textbf{Когда опасно.} Опасны неоговорённые деления на буквы.\par \textbf{Мини-проба.} \muid{MB-W01-L1-HR-06} Проверьте $(a-2)^2=a^2-4a+4$.


\section{Приём 7: Появляется неравенство}

\begin{keybox}\textbf{Сигнал в условии.} Появляется неравенство\par \textbf{Что пробовать.} Сравните разность сторон с нулём; ищите полный квадрат.\end{keybox}

\textbf{Когда опасно.} Не забывайте знак множителя и равенство.\par \textbf{Мини-проба.} \muid{MB-W01-L1-HR-07} Докажите $x^2+2x+1\ge0$.


\section{Приём 8: Нужно разбить на случаи}

\begin{keybox}\textbf{Сигнал в условии.} Нужно разбить на случаи\par \textbf{Что пробовать.} Проверьте, что случаи покрывают всю область и не оставляют граничной щели.\end{keybox}

\textbf{Когда опасно.} Опасен пропущенный нуль или граница.\par \textbf{Мини-проба.} \muid{MB-W01-L1-HR-08} Докажите $|x|\ge0$ при $x<0$, $x=0$, $x>0$.


\section{Приём 9: Используется рисунок}

\begin{keybox}\textbf{Сигнал в условии.} Используется рисунок\par \textbf{Что пробовать.} Отделите данные геометрии от визуальных впечатлений; подпишите свойства.\end{keybox}

\textbf{Когда опасно.} Пропорции на рисунке не гарантируют равенства.\par \textbf{Мини-проба.} \muid{MB-W01-L1-HR-09} Придумайте трапецию, которая не является прямоугольником.


\section{Приём 10: Подозрительная короткая выкладка}

\begin{keybox}\textbf{Сигнал в условии.} Подозрительная короткая выкладка\par \textbf{Что пробовать.} Проверьте каждую стрелку, значение делителя, кванторы и полноту ответа.\end{keybox}

\textbf{Когда опасно.} Опасен незаметный переход из следствия в эквивалентность.\par \textbf{Мини-проба.} \muid{MB-W01-L1-HR-10} Почему из $x^2=4$ нельзя заключить только $x=2$?


\section{Приём 11: Чисел в условии много}

\begin{keybox}\textbf{Сигнал в условии.} Чисел в условии много\par \textbf{Что пробовать.} Сначала упростите структуру: сумму, разность, произведение или общую переменную.\end{keybox}

\textbf{Когда опасно.} Не превращайте универсальную задачу в набор вычислений.\par \textbf{Мини-проба.} \muid{MB-W01-L1-HR-11} Преобразуйте $(n+2)+(n+3)$.


\section{Четыре алгоритма выбора}
\begin{enumerate}
\item \textbf{Проверка формулировки.} Найдите переменную; обозначьте область; восстановите квантор; проверьте, осталось ли свободное место в условии.
\item \textbf{Поиск универсального доказательства.} Возьмите произвольный допустимый объект; запишите свойства из определения; преобразуйте выражение; сделайте общий вывод.
\item \textbf{Поиск контрпримера.} Подчеркните все условия; испытайте границы и малые значения; предъявите допустимый объект; отдельно покажите ложность заключения.
\item \textbf{Аудит решения.} Проверьте каждую стрелку, возможное деление на нуль, граничные случаи, обратную проверку, полноту ответа и точную заключительную фразу.
\end{enumerate}
\section{Рубрика оценивания математического текста}
\begin{longtable}{p{.15\textwidth}p{.74\textwidth}}\toprule Балл & Критерий\\\midrule\endhead
0 & Нет работающей идеи; условия переписаны без полезного перехода.\\
1 & Найден частный случай или правильная идея, но она не обоснована для требуемой области.\\
2 & Присутствует основной механизм, однако нарушена полнота случаев или есть существенный логический пробел.\\
3 & Доказательство по существу верно, остаются небольшие недостатки оформления или непрояснённая очевидная ссылка.\\
4 & Доказаны все требуемые направления, условия использованы корректно, вывод совпадает с целью, нет необоснованных переходов.\\\bottomrule\end{longtable}
\section{Режим повторения}
В день занятия вернитесь к пяти наиболее трудным задачам. Через 24 часа запишите решения заново без подсказок, через неделю решите выборку из задач A, B, C и D. В журнале прогресса фиксируйте «решил сам», «после подсказки», «прочитал решение» и «вернуться». Понимание доказательства можно считать устойчивым, когда вы способны воспроизвести ключевой довод и объяснить границы применимости.

\chapter{Полные решения всех 90 задач}
\begin{keybox}
В этом разделе каждый номер с \texttt{PR} имеет ровно одно эталонное полное решение. Пояснения учитывают область, основания переходов и характер требуемого вывода. Для доказательства достаточно одного строгого метода; некоторые решения допускают альтернативы, которые полезно обсудить самостоятельно.
\end{keybox}

\section{Решения уровня A}

\subsection*{Числовая истина · № 1}\hypertarget{solution:MB-W01-L1-PR-A-001}{}

\muid{MB-W01-L1-PR-A-001}\par

\textbf{Условие.} Определите истинность высказывания $14-6=8$ и объясните ответ.\par

\textbf{Полное решение.} Высказывание полностью определено. Выполняем вычитание: $14-6=8$. Левая и правая части равны; следовательно, утверждение истинно. В данном случае достаточно одного вычисления, потому что числовое выражение не содержит переменной.


\begin{learnbox}\textbf{Ключевой итог.} Истина: $14-6=8$.\par \textbf{Проверка доказательства.} 
Следите за условиями применимости правил и за переходом от произвольного объекта к общему выводу. 
\end{learnbox}

\subsection*{Числовая ложь · № 2}\hypertarget{solution:MB-W01-L1-PR-A-002}{}

\muid{MB-W01-L1-PR-A-002}\par

\textbf{Условие.} Верно ли $8\cdot7=54$? Укажите точную причину.\par

\textbf{Полное решение.} Произведение $8\cdot7=56$. Поскольку $56\ne54$, заданное числовое высказывание ложно. Здесь не требуется контрпример с переменной: обе части обозначают конкретные числа.


\begin{learnbox}\textbf{Ключевой итог.} Ложь: $8\cdot7=56\ne54$.\par \textbf{Проверка доказательства.} 
Следите за условиями применимости правил и за переходом от произвольного объекта к общему выводу. 
\end{learnbox}

\subsection*{Открытая формула · № 3}\hypertarget{solution:MB-W01-L1-PR-A-003}{}

\muid{MB-W01-L1-PR-A-003}\par

\textbf{Условие.} Является ли запись $x-4=1$ высказыванием без дополнительных условий?\par

\textbf{Полное решение.} Буква $x$ не связана квантором и ей не присвоено значение. При $x=5$ равенство верно, поскольку $5-4=1$; при $x=2$ ложно, поскольку $2-4=-2$. Значит, сама запись является открытой формулой, истинность которой зависит от выбора $x$.


\begin{learnbox}\textbf{Ключевой итог.} Открытая формула; нужен выбор $x$ или квантор.\par \textbf{Проверка доказательства.} 
Следите за условиями применимости правил и за переходом от произвольного объекта к общему выводу. 
\end{learnbox}

\subsection*{Область · № 4}\hypertarget{solution:MB-W01-L1-PR-A-004}{}

\muid{MB-W01-L1-PR-A-004}\par

\textbf{Условие.} Истинно ли $\forall x\in\mathbb R\colon x+2>x$?\par

\textbf{Полное решение.} Для произвольного вещественного $x$ вычтем $x$ из обеих частей предполагаемого неравенства. Получим $2>0$, что верно. Равносильные преобразования действуют для всех вещественных чисел; исходное универсальное утверждение истинно.


\begin{learnbox}\textbf{Ключевой итог.} Истина для каждого $x\in\mathbb R$.\par \textbf{Проверка доказательства.} 
Следите за условиями применимости правил и за переходом от произвольного объекта к общему выводу. 
\end{learnbox}

\subsection*{Свидетель · № 5}\hypertarget{solution:MB-W01-L1-PR-A-005}{}

\muid{MB-W01-L1-PR-A-005}\par

\textbf{Условие.} Докажите: $\exists n\in\mathbb N\colon 3n=15$.\par

\textbf{Полное решение.} Выберем $n=5$. Это натуральное число, и $3\cdot5=15$. Мы предъявили объект из указанной области и проверили нужное равенство, поэтому существование доказано. Перечислять другие числа не требуется.


\begin{learnbox}\textbf{Ключевой итог.} Свидетель $n=5$.\par \textbf{Проверка доказательства.} 
Следите за условиями применимости правил и за переходом от произвольного объекта к общему выводу. 
\end{learnbox}

\subsection*{Несуществование · № 6}\hypertarget{solution:MB-W01-L1-PR-A-006}{}

\muid{MB-W01-L1-PR-A-006}\par

\textbf{Условие.} Существует ли натуральное $n$ с $n+2=1$?\par

\textbf{Полное решение.} Из $n+2=1$ вычитанием $2$ получаем $n=-1$. Но $-1\notin\mathbb N=\{1,2,\ldots\}$. Следовательно, подходящего натурального числа нет. Обратите внимание: в целых числах решение существует.


\begin{learnbox}\textbf{Ключевой итог.} Нет; единственное числовое решение $n=-1$.\par \textbf{Проверка доказательства.} 
Следите за условиями применимости правил и за переходом от произвольного объекта к общему выводу. 
\end{learnbox}

\subsection*{Допустимость примера · № 7}\hypertarget{solution:MB-W01-L1-PR-A-007}{}

\muid{MB-W01-L1-PR-A-007}\par

\textbf{Условие.} Можно ли использовать $x=-3$ как контрпример к утверждению «для всех натуральных $x$ выполняется $x>0$»?\par

\textbf{Полное решение.} Нет. Число $-3$ не принадлежит $\mathbb N$, поэтому не является допустимым объектом для универсального высказывания о натуральных числах. Более того, по нашему определению каждое натуральное положительно, так что утверждение истинно.


\begin{learnbox}\textbf{Ключевой итог.} Нельзя: $-3\notin\mathbb N$.\par \textbf{Проверка доказательства.} 
Помимо проверки найденных значений, аргумент исключает остальные варианты. 
\end{learnbox}

\subsection*{Квадраты · № 8}\hypertarget{solution:MB-W01-L1-PR-A-008}{}

\muid{MB-W01-L1-PR-A-008}\par

\textbf{Условие.} Истинно ли «каждый квадрат вещественного числа неотрицателен»?\par

\textbf{Полное решение.} Если $x>0$, квадрат положителен; если $x=0$, квадрат равен нулю; если $x<0$, произведение двух отрицательных множителей положительно. Таким образом, во всех трёх исчерпывающих случаях $x^2\ge0$. Высказывание истинно.


\begin{learnbox}\textbf{Ключевой итог.} Истина: $x^2\ge0$ при всех $x\in\mathbb R$.\par \textbf{Проверка доказательства.} 
Следите за условиями применимости правил и за переходом от произвольного объекта к общему выводу. 
\end{learnbox}

\subsection*{Строгость · № 9}\hypertarget{solution:MB-W01-L1-PR-A-009}{}

\muid{MB-W01-L1-PR-A-009}\par

\textbf{Условие.} Опровергните утверждение $\forall x\in\mathbb R\colon x^2>0$.\par

\textbf{Полное решение.} Возьмём $x=0$, это допустимое вещественное число. Его квадрат равен $0$, а не является строго положительным. Значит, универсальное утверждение ложно. Исправление состоит в замене строгого $>0$ на нестрогое $\ge0$.


\begin{learnbox}\textbf{Ключевой итог.} Ложь; контрпример $x=0$.\par \textbf{Проверка доказательства.} 
Следите за условиями применимости правил и за переходом от произвольного объекта к общему выводу. 
\end{learnbox}

\subsection*{Отрицательные числа · № 10}\hypertarget{solution:MB-W01-L1-PR-A-010}{}

\muid{MB-W01-L1-PR-A-010}\par

\textbf{Условие.} Верно ли $\forall z\in\mathbb Z\colon z+1>z$?\par

\textbf{Полное решение.} Для произвольного целого $z$ равенство $(z+1)-z=1$ показывает, что разность строго положительна. Следовательно, $z+1>z$ для всех целых значений. Отрицательные числа не являются исключением.


\begin{learnbox}\textbf{Ключевой итог.} Истина; разность равна $1>0$.\par \textbf{Проверка доказательства.} 
Следите за условиями применимости правил и за переходом от произвольного объекта к общему выводу. 
\end{learnbox}

\subsection*{Квантор существования · № 11}\hypertarget{solution:MB-W01-L1-PR-A-011}{}

\muid{MB-W01-L1-PR-A-011}\par

\textbf{Условие.} Верно ли $\exists x\in\mathbb Z\colon x^2=16$?\par

\textbf{Полное решение.} Выберем $x=4$. Это целое число, и $x^2=4^2=16$. Тем самым существование доказано. Значение $-4$ также подходит, но для утверждения с квантором существования второго свидетеля искать необязательно.


\begin{learnbox}\textbf{Ключевой итог.} Истина; например $x=4$.\par \textbf{Проверка доказательства.} 
Отмечены зависимости переменных и порядок выбора свидетеля. 
\end{learnbox}

\subsection*{Конечная область · № 12}\hypertarget{solution:MB-W01-L1-PR-A-012}{}

\muid{MB-W01-L1-PR-A-012}\par

\textbf{Условие.} Проверьте $x^2\le4$ для всех $x\in\{-2,-1,0,1,2\}$.\par

\textbf{Полное решение.} Квадраты перечисленных значений: $4,1,0,1,4$. Каждый из них не превышает $4$. Проверены все пять элементов заданной конечной области, поэтому утверждение доказано полным перебором; никакие другие $x$ сюда не входят.


\begin{learnbox}\textbf{Ключевой итог.} Истина на указанном пятиэлементном множестве.\par \textbf{Проверка доказательства.} 
Помимо проверки найденных значений, аргумент исключает остальные варианты. 
\end{learnbox}

\subsection*{Обратное · № 13}\hypertarget{solution:MB-W01-L1-PR-A-013}{}

\muid{MB-W01-L1-PR-A-013}\par

\textbf{Условие.} Сформулируйте обратное к «если $x>9$, то $x>0$» и оцените его истинность.\par

\textbf{Полное решение.} Обратное утверждение: «если $x>0$, то $x>9$». Оно ложно: число $x=1$ положительно, но не больше $9$. Изначальная стрелка верна, поскольку $9>0$; истинность прямой стрелки не гарантирует истинности обратной.


\begin{learnbox}\textbf{Ключевой итог.} Обратное ложно; $x=1$.\par \textbf{Проверка доказательства.} 
Следите за условиями применимости правил и за переходом от произвольного объекта к общему выводу. 
\end{learnbox}

\subsection*{Переменная в задаче · № 14}\hypertarget{solution:MB-W01-L1-PR-A-014}{}

\muid{MB-W01-L1-PR-A-014}\par

\textbf{Условие.} Является ли предложение «найдите $x$, если $x+7=10$» математическим высказыванием?\par

\textbf{Полное решение.} Фраза представляет собой задание найти неизвестное, а не повествовательное утверждение с определённым значением истинности. Само равенство $x+7=10$ является открытой формулой. При $x=3$ оно верно, но просьба «найдите» не бывает истинной или ложной.


\begin{learnbox}\textbf{Ключевой итог.} Нет; это инструкция к решению.\par \textbf{Проверка доказательства.} 
Следите за условиями применимости правил и за переходом от произвольного объекта к общему выводу. 
\end{learnbox}

\subsection*{Неравенство · № 15}\hypertarget{solution:MB-W01-L1-PR-A-015}{}

\muid{MB-W01-L1-PR-A-015}\par

\textbf{Условие.} Докажите $2a+5>2a+4$ для любого вещественного $a$.\par

\textbf{Полное решение.} Вычитание $2a$ из обеих частей неравенства равносильно переходу к $5>4$. Последнее истинно. Следовательно, исходное неравенство выполнено при каждом вещественном $a$, включая отрицательные значения.


\begin{learnbox}\textbf{Ключевой итог.} Истина: разность сторон равна $1$.\par \textbf{Проверка доказательства.} 
Следите за условиями применимости правил и за переходом от произвольного объекта к общему выводу. 
\end{learnbox}

\subsection*{Вычисление и вывод · № 16}\hypertarget{solution:MB-W01-L1-PR-A-016}{}

\muid{MB-W01-L1-PR-A-016}\par

\textbf{Условие.} Оцените фразу: «$1^2+1=2$ и $2^2+2=6$, значит $n^2+n$ чётно при любом $n$».\par

\textbf{Полное решение.} Проверены только $n=1$ и $n=2$. Такие вычисления не дают универсального доказательства для бесконечной области $\mathbb N$. Заключение действительно верно, но нуждается в общем рассуждении, например через разбор чётности соседних чисел; приведённые две проверки недостаточны.


\begin{learnbox}\textbf{Ключевой итог.} Аргумент недостаточен: проверены два значения.\par \textbf{Проверка доказательства.} 
Следите за условиями применимости правил и за переходом от произвольного объекта к общему выводу. 
\end{learnbox}

\subsection*{Утверждение о нуле · № 17}\hypertarget{solution:MB-W01-L1-PR-A-017}{}

\muid{MB-W01-L1-PR-A-017}\par

\textbf{Условие.} Верно ли $\exists x\in\mathbb R\colon x^2=0$?\par

\textbf{Полное решение.} Число $0$ принадлежит $\mathbb R$ и удовлетворяет равенству $0^2=0$. Поэтому существование доказано. При необходимости можно дополнительно доказать единственность, но заданный вопрос требует только существования.


\begin{learnbox}\textbf{Ключевой итог.} Истина; $x=0$.\par \textbf{Проверка доказательства.} 
Следите за условиями применимости правил и за переходом от произвольного объекта к общему выводу. 
\end{learnbox}

\subsection*{Положительный корень · № 18}\hypertarget{solution:MB-W01-L1-PR-A-018}{}

\muid{MB-W01-L1-PR-A-018}\par

\textbf{Условие.} Существует ли натуральное $n$, для которого $n^2=25$?\par

\textbf{Полное решение.} Положительное число $n=5$ натуральное и даёт $5^2=25$. Следовательно, такой объект существует. Отрицательное $-5$ тоже имеет квадрат $25$, но не принадлежит выбранной натуральной области.


\begin{learnbox}\textbf{Ключевой итог.} Да; $n=5$.\par \textbf{Проверка доказательства.} 
Следите за условиями применимости правил и за переходом от произвольного объекта к общему выводу. 
\end{learnbox}

\subsection*{Равенство · № 19}\hypertarget{solution:MB-W01-L1-PR-A-019}{}

\muid{MB-W01-L1-PR-A-019}\par

\textbf{Условие.} Докажите тождество $a+0=a$ для всех вещественных $a$.\par

\textbf{Полное решение.} Нуль является нейтральным элементом относительно сложения: прибавление нуля не меняет число. Поэтому для произвольного $a\in\mathbb R$ выполняется $a+0=a$. Свободного выбора конкретного $a$ не потребовалось.


\begin{learnbox}\textbf{Ключевой итог.} Истина по свойству нуля при сложении.\par \textbf{Проверка доказательства.} 
Помимо проверки найденных значений, аргумент исключает остальные варианты. 
\end{learnbox}

\subsection*{Контрпример по области · № 20}\hypertarget{solution:MB-W01-L1-PR-A-020}{}

\muid{MB-W01-L1-PR-A-020}\par

\textbf{Условие.} Опровергните утверждение «любой целый $n$ положителен».\par

\textbf{Полное решение.} Возьмём $n=0$, это целое число. Оно не положительно, поскольку строгий знак $0>0$ ложен. Одного допустимого значения достаточно, чтобы опровергнуть универсальную формулировку.


\begin{learnbox}\textbf{Ключевой итог.} Ложь; $n=0$.\par \textbf{Проверка доказательства.} 
Числовой пример проверен на принадлежность области и на нарушение именно требуемого вывода. 
\end{learnbox}

\subsection*{Сумма · № 21}\hypertarget{solution:MB-W01-L1-PR-A-021}{}

\muid{MB-W01-L1-PR-A-021}\par

\textbf{Условие.} Покажите, что $(a+1)+(b+1)=a+b+2$ при любых вещественных $a,b$.\par

\textbf{Полное решение.} По сочетательному и переместительному законам сложения $(a+1)+(b+1)=a+b+(1+1)=a+b+2$. Эти законы верны при любых вещественных $a$ и $b$, поэтому получено тождество.


\begin{learnbox}\textbf{Ключевой итог.} Истина: обе части тождественны.\par \textbf{Проверка доказательства.} 
Следите за условиями применимости правил и за переходом от произвольного объекта к общему выводу. 
\end{learnbox}

\subsection*{Неверная формула · № 22}\hypertarget{solution:MB-W01-L1-PR-A-022}{}

\muid{MB-W01-L1-PR-A-022}\par

\textbf{Условие.} Опровергните формулу $(a+b)^2=a^2+b^2$ для всех вещественных $a,b$.\par

\textbf{Полное решение.} При $a=b=1$ левая часть равна $(1+1)^2=4$, правая равна $1^2+1^2=2$. Значения различны. Значит, универсальная формула ложна. Корректное раскрытие квадрата суммы включает также $2ab$.


\begin{learnbox}\textbf{Ключевой итог.} Ложь; при $(1,1)$ имеем $4\ne2$.\par \textbf{Проверка доказательства.} 
Помимо проверки найденных значений, аргумент исключает остальные варианты. 
\end{learnbox}

\subsection*{Части доказательства · № 23}\hypertarget{solution:MB-W01-L1-PR-A-023}{}

\muid{MB-W01-L1-PR-A-023}\par

\textbf{Условие.} В фразе «пусть $n\in\mathbb N$; тогда $n>0$; следовательно, $n+2>2$» назовите данное, основание и вывод.\par

\textbf{Полное решение.} Дано $n\in\mathbb N$, то есть натуральное положительное число. Основание: по определению $n>0$. Прибавление $2$ к обеим частям даёт $n+2>2$. Последняя строка — требуемый вывод; все переходы оправданы.


\begin{learnbox}\textbf{Ключевой итог.} Дано $n\in\mathbb N$; основание $n>0$; вывод $n+2>2$.\par \textbf{Проверка доказательства.} 
Следите за условиями применимости правил и за переходом от произвольного объекта к общему выводу. 
\end{learnbox}

\subsection*{Строгое и нестрогое · № 24}\hypertarget{solution:MB-W01-L1-PR-A-024}{}

\muid{MB-W01-L1-PR-A-024}\par

\textbf{Условие.} Определите истинность $\forall x\in\mathbb R\colon x^2+1\ge1$.\par

\textbf{Полное решение.} Для любого вещественного $x$ выполняется $x^2\ge0$. Прибавляя $1$ к обеим частям, получаем $x^2+1\ge1$. При $x=0$ достигается равенство; потому знак $\ge$ здесь существенен.


\begin{learnbox}\textbf{Ключевой итог.} Истина; равенство при $x=0$.\par \textbf{Проверка доказательства.} 
Следите за условиями применимости правил и за переходом от произвольного объекта к общему выводу. 
\end{learnbox}

\subsection*{Составное число · № 25}\hypertarget{solution:MB-W01-L1-PR-A-025}{}

\muid{MB-W01-L1-PR-A-025}\par

\textbf{Условие.} Опровергните «каждое натуральное число больше единицы является простым».\par

\textbf{Полное решение.} Число $4$ натурально и больше единицы, то есть удовлетворяет предпосылкам. Однако $4=2\cdot2$ и имеет положительный делитель $2$, отличный от $1$ и $4$. Поэтому $4$ составное, а универсальное утверждение ложно.


\begin{learnbox}\textbf{Ключевой итог.} Ложь; контрпример $4$.\par \textbf{Проверка доказательства.} 
Следите за условиями применимости правил и за переходом от произвольного объекта к общему выводу. 
\end{learnbox}

\subsection*{Свидетель и проверка · № 26}\hypertarget{solution:MB-W01-L1-PR-A-026}{}

\muid{MB-W01-L1-PR-A-026}\par

\textbf{Условие.} Докажите $\exists k\in\mathbb Z\colon k+5=0$.\par

\textbf{Полное решение.} Уравнение $k+5=0$ имеет решение $k=-5$. Поскольку $-5\in\mathbb Z$ и $-5+5=0$, свидетель найден. Следовательно, квантор существования оправдан.


\begin{learnbox}\textbf{Ключевой итог.} Истина; $k=-5$.\par \textbf{Проверка доказательства.} 
Следите за условиями применимости правил и за переходом от произвольного объекта к общему выводу. 
\end{learnbox}

\subsection*{Порядок величин · № 27}\hypertarget{solution:MB-W01-L1-PR-A-027}{}

\muid{MB-W01-L1-PR-A-027}\par

\textbf{Условие.} Верно ли «если $x\ge4$, то $x\ge1$»?\par

\textbf{Полное решение.} При выполнении условия $x\ge4$ и числовом факте $4\ge1$ транзитивность даёт $x\ge1$. Вывод верен для каждого вещественного $x$, удовлетворяющего предпосылке. Вне этой предпосылки ничего не требуется.


\begin{learnbox}\textbf{Ключевой итог.} Истина; $x\ge4\ge1$.\par \textbf{Проверка доказательства.} 
Следите за условиями применимости правил и за переходом от произвольного объекта к общему выводу. 
\end{learnbox}

\subsection*{Ложная импликация · № 28}\hypertarget{solution:MB-W01-L1-PR-A-028}{}

\muid{MB-W01-L1-PR-A-028}\par

\textbf{Условие.} Опровергните «если целое $n$ кратно $4$, то оно кратно $8$».\par

\textbf{Полное решение.} Число $n=4$ целое и кратно $4$, поскольку $4=4\cdot1$. Однако оно не кратно $8$, так как $4/8=1/2$ не целое. Предпосылка истинна, заключение ложно; значит, импликация неверна.


\begin{learnbox}\textbf{Ключевой итог.} Ложь; $n=4$.\par \textbf{Проверка доказательства.} 
Числовой пример проверен на принадлежность области и на нарушение именно требуемого вывода. 
\end{learnbox}

\subsection*{Уточнение формулировки · № 29}\hypertarget{solution:MB-W01-L1-PR-A-029}{}

\muid{MB-W01-L1-PR-A-029}\par

\textbf{Условие.} Исправьте «каждое вещественное число имеет положительный квадрат», изменив одно слово.\par

\textbf{Полное решение.} Слово «положительный» заменяем на «неотрицательный». Получаем истинное высказывание: квадрат любого вещественного числа неотрицателен. Доказательство состоит в том, что произведение двух одинаковых чисел не бывает отрицательным; нуль даёт равенство.


\begin{learnbox}\textbf{Ключевой итог.} «Каждое вещественное число имеет неотрицательный квадрат».\par \textbf{Проверка доказательства.} 
Следите за условиями применимости правил и за переходом от произвольного объекта к общему выводу. 
\end{learnbox}

\subsection*{Множество решений · № 30}\hypertarget{solution:MB-W01-L1-PR-A-030}{}

\muid{MB-W01-L1-PR-A-030}\par

\textbf{Условие.} Найдите все целые $x$, удовлетворяющие $x^2=1$.\par

\textbf{Полное решение.} Уравнение равносильно $(x-1)(x+1)=0$. Значит, $x=1$ либо $x=-1$. Оба значения целые, и каждый квадрат равен единице. Благодаря свойству нулевого произведения других решений нет.


\begin{learnbox}\textbf{Ключевой итог.} $x\in\{-1,1\}$.\par \textbf{Проверка доказательства.} 
Помимо проверки найденных значений, аргумент исключает остальные варианты. 
\end{learnbox}

\section{Решения уровня B}

\subsection*{Прямое доказательство · № 1}\hypertarget{solution:MB-W01-L1-PR-B-031}{}

\muid{MB-W01-L1-PR-B-031}\par

\textbf{Условие.} Докажите $5(n+2)-5n=10$ для любого целого $n$.\par

\textbf{Полное решение.} Пусть $n\in\mathbb Z$ произвольно. Тогда $5(n+2)-5n=5n+10-5n=10$. Мы использовали распределительный закон и сокращение противоположных слагаемых. Результат не зависит от $n$, значит равенство является тождеством на целых.


\begin{learnbox}\textbf{Ключевой итог.} Тождество; результат всегда $10$.\par \textbf{Проверка доказательства.} 
Следите за условиями применимости правил и за переходом от произвольного объекта к общему выводу. 
\end{learnbox}

\subsection*{Область истинности · № 2}\hypertarget{solution:MB-W01-L1-PR-B-032}{}

\muid{MB-W01-L1-PR-B-032}\par

\textbf{Условие.} На какой из областей $\mathbb N$, $\mathbb Z$, $\mathbb R$ истинно $\forall x\colon x^2\ge x$?\par

\textbf{Полное решение.} Для натуральных $x\ge1$, поэтому $x(x-1)\ge0$. Для целых $x$ дополнительно возможны $x\le0$, когда оба множителя неположительны, а их произведение неотрицательно. Для вещественного $x=\tfrac12$ имеем $x^2=\tfrac14<\tfrac12=x$. Следовательно, формула универсально истинна на $\mathbb N$ и $\mathbb Z$, но ложна на $\mathbb R$.


\begin{learnbox}\textbf{Ключевой итог.} Истинно в $\mathbb N,\mathbb Z$; ложно в $\mathbb R$.\par \textbf{Проверка доказательства.} 
Следите за условиями применимости правил и за переходом от произвольного объекта к общему выводу. 
\end{learnbox}

\subsection*{Слабый вывод · № 3}\hypertarget{solution:MB-W01-L1-PR-B-033}{}

\muid{MB-W01-L1-PR-B-033}\par

\textbf{Условие.} Если $x>8$, докажите $2x>12$.\par

\textbf{Полное решение.} Пусть $x>8$. Так как $2>0$, умножение сохраняет знак: $2x>16$. Но $16>12$, значит по транзитивности $2x>12$. Все шаги работают для произвольного $x$, удовлетворяющего условию.


\begin{learnbox}\textbf{Ключевой итог.} Верно; $2x>16>12$.\par \textbf{Проверка доказательства.} 
Следите за условиями применимости правил и за переходом от произвольного объекта к общему выводу. 
\end{learnbox}

\subsection*{Граница условия · № 4}\hypertarget{solution:MB-W01-L1-PR-B-034}{}

\muid{MB-W01-L1-PR-B-034}\par

\textbf{Условие.} Опровергните «если $x^2\ge4$, то $x\ge2$» над вещественными числами.\par

\textbf{Полное решение.} При $x=-2$ квадрат равен $4$, поэтому условие $x^2\ge4$ выполнено. Однако $-2\ge2$ ложно. Следовательно, импликация неверна. Исправление: из $x^2\ge4$ следует $x\le-2$ или $x\ge2$.


\begin{learnbox}\textbf{Ключевой итог.} Ложь; $x=-2$.\par \textbf{Проверка доказательства.} 
Следите за условиями применимости правил и за переходом от произвольного объекта к общему выводу. 
\end{learnbox}

\subsection*{Целый множитель · № 5}\hypertarget{solution:MB-W01-L1-PR-B-035}{}

\muid{MB-W01-L1-PR-B-035}\par

\textbf{Условие.} Докажите, что $15k+25$ кратно $5$ при любом $k\in\mathbb Z$.\par

\textbf{Полное решение.} Для произвольного целого $k$ имеем $15k+25=5(3k+5)$. Число $3k+5$ целое. По определению делимости выражение $15k+25$ делится на $5$ для всех $k\in\mathbb Z$.


\begin{learnbox}\textbf{Ключевой итог.} Истина: $15k+25=5(3k+5)$.\par \textbf{Проверка доказательства.} 
Следите за условиями применимости правил и за переходом от произвольного объекта к общему выводу. 
\end{learnbox}

\subsection*{Неверное сокращение · № 6}\hypertarget{solution:MB-W01-L1-PR-B-036}{}

\muid{MB-W01-L1-PR-B-036}\par

\textbf{Условие.} Найдите ошибку в рассуждении: из $0\cdot a=0\cdot b$ следует $a=b$.\par

\textbf{Полное решение.} Равенство $0\cdot a=0\cdot b$ истинно при любых $a,b$, например при $a=1$, $b=2$. Деление на нулевой общий множитель запрещено. Поэтому заключение $a=b$ без дополнительных условий получить нельзя.


\begin{learnbox}\textbf{Ключевой итог.} Ошибка: сокращение на нуль.\par \textbf{Проверка доказательства.} 
Следите за условиями применимости правил и за переходом от произвольного объекта к общему выводу. 
\end{learnbox}

\subsection*{Идентичность · № 7}\hypertarget{solution:MB-W01-L1-PR-B-037}{}

\muid{MB-W01-L1-PR-B-037}\par

\textbf{Условие.} Докажите $(x+3)^2-(x-3)^2=12x$ для всех вещественных $x$.\par

\textbf{Полное решение.} Левая часть равна $(x^2+6x+9)-(x^2-6x+9)$. После раскрытия скобок получаем $x^2+6x+9-x^2+6x-9=12x$. Равенство получено формальными тождественными преобразованиями, верными для любого вещественного $x$.


\begin{learnbox}\textbf{Ключевой итог.} Тождество: $12x$.\par \textbf{Проверка доказательства.} 
Помимо проверки найденных значений, аргумент исключает остальные варианты. 
\end{learnbox}

\subsection*{Условная кратность · № 8}\hypertarget{solution:MB-W01-L1-PR-B-038}{}

\muid{MB-W01-L1-PR-B-038}\par

\textbf{Условие.} Докажите: если целое $n$ кратно $7$, то $3n$ кратно $7$.\par

\textbf{Полное решение.} Пусть $n=7k$ для некоторого целого $k$, что означает делимость $n$ на $7$. Тогда $3n=3\cdot7k=7(3k)$, а $3k\in\mathbb Z$. По определению делимости $3n$ кратно $7$.


\begin{learnbox}\textbf{Ключевой итог.} Истина; $3n=7(3k)$.\par \textbf{Проверка доказательства.} 
Следите за условиями применимости правил и за переходом от произвольного объекта к общему выводу. 
\end{learnbox}

\subsection*{Контрпример к квадрату · № 9}\hypertarget{solution:MB-W01-L1-PR-B-039}{}

\muid{MB-W01-L1-PR-B-039}\par

\textbf{Условие.} Опровергните «если $a^2>b^2$, то $a>b$» для вещественных $a,b$.\par

\textbf{Полное решение.} При $a=-3$, $b=2$ получаем $a^2=9>4=b^2$. Но $a=-3$ меньше $b=2$. Мы предъявили случай, удовлетворяющий предпосылке и нарушающий заключение, поэтому импликация ложна.


\begin{learnbox}\textbf{Ключевой итог.} Ложь; $(a,b)=(-3,2)$.\par \textbf{Проверка доказательства.} 
Числовой пример проверен на принадлежность области и на нарушение именно требуемого вывода. 
\end{learnbox}

\subsection*{Существование целого · № 10}\hypertarget{solution:MB-W01-L1-PR-B-040}{}

\muid{MB-W01-L1-PR-B-040}\par

\textbf{Условие.} Найдите целое число $m$, для которого $m^2-5m=0$, и докажите существование.\par

\textbf{Полное решение.} Уравнение равносильно $m(m-5)=0$, откуда $m=0$ или $m=5$. Оба числа целые и обращают исходное выражение в нуль. Уже число $m=0$ даёт свидетель существования; перечисление обоих решений дополнительно показывает полный набор.


\begin{learnbox}\textbf{Ключевой итог.} Существуют $m=0$ и $m=5$.\par \textbf{Проверка доказательства.} 
Следите за условиями применимости правил и за переходом от произвольного объекта к общему выводу. 
\end{learnbox}

\subsection*{Усиление тезиса · № 11}\hypertarget{solution:MB-W01-L1-PR-B-041}{}

\muid{MB-W01-L1-PR-B-041}\par

\textbf{Условие.} Докажите, что $x^2+4\ge4$ для любого вещественного $x$, и укажите равенство.\par

\textbf{Полное решение.} Для каждого $x\in\mathbb R$ имеем $x^2\ge0$. Прибавление $4$ даёт $x^2+4\ge4$. Равенство возможно тогда и только тогда, когда $x^2=0$, то есть $x=0$.


\begin{learnbox}\textbf{Ключевой итог.} Истина; равенство только при $x=0$.\par \textbf{Проверка доказательства.} 
Следите за условиями применимости правил и за переходом от произвольного объекта к общему выводу. 
\end{learnbox}

\subsection*{Сумма последовательных · № 12}\hypertarget{solution:MB-W01-L1-PR-B-042}{}

\muid{MB-W01-L1-PR-B-042}\par

\textbf{Условие.} Докажите, что $k+(k+1)+(k+2)+(k+3)=4k+6$ для всех целых $k$.\par

\textbf{Полное решение.} Складываем левые выражения: $k+k+1+k+2+k+3=4k+(1+2+3)=4k+6$. Равенство основано только на законах сложения, поэтому верно при любом целом $k$. Указание $k\in\mathbb Z$ здесь не ограничивает алгебраическую формулу, но соответствует условию о последовательных целых.


\begin{learnbox}\textbf{Ключевой итог.} Тождество; $4k+6$.\par \textbf{Проверка доказательства.} 
Помимо проверки найденных значений, аргумент исключает остальные варианты. 
\end{learnbox}

\subsection*{Неравенство с двумя условиями · № 13}\hypertarget{solution:MB-W01-L1-PR-B-043}{}

\muid{MB-W01-L1-PR-B-043}\par

\textbf{Условие.} Если $a\ge2$ и $b\ge3$, докажите $a+b\ge5$.\par

\textbf{Полное решение.} Из $a\ge2$ и $b\ge3$ сложением двух неравенств одинакового направления получаем $a+b\ge2+3=5$. Это прямое доказательство: использованы обе предпосылки. Равенство возможно при $a=2$ и $b=3$.


\begin{learnbox}\textbf{Ключевой итог.} Истина; $a+b\ge5$.\par \textbf{Проверка доказательства.} 
Следите за условиями применимости правил и за переходом от произвольного объекта к общему выводу. 
\end{learnbox}

\subsection*{Неполное доказательство · № 14}\hypertarget{solution:MB-W01-L1-PR-B-044}{}

\muid{MB-W01-L1-PR-B-044}\par

\textbf{Условие.} Оцените аргумент «$3\cdot4$ кратно $3$, значит произведение любых двух натуральных чисел кратно $3$».\par

\textbf{Полное решение.} Рассуждение неверно и как доказательство, и по выводу. Пара $(2,2)$ состоит из натуральных чисел, но произведение $4$ не делится на $3$. Значит, универсальный тезис опровергнут, а частный пример $3\cdot4$ ничего не гарантирует для других пар.


\begin{learnbox}\textbf{Ключевой итог.} Ложь; контрпример $2\cdot2=4$.\par \textbf{Проверка доказательства.} 
Следите за условиями применимости правил и за переходом от произвольного объекта к общему выводу. 
\end{learnbox}

\subsection*{Проверка обратного · № 15}\hypertarget{solution:MB-W01-L1-PR-B-045}{}

\muid{MB-W01-L1-PR-B-045}\par

\textbf{Условие.} Из $a=b$ следует $a^2=b^2$. Верно ли обратное для вещественных $a,b$?\par

\textbf{Полное решение.} Прямая стрелка верна: равные числа имеют равные квадраты. Обратная стрелка ложна: $a=1$, $b=-1$ дают $a^2=b^2=1$, хотя $a\ne b$. Следовательно, равенство квадратов означает $a=b$ или $a=-b$.


\begin{learnbox}\textbf{Ключевой итог.} Обратное ложно; $(1,-1)$.\par \textbf{Проверка доказательства.} 
Следите за условиями применимости правил и за переходом от произвольного объекта к общему выводу. 
\end{learnbox}

\subsection*{Отрицательный множитель · № 16}\hypertarget{solution:MB-W01-L1-PR-B-046}{}

\muid{MB-W01-L1-PR-B-046}\par

\textbf{Условие.} Если $a>b$, докажите $-2a<-2b$.\par

\textbf{Полное решение.} Из условия $a>b$ умножением обеих частей на отрицательное число $-2$ получаем $-2a<-2b$. Правило изменения направления является обязательным основанием. Утверждение верно для любой вещественной пары $a,b$ с $a>b$.


\begin{learnbox}\textbf{Ключевой итог.} Истина: $-2a<-2b$.\par \textbf{Проверка доказательства.} 
Следите за условиями применимости правил и за переходом от произвольного объекта к общему выводу. 
\end{learnbox}

\subsection*{Уравнение с двумя решениями · № 17}\hypertarget{solution:MB-W01-L1-PR-B-047}{}

\muid{MB-W01-L1-PR-B-047}\par

\textbf{Условие.} Найдите все вещественные $x$ с $(x-4)(x+1)=0$.\par

\textbf{Полное решение.} Если произведение равно нулю, то $x-4=0$ или $x+1=0$. Это даёт $x=4$ и $x=-1$. Подстановка каждого значения делает один множитель нулём. Других вещественных решений нет по свойству нулевого произведения.


\begin{learnbox}\textbf{Ключевой итог.} $x\in\{-1,4\}$.\par \textbf{Проверка доказательства.} 
Помимо проверки найденных значений, аргумент исключает остальные варианты. 
\end{learnbox}

\subsection*{Двойная квантификация · № 18}\hypertarget{solution:MB-W01-L1-PR-B-048}{}

\muid{MB-W01-L1-PR-B-048}\par

\textbf{Условие.} Докажите $\forall n\in\mathbb N\ \exists m\in\mathbb N\colon m=n+2$.\par

\textbf{Полное решение.} Пусть $n$ — произвольное натуральное число. Положим $m=n+2$. Поскольку сумма натурального и положительного целого натуральна, $m\in\mathbb N$. По построению выполнено требуемое равенство. Значит, после каждого $n$ можно выбрать подходящее $m$.


\begin{learnbox}\textbf{Ключевой итог.} Истина; $m=n+2$.\par \textbf{Проверка доказательства.} 
Следите за условиями применимости правил и за переходом от произвольного объекта к общему выводу. 
\end{learnbox}

\subsection*{Кванторы в обратном порядке · № 19}\hypertarget{solution:MB-W01-L1-PR-B-049}{}

\muid{MB-W01-L1-PR-B-049}\par

\textbf{Условие.} Опровергните $\exists m\in\mathbb N\ \forall n\in\mathbb N\colon m>n$.\par

\textbf{Полное решение.} Предположим, что существует фиксированное натуральное $m$, превосходящее любое натуральное $n$. Тогда разрешено выбрать $n=m+1$, и получится требование $m>m+1$, что ложно. Следовательно, универсально превосходящего натурального числа не существует.


\begin{learnbox}\textbf{Ключевой итог.} Ложь; для любого $m$ подходит $n=m+1$.\par \textbf{Проверка доказательства.} 
Отмечены зависимости переменных и порядок выбора свидетеля. 
\end{learnbox}

\subsection*{Логика «и» · № 20}\hypertarget{solution:MB-W01-L1-PR-B-050}{}

\muid{MB-W01-L1-PR-B-050}\par

\textbf{Условие.} Опровергните «любое натуральное число чётно и кратно $3$».\par

\textbf{Полное решение.} Число $1$ натуральное, но не чётное и не кратно $3$. Конъюнкция («и») требует обеих частей одновременно; при $n=1$ это условие не выполнено. Поэтому универсальное утверждение ложно.


\begin{learnbox}\textbf{Ключевой итог.} Ложь; $n=1$.\par \textbf{Проверка доказательства.} 
Следите за условиями применимости правил и за переходом от произвольного объекта к общему выводу. 
\end{learnbox}

\subsection*{Логика «или» · № 21}\hypertarget{solution:MB-W01-L1-PR-B-051}{}

\muid{MB-W01-L1-PR-B-051}\par

\textbf{Условие.} Докажите: для любого целого $n$ верно $n\ge0$ или $n<0$.\par

\textbf{Полное решение.} Для каждого целого $n$ по свойству линейного порядка выполняется ровно одно из двух: $n<0$ или $n\ge0$. Это исчерпывающее разбиение на случаи. Следовательно, дизъюнкция истинна без каких-либо дополнительных ограничений.


\begin{learnbox}\textbf{Ключевой итог.} Истина; два случая исчерпывают $\mathbb Z$.\par \textbf{Проверка доказательства.} 
Следите за условиями применимости правил и за переходом от произвольного объекта к общему выводу. 
\end{learnbox}

\subsection*{Деление и область · № 22}\hypertarget{solution:MB-W01-L1-PR-B-052}{}

\muid{MB-W01-L1-PR-B-052}\par

\textbf{Условие.} Можно ли из $ax=ay$ заключить $x=y$, если известно $a\ne0$?\par

\textbf{Полное решение.} Поскольку $a\ne0$, деление обеих частей $ax=ay$ на $a$ допустимо. Получаем $x=y$. Равенство доказано при заявленном условии; без него при $a=0$ вывод мог быть ложным.


\begin{learnbox}\textbf{Ключевой итог.} Да; делим на ненулевое $a$.\par \textbf{Проверка доказательства.} 
Следите за условиями применимости правил и за переходом от произвольного объекта к общему выводу. 
\end{learnbox}

\subsection*{Квадрат разности · № 23}\hypertarget{solution:MB-W01-L1-PR-B-053}{}

\muid{MB-W01-L1-PR-B-053}\par

\textbf{Условие.} Докажите $x^2-2x+1\ge0$ при любом $x\in\mathbb R$.\par

\textbf{Полное решение.} По формуле квадрата разности $x^2-2x+1=(x-1)^2$. Квадрат любого вещественного числа неотрицателен. Поэтому исходное выражение неотрицательно при всех $x$. Равенство имеет место только при $x=1$.


\begin{learnbox}\textbf{Ключевой итог.} Истина; $(x-1)^2\ge0$.\par \textbf{Проверка доказательства.} 
Следите за условиями применимости правил и за переходом от произвольного объекта к общему выводу. 
\end{learnbox}

\subsection*{Неверная универсальность · № 24}\hypertarget{solution:MB-W01-L1-PR-B-054}{}

\muid{MB-W01-L1-PR-B-054}\par

\textbf{Условие.} Опровергните «сумма двух положительных чисел меньше их произведения».\par

\textbf{Полное решение.} При $a=b=1$ оба числа положительны, но сумма равна $2$, а произведение равно $1$. Требуемое $2<1$ ложно. Предпосылки выполнены, следовательно, единичная пара опровергает универсальную формулировку.


\begin{learnbox}\textbf{Ключевой итог.} Ложь; $1+1=2>1\cdot1$.\par \textbf{Проверка доказательства.} 
Следите за условиями применимости правил и за переходом от произвольного объекта к общему выводу. 
\end{learnbox}

\subsection*{Связанные неравенства · № 25}\hypertarget{solution:MB-W01-L1-PR-B-055}{}

\muid{MB-W01-L1-PR-B-055}\par

\textbf{Условие.} Если $a\ge b$ и $b\ge0$, докажите $a\ge0$.\par

\textbf{Полное решение.} Обе предпосылки относятся к одним и тем же вещественным $a,b$. Их последовательное применение даёт $a\ge b\ge0$. По транзитивности $a\ge0$. Здесь нельзя опустить условие $b\ge0$: оно обеспечивает нижнюю границу.


\begin{learnbox}\textbf{Ключевой итог.} Истина; $a\ge b\ge0$.\par \textbf{Проверка доказательства.} 
Следите за условиями применимости правил и за переходом от произвольного объекта к общему выводу. 
\end{learnbox}

\subsection*{Чётный показатель · № 26}\hypertarget{solution:MB-W01-L1-PR-B-056}{}

\muid{MB-W01-L1-PR-B-056}\par

\textbf{Условие.} Докажите $x^4\ge0$ при любом вещественном $x$.\par

\textbf{Полное решение.} Число $x^2$ вещественно; его квадрат неотрицателен. Поскольку $x^4=(x^2)^2$, получаем $x^4\ge0$ для всех $x\in\mathbb R$. Равенство возможно только при $x=0$.


\begin{learnbox}\textbf{Ключевой итог.} Истина; $x^4=(x^2)^2\ge0$.\par \textbf{Проверка доказательства.} 
Следите за условиями применимости правил и за переходом от произвольного объекта к общему выводу. 
\end{learnbox}

\subsection*{Неравенство с параметром · № 27}\hypertarget{solution:MB-W01-L1-PR-B-057}{}

\muid{MB-W01-L1-PR-B-057}\par

\textbf{Условие.} Докажите: если $t\ge1$, то $t^2\ge t$.\par

\textbf{Полное решение.} При $t\ge1$ имеем $t\ge0$ и $t-1\ge0$. Поэтому $t(t-1)\ge0$. Раскрытие скобок даёт $t^2-t\ge0$, что равносильно $t^2\ge t$. Граница $t=1$ даёт равенство.


\begin{learnbox}\textbf{Ключевой итог.} Истина при $t\ge1$.\par \textbf{Проверка доказательства.} 
Следите за условиями применимости правил и за переходом от произвольного объекта к общему выводу. 
\end{learnbox}

\subsection*{Доказательство единственности · № 28}\hypertarget{solution:MB-W01-L1-PR-B-058}{}

\muid{MB-W01-L1-PR-B-058}\par

\textbf{Условие.} Докажите, что уравнение $x+7=3$ имеет ровно одно вещественное решение.\par

\textbf{Полное решение.} Из $x+7=3$ вычитанием $7$ следует $x=-4$. Обратно, $-4+7=3$, поэтому найденное число подходит. Любое решение обязано равняться $-4$, поскольку вычитание $7$ было допустимо и обратимо. Следовательно, решение существует и единственно.


\begin{learnbox}\textbf{Ключевой итог.} Единственное решение $x=-4$.\par \textbf{Проверка доказательства.} 
Следите за условиями применимости правил и за переходом от произвольного объекта к общему выводу. 
\end{learnbox}

\subsection*{Круговое рассуждение · № 29}\hypertarget{solution:MB-W01-L1-PR-B-059}{}

\muid{MB-W01-L1-PR-B-059}\par

\textbf{Условие.} Почему «$a+b=b+a$, потому что сумма не меняется при перестановке слагаемых» допустимо, если правило уже изучено, но недостаточно для доказательства самого правила?\par

\textbf{Полное решение.} Если переместительный закон сложения принят или уже доказан, ссылка на него является корректным основанием для данного равенства. Если же требуется доказать сам переместительный закон, фраза лишь повторяет утверждение другими словами и создаёт логический круг. Основание должно быть независимым от искомого тезиса.


\begin{learnbox}\textbf{Ключевой итог.} Допустимо как ссылка на известный закон; кругово как его доказательство.\par \textbf{Проверка доказательства.} 
Следите за условиями применимости правил и за переходом от произвольного объекта к общему выводу. 
\end{learnbox}

\subsection*{Проверка краёв · № 30}\hypertarget{solution:MB-W01-L1-PR-B-060}{}

\muid{MB-W01-L1-PR-B-060}\par

\textbf{Условие.} Покажите, почему утверждение «если $x>0$, то $x^2>0$» истинно, а замена условия на $x\ge0$ делает его ложным.\par

\textbf{Полное решение.} При $x>0$ оба множителя в $x\cdot x$ положительны, значит $x^2>0$. Если условие ослабить до $x\ge0$, допускается $x=0$, и тогда заключение $x^2>0$ ложно. Поэтому строгая граница в первоначальной формулировке существенна.


\begin{learnbox}\textbf{Ключевой итог.} Первая импликация истинна; вторая ложна при $x=0$.\par \textbf{Проверка доказательства.} 
Следите за условиями применимости правил и за переходом от произвольного объекта к общему выводу. 
\end{learnbox}

\section{Решения уровня C}

\subsection*{Смена области · № 1}\hypertarget{solution:MB-W01-L1-PR-C-061}{}

\muid{MB-W01-L1-PR-C-061}\par

\textbf{Условие.} Докажите $n(n-2)\ge0$ для целых $n\le0$ или $n\ge2$; найдите целое исключение для всех $\mathbb Z$.\par

\textbf{Полное решение.} Если $n\le0$, то оба множителя $n$ и $n-2$ неположительны, их произведение неотрицательно. Если $n\ge2$, оба множителя неотрицательны. Однако для $n=1$ получаем $1(1-2)=-1<0$. Поэтому доказанный вывод верен на объединении указанных диапазонов и ложен на всех целых.


\begin{learnbox}\textbf{Ключевой итог.} Верно на $n\le0$ или $n\ge2$; исключение $n=1$.\par \textbf{Проверка доказательства.} 
Помимо проверки найденных значений, аргумент исключает остальные варианты. 
\end{learnbox}

\subsection*{Квадратная оценка · № 2}\hypertarget{solution:MB-W01-L1-PR-C-062}{}

\muid{MB-W01-L1-PR-C-062}\par

\textbf{Условие.} Докажите $x^2+6x+10\ge1$ для каждого вещественного $x$ и найдите случай равенства.\par

\textbf{Полное решение.} Преобразуем выражение: $x^2+6x+10=(x+3)^2+1$. Так как квадрат неотрицателен, сумма не меньше $1$. Равенство достигается только при $x+3=0$, то есть $x=-3$. Область вещественных значений допускает этот случай.


\begin{learnbox}\textbf{Ключевой итог.} Минимум $1$ при $x=-3$.\par \textbf{Проверка доказательства.} 
Следите за условиями применимости правил и за переходом от произвольного объекта к общему выводу. 
\end{learnbox}

\subsection*{Ошибка в направлении · № 3}\hypertarget{solution:MB-W01-L1-PR-C-063}{}

\muid{MB-W01-L1-PR-C-063}\par

\textbf{Условие.} Рецензируйте «если $n$ кратно $12$, то кратно $4$; $20$ кратно $4$; значит, $20$ кратно $12$».\par

\textbf{Полное решение.} Правильно: из $12\mid n$ следует $4\mid n$, поскольку $12=3\cdot4$. Рассуждение применяет обратную стрелку без доказательства. Число $20=4\cdot5$ кратно $4$, но не кратно $12$. Значит, итог ложен; дефект — утверждение следствия как достаточного условия.


\begin{learnbox}\textbf{Ключевой итог.} Ошибка: незаконная обратная импликация.\par \textbf{Проверка доказательства.} 
Следите за условиями применимости правил и за переходом от произвольного объекта к общему выводу. 
\end{learnbox}

\subsection*{Три части множества · № 4}\hypertarget{solution:MB-W01-L1-PR-C-064}{}

\muid{MB-W01-L1-PR-C-064}\par

\textbf{Условие.} Для $x\in\mathbb R$ докажите, что ровно одна ситуация $x<-2$, $x=-2$, $x>-2$ реализуется.\par

\textbf{Полное решение.} По свойству линейного порядка любые два вещественных числа сравнимы. Для $x$ и $-2$ имеет место хотя бы одно из трёх отношений $<$, $=$, $>$. Они попарно несовместимы: число не может одновременно быть меньше и равно $-2$ и т.д. Следовательно, реализуется ровно один из трёх случаев.


\begin{learnbox}\textbf{Ключевой итог.} Три случая образуют полное разбиение.\par \textbf{Проверка доказательства.} 
Следите за условиями применимости правил и за переходом от произвольного объекта к общему выводу. 
\end{learnbox}

\subsection*{Равенство двух квадратов · № 5}\hypertarget{solution:MB-W01-L1-PR-C-065}{}

\muid{MB-W01-L1-PR-C-065}\par

\textbf{Условие.} Найдите все пары вещественных $a,b$, для которых $a^2=b^2$, в виде условия на $a,b$.\par

\textbf{Полное решение.} Равенство $a^2=b^2$ равносильно $a^2-b^2=0$, то есть $(a-b)(a+b)=0$. По свойству нулевого произведения $a-b=0$ или $a+b=0$. Это означает $a=b$ либо $a=-b$. Обратно, в каждом случае квадраты равны; значит, описание полное.


\begin{learnbox}\textbf{Ключевой итог.} $a=b$ или $a=-b$.\par \textbf{Проверка доказательства.} 
Помимо проверки найденных значений, аргумент исключает остальные варианты. 
\end{learnbox}

\subsection*{Двойной контрпример · № 6}\hypertarget{solution:MB-W01-L1-PR-C-066}{}

\muid{MB-W01-L1-PR-C-066}\par

\textbf{Условие.} Опровергните «если два вещественных числа имеют одинаковые модули, то они равны»; запишите верное исправление.\par

\textbf{Полное решение.} При $a=3$, $b=-3$ выполнено $|a|=|b|=3$, но $a\ne b$. Универсальная импликация ложна. Правильное условие: из $|a|=|b|$ следует $a=b$ или $a=-b$; это эквивалентно $a^2=b^2$ и доказывается разложением разности квадратов.


\begin{learnbox}\textbf{Ключевой итог.} Контрпример $(3,-3)$; исправление $a=\pm b$.\par \textbf{Проверка доказательства.} 
Числовой пример проверен на принадлежность области и на нарушение именно требуемого вывода. 
\end{learnbox}

\subsection*{Квантор и зависимость · № 7}\hypertarget{solution:MB-W01-L1-PR-C-067}{}

\muid{MB-W01-L1-PR-C-067}\par

\textbf{Условие.} Докажите $\forall a\in\mathbb Z\ \exists b\in\mathbb Z\colon a+b=0$ и объясните, почему свидетель зависит от $a$.\par

\textbf{Полное решение.} Пусть $a$ произвольное целое. Выберем $b=-a$; противоположное целому число тоже целое. Тогда $a+b=a-a=0$. Для разных $a$ требуется, вообще говоря, разное $b$; это допускается порядком кванторов, когда существование стоит после универсальности.


\begin{learnbox}\textbf{Ключевой итог.} Истина; для данного $a$ берём $b=-a$.\par \textbf{Проверка доказательства.} 
Отмечены зависимости переменных и порядок выбора свидетеля. 
\end{learnbox}

\subsection*{Ложная перестановка · № 8}\hypertarget{solution:MB-W01-L1-PR-C-068}{}

\muid{MB-W01-L1-PR-C-068}\par

\textbf{Условие.} Опровергните $\exists b\in\mathbb Z\ \forall a\in\mathbb Z\colon a+b=0$.\par

\textbf{Полное решение.} Если один фиксированный $b$ удовлетворяет равенству при всех $a$, то при $a=0$ получаем $b=0$. При $a=1$ тогда $1+b=1\ne0$, противоречие. Поэтому одного целого $b$, компенсирующего все целые $a$, не существует.


\begin{learnbox}\textbf{Ключевой итог.} Ложь: $a=0$ требует $b=0$, а $a=1$ — $b=-1$.\par \textbf{Проверка доказательства.} 
Числовой пример проверен на принадлежность области и на нарушение именно требуемого вывода. 
\end{learnbox}

\subsection*{Неравенство по условию · № 9}\hypertarget{solution:MB-W01-L1-PR-C-069}{}

\muid{MB-W01-L1-PR-C-069}\par

\textbf{Условие.} Если $u>v\ge0$, докажите $u^2>v^2$ и объясните, почему $u\ge0$ следует из условия.\par

\textbf{Полное решение.} Из $u>v\ge0$ сразу получаем $u>0$ и $u-v>0$. Поэтому $u+v>0$. По формуле разности квадратов $u^2-v^2=(u-v)(u+v)>0$, следовательно, $u^2>v^2$. Все предпосылки использованы при обосновании знаков множителей.


\begin{learnbox}\textbf{Ключевой итог.} Истина; $(u-v)(u+v)>0$.\par \textbf{Проверка доказательства.} 
Следите за условиями применимости правил и за переходом от произвольного объекта к общему выводу. 
\end{learnbox}

\subsection*{Эквивалентность границ · № 10}\hypertarget{solution:MB-W01-L1-PR-C-070}{}

\muid{MB-W01-L1-PR-C-070}\par

\textbf{Условие.} Докажите: для вещественного $x$ условие $|x|\le3$ равносильно $-3\le x\le3$.\par

\textbf{Полное решение.} Если $|x|\le3$, то $-\lvert x\rvert\le x\le\lvert x\rvert$ даёт $-3\le x\le3$. Обратно, при $x\ge0$ из $x\le3$ следует $|x|=x\le3$; при $x<0$ из $x\ge-3$ следует $|x|=-x\le3$. Оба направления доказаны, значит условия равносильны.


\begin{learnbox}\textbf{Ключевой итог.} $|x|\le3\Leftrightarrow-3\le x\le3$.\par \textbf{Проверка доказательства.} 
Следите за условиями применимости правил и за переходом от произвольного объекта к общему выводу. 
\end{learnbox}

\subsection*{Разбор ошибки с корнем · № 11}\hypertarget{solution:MB-W01-L1-PR-C-071}{}

\muid{MB-W01-L1-PR-C-071}\par

\textbf{Условие.} Найдите недостающее условие в выводе «из $x^2=9$ следует $x=3$».\par

\textbf{Полное решение.} Уравнение $x^2=9$ равносильно $(x-3)(x+3)=0$, значит $x=3$ либо $x=-3$. Для заключения $x=3$ нужно дополнительно предположить $x\ge0$. При таком условии отрицательное решение исключается, а положительное проверяется подстановкой.


\begin{learnbox}\textbf{Ключевой итог.} Нужно $x\ge0$; без него $x=\pm3$.\par \textbf{Проверка доказательства.} 
Следите за условиями применимости правил и за переходом от произвольного объекта к общему выводу. 
\end{learnbox}

\subsection*{Усреднение · № 12}\hypertarget{solution:MB-W01-L1-PR-C-072}{}

\muid{MB-W01-L1-PR-C-072}\par

\textbf{Условие.} Докажите: если $a\ge b$, то $a\ge\frac{a+b}{2}\ge b$ для вещественных $a,b$.\par

\textbf{Полное решение.} Поскольку $a-b\ge0$, имеем $a-\frac{a+b}{2}=\frac{a-b}{2}\ge0$, то есть $a\ge\frac{a+b}{2}$. Аналогично $\frac{a+b}{2}-b=\frac{a-b}{2}\ge0$. Оба неравенства следуют из одного условия, значит среднее лежит между числами.


\begin{learnbox}\textbf{Ключевой итог.} $a\ge(a+b)/2\ge b$.\par \textbf{Проверка доказательства.} 
Следите за условиями применимости правил и за переходом от произвольного объекта к общему выводу. 
\end{learnbox}

\subsection*{Единственность результата · № 13}\hypertarget{solution:MB-W01-L1-PR-C-073}{}

\muid{MB-W01-L1-PR-C-073}\par

\textbf{Условие.} Докажите, что уравнение $3x-5=7$ имеет единственное вещественное решение.\par

\textbf{Полное решение.} Прибавляя $5$, получаем $3x=12$. Деление на ненулевое $3$ даёт $x=4$. Обратно, $3\cdot4-5=7$, следовательно, найденное значение действительно решение. Поскольку каждое преобразование равносильно, любой другой кандидат обязан равняться $4$.


\begin{learnbox}\textbf{Ключевой итог.} $x=4$, единственное решение.\par \textbf{Проверка доказательства.} 
Помимо проверки найденных значений, аргумент исключает остальные варианты. 
\end{learnbox}

\subsection*{Невозможность · № 14}\hypertarget{solution:MB-W01-L1-PR-C-074}{}

\muid{MB-W01-L1-PR-C-074}\par

\textbf{Условие.} Докажите, что не существуют вещественные $x,y$ с $x^2+y^2=-1$.\par

\textbf{Полное решение.} Для любых вещественных $x,y$ выполняются $x^2\ge0$ и $y^2\ge0$. Складывая, получаем $x^2+y^2\ge0$. Поэтому эта сумма не может равняться $-1<0$. Следовательно, искомых вещественных пар не существует.


\begin{learnbox}\textbf{Ключевой итог.} Невозможно: $x^2+y^2\ge0$.\par \textbf{Проверка доказательства.} 
Следите за условиями применимости правил и за переходом от произвольного объекта к общему выводу. 
\end{learnbox}

\subsection*{Контрпример к нестрогому · № 15}\hypertarget{solution:MB-W01-L1-PR-C-075}{}

\muid{MB-W01-L1-PR-C-075}\par

\textbf{Условие.} Опровергните «если $a>b$, то $a^2\ge b^2$» без ограничений на знаки.\par

\textbf{Полное решение.} Выберем $a=-1$, $b=-2$. Тогда $a>b$ истинно, но $a^2=1<4=b^2$, так что заключение $a^2\ge b^2$ ложно. Дефект формулировки связан с тем, что возведение в квадрат не сохраняет порядок отрицательных чисел.


\begin{learnbox}\textbf{Ключевой итог.} Ложь; $(-1,-2)$.\par \textbf{Проверка доказательства.} 
Числовой пример проверен на принадлежность области и на нарушение именно требуемого вывода. 
\end{learnbox}

\subsection*{Разбиение условия · № 16}\hypertarget{solution:MB-W01-L1-PR-C-076}{}

\muid{MB-W01-L1-PR-C-076}\par

\textbf{Условие.} Найдите все целые $n$, при которых $n^2\le4$, и докажите полноту списка.\par

\textbf{Полное решение.} Неравенство $n^2\le4$ для вещественных $n$ равносильно $|n|\le2$, то есть $-2\le n\le2$. Целые числа внутри отрезка: $-2,-1,0,1,2$. Каждый квадрат равен соответственно $4,1,0,1,4$ и не превосходит $4$; вне отрезка квадрат больше $4$. Список исчерпывающий.


\begin{learnbox}\textbf{Ключевой итог.} $n\in\{-2,-1,0,1,2\}$.\par \textbf{Проверка доказательства.} 
Помимо проверки найденных значений, аргумент исключает остальные варианты. 
\end{learnbox}

\subsection*{Сравнение долей · № 17}\hypertarget{solution:MB-W01-L1-PR-C-077}{}

\muid{MB-W01-L1-PR-C-077}\par

\textbf{Условие.} Докажите: если $a>b>0$, то $\frac1a<\frac1b$.\par

\textbf{Полное решение.} По условию $a-b>0$, а произведение $ab>0$. Значит, $\frac1b-\frac1a=\frac{a-b}{ab}>0$. Отсюда $\frac1b>\frac1a$. Строгое требование положительности знаменателей используется при проверке знака дроби.


\begin{learnbox}\textbf{Ключевой итог.} Истина: $1/a<1/b$.\par \textbf{Проверка доказательства.} 
Следите за условиями применимости правил и за переходом от произвольного объекта к общему выводу. 
\end{learnbox}

\subsection*{Строгость доказательства · № 18}\hypertarget{solution:MB-W01-L1-PR-C-078}{}

\muid{MB-W01-L1-PR-C-078}\par

\textbf{Условие.} Прочитайте «если $a=b$, то $a^3=b^3$» и напишите обоснование, применимое ко всем вещественным $a,b$.\par

\textbf{Полное решение.} Если $a=b$, то замена каждого вхождения $a$ равным ему $b$ в произведении допустима: $a^3=a\cdot a\cdot a=b\cdot b\cdot b=b^3$. Это свойство равенства; конкретных значений не требуется. Следовательно, импликация истинна для всех вещественных $a,b$.


\begin{learnbox}\textbf{Ключевой итог.} Истина: подстановка равных величин сохраняет куб.\par \textbf{Проверка доказательства.} 
Помимо проверки найденных значений, аргумент исключает остальные варианты. 
\end{learnbox}

\subsection*{Анализ ложного доказательства · № 19}\hypertarget{solution:MB-W01-L1-PR-C-079}{}

\muid{MB-W01-L1-PR-C-079}\par

\textbf{Условие.} Ученик пишет: «Для любого $x$ имеем $x^2+2x+1>0$, потому что это квадрат». Найдите ошибку и исправьте тезис.\par

\textbf{Полное решение.} Выражение действительно равно $(x+1)^2$, но квадрат гарантирует лишь неотрицательность. При $x=-1$ имеем $(x+1)^2=0$, поэтому знак $>0$ ложен. Корректная формулировка: $x^2+2x+1\ge0$ при всех вещественных $x$, с равенством при $x=-1$.


\begin{learnbox}\textbf{Ключевой итог.} Ошибка в строгости: верно $\ge0$, равенство при $x=-1$.\par \textbf{Проверка доказательства.} 
Следите за условиями применимости правил и за переходом от произвольного объекта к общему выводу. 
\end{learnbox}

\subsection*{Оценка по двум условиям · № 20}\hypertarget{solution:MB-W01-L1-PR-C-080}{}

\muid{MB-W01-L1-PR-C-080}\par

\textbf{Условие.} Если $0\le a\le2$ и $0\le b\le3$, докажите $0\le ab\le6$.\par

\textbf{Полное решение.} Из $a,b\ge0$ имеем $ab\ge0$. Умножив $a\le2$ на неотрицательное $b$, получим $ab\le2b$. Из $b\le3$ следует $2b\le6$. По транзитивности $0\le ab\le6$. Равенство справа возможно при $a=2$, $b=3$.


\begin{learnbox}\textbf{Ключевой итог.} $0\le ab\le6$.\par \textbf{Проверка доказательства.} 
Следите за условиями применимости правил и за переходом от произвольного объекта к общему выводу. 
\end{learnbox}

\section{Решения уровня D}

\subsection*{Две границы · № 1}\hypertarget{solution:MB-W01-L1-PR-D-081}{}

\muid{MB-W01-L1-PR-D-081}\par

\textbf{Условие.} Для $x\in\mathbb R$ докажите эквивалентность: $x^2-4x+3\le0$ тогда и только тогда, когда $1\le x\le3$.\par

\textbf{Полное решение.} Раскладываем $x^2-4x+3=(x-1)(x-3)$. Если $1\le x\le3$, то $x-1\ge0$, а $x-3\le0$, поэтому произведение неположительно. Обратно, вне отрезка оба множителя имеют одинаковый знак: при $x<1$ оба отрицательны, при $x>3$ оба положительны. Тогда произведение строго положительно. Следовательно, неравенство выполняется ровно на отрезке $[1,3]$.


\begin{learnbox}\textbf{Ключевой итог.} Равносильно $x\in[1,3]$.\par \textbf{Проверка доказательства.} 
Следите за условиями применимости правил и за переходом от произвольного объекта к общему выводу. 
\end{learnbox}

\subsection*{Зависимость кванторов · № 2}\hypertarget{solution:MB-W01-L1-PR-D-082}{}

\muid{MB-W01-L1-PR-D-082}\par

\textbf{Условие.} Сравните истинность $\forall x\in\mathbb R\ \exists y\in\mathbb R\colon x+y=1$ и $\exists y\in\mathbb R\ \forall x\in\mathbb R\colon x+y=1$.\par

\textbf{Полное решение.} Первое высказывание истинно: для каждого заданного $x$ выбираем $y=1-x$, тогда $x+y=1$. Второе ложно: из подстановки $x=0$ получаем $y=1$, а из $x=1$ получаем $y=0$. Один фиксированный $y$ не может быть одновременно $1$ и $0$. Значение истинности меняется при перестановке кванторов.


\begin{learnbox}\textbf{Ключевой итог.} Первое истинно ($y=1-x$); второе ложно.\par \textbf{Проверка доказательства.} 
Отмечены зависимости переменных и порядок выбора свидетеля. 
\end{learnbox}

\subsection*{Несколько условий · № 3}\hypertarget{solution:MB-W01-L1-PR-D-083}{}

\muid{MB-W01-L1-PR-D-083}\par

\textbf{Условие.} Если $a,b,c\ge0$ и $a+b+c=0$, докажите $a=b=c=0$.\par

\textbf{Полное решение.} Из $a,b,c\ge0$ следует $a\le a+b+c=0$, поэтому $a=0$ с учётом $a\ge0$. Аналогично $b\le a+b+c=0$ и $c\le a+b+c=0$, откуда $b=c=0$. Обратно, нулевые значения действительно удовлетворяют условию. Важна неотрицательность каждого слагаемого: при отрицательных числах сумма могла бы обращаться в нуль иначе.


\begin{learnbox}\textbf{Ключевой итог.} $a=b=c=0$.\par \textbf{Проверка доказательства.} 
Следите за условиями применимости правил и за переходом от произвольного объекта к общему выводу. 
\end{learnbox}

\subsection*{Разбор ложного обобщения · № 4}\hypertarget{solution:MB-W01-L1-PR-D-084}{}

\muid{MB-W01-L1-PR-D-084}\par

\textbf{Условие.} Числа $1^2+1+17$, $2^2+2+17$ и $3^2+3+17$ просты. Следует ли, что $n^2+n+17$ всегда простое при $n\in\mathbb N$?\par

\textbf{Полное решение.} Наблюдаемые первые значения не дают доказательства универсальности. Более того, при $n=17$ выражение равно $17^2+17+17=17(17+2)=17\cdot19$, что составное число. Число $17$ принадлежит $\mathbb N$ и удовлетворяет всем условиям, поэтому найден полноценный контрпример. Данный опыт демонстрирует, почему длинная серия простых значений не гарантирует формулу для всех $n$.


\begin{learnbox}\textbf{Ключевой итог.} Нет; при $n=17$ значение $17\cdot19$ составное.\par \textbf{Проверка доказательства.} 
Помимо проверки найденных значений, аргумент исключает остальные варианты. 
\end{learnbox}

\subsection*{Геометрический контрпример · № 5}\hypertarget{solution:MB-W01-L1-PR-D-085}{}

\muid{MB-W01-L1-PR-D-085}\par

\textbf{Условие.} Опровергните «если у четырёхугольника диагонали равны, то он прямоугольник» на примере равнобедренной трапеции.\par

\textbf{Полное решение.} Возьмём $A=(-2,0)$, $B=(2,0)$, $C=(1,2)$, $D=(-1,2)$. Основания $AB$ и $CD$ горизонтальны, их длины $4$ и $2$, значит фигура является трапецией с неравными основаниями, а не прямоугольником. Диагонали $AC$ и $BD$ обе имеют длину $\sqrt{3^2+2^2}=\sqrt{13}$. Условия о равных диагоналях выполнены, заключение о прямоугольнике ложно.


\begin{learnbox}\textbf{Ключевой итог.} Ложь; симметричная трапеция с диагоналями $\sqrt{13}$.\par \textbf{Проверка доказательства.} 
Числовой пример проверен на принадлежность области и на нарушение именно требуемого вывода. 
\end{learnbox}

\subsection*{Невозможность с условиями · № 6}\hypertarget{solution:MB-W01-L1-PR-D-086}{}

\muid{MB-W01-L1-PR-D-086}\par

\textbf{Условие.} Докажите, что не существует вещественного $x$, для которого одновременно $x>2$ и $x^2<4$.\par

\textbf{Полное решение.} Предположим, что $x>2$ и $x^2<4$. Тогда $x>0$, поэтому $x^2-4=(x-2)(x+2)>0$ как произведение двух положительных чисел. Значит, $x^2>4$, что несовместимо с $x^2<4$. Такой вещественный $x$ не существует. Более короткое обоснование использует монотонность квадрата на положительной полуоси.


\begin{learnbox}\textbf{Ключевой итог.} Невозможно: из $x>2$ следует $x^2>4$.\par \textbf{Проверка доказательства.} 
Следите за условиями применимости правил и за переходом от произвольного объекта к общему выводу. 
\end{learnbox}

\subsection*{Уточнение области · № 7}\hypertarget{solution:MB-W01-L1-PR-D-087}{}

\muid{MB-W01-L1-PR-D-087}\par

\textbf{Условие.} Исследуйте истинность $\forall x\in A\colon x^2\ge2x$ для $A=\mathbb N$ и для $A=\{2,3,4,\ldots\}$; приведите верную максимальную целую область.\par

\textbf{Полное решение.} Неравенство равносильно $x(x-2)\ge0$, то есть $x\le0$ или $x\ge2$. На $\mathbb N$ оно ложно при $x=1$. На $\{2,3,4,\ldots\}$ истинно, поскольку оба множителя неотрицательны. Среди целых максимальная область, где оно выполнено, есть $\{\ldots,-2,-1,0\}\cup\{2,3,\ldots\}$. При отрицательном $x$ оба множителя отрицательны и произведение положительно.


\begin{learnbox}\textbf{Ключевой итог.} Ложно на $\mathbb N$; истинно при $x\in\mathbb Z$, $x\le0$ или $x\ge2$.\par \textbf{Проверка доказательства.} 
Следите за условиями применимости правил и за переходом от произвольного объекта к общему выводу. 
\end{learnbox}

\subsection*{Правка доказательства · № 8}\hypertarget{solution:MB-W01-L1-PR-D-088}{}

\muid{MB-W01-L1-PR-D-088}\par

\textbf{Условие.} Ученик получил из $ab=ac$ равенство $b=c$. Уточните необходимую предпосылку и обоснуйте её необходимость примером.\par

\textbf{Полное решение.} Если $a\ne0$, из $ab=ac$ делением на $a$ действительно получаем $b=c$. Если $a=0$, равенство произведений верно при любых $b,c$; например, $b=1$, $c=2$ дают $0\cdot1=0\cdot2$, хотя $b\ne c$. Следовательно, предпосылка $a\ne0$ необходима для вывода при произвольных $b,c$.


\begin{learnbox}\textbf{Ключевой итог.} Нужно $a\ne0$; при $a=0$, $b=1$, $c=2$ вывод неверен.\par \textbf{Проверка доказательства.} 
Следите за условиями применимости правил и за переходом от произвольного объекта к общему выводу. 
\end{learnbox}

\subsection*{Синтез ограничений · № 9}\hypertarget{solution:MB-W01-L1-PR-D-089}{}

\muid{MB-W01-L1-PR-D-089}\par

\textbf{Условие.} Докажите, что при $a,b\in\mathbb R$ и $a+b=6$ выполняется $ab\le9$, и найдите условие равенства.\par

\textbf{Полное решение.} Для вещественных $a,b$ имеем $(a-b)^2\ge0$. По тождеству $(a-b)^2=(a+b)^2-4ab=36-4ab$. Значит, $36-4ab\ge0$, откуда $ab\le9$. Равенство наступает тогда и только тогда, когда $a-b=0$, то есть $a=b=3$. Подстановка подтверждает выполнение обоих условий.


\begin{learnbox}\textbf{Ключевой итог.} $ab\le9$; равенство при $a=b=3$.\par \textbf{Проверка доказательства.} 
Следите за условиями применимости правил и за переходом от произвольного объекта к общему выводу. 
\end{learnbox}

\subsection*{Полное доказательство · № 10}\hypertarget{solution:MB-W01-L1-PR-D-090}{}

\muid{MB-W01-L1-PR-D-090}\par

\textbf{Условие.} Найдите все вещественные решения системы $x+y=4$, $x^2+y^2=8$ и строго докажите полноту ответа.\par

\textbf{Полное решение.} Из $(x+y)^2=x^2+2xy+y^2$ получаем $16=8+2xy$, поэтому $xy=4$. Теперь $(x-y)^2=x^2-2xy+y^2=8-8=0$, следовательно, $x=y$. С учётом $x+y=4$ получаем $x=y=2$. Проверка: $2+2=4$, $2^2+2^2=8$. Все преобразования являются необходимыми следствиями системы, а проверка подтверждает существование решения; других пар нет.


\begin{learnbox}\textbf{Ключевой итог.} Единственная пара $(x,y)=(2,2)$.\par \textbf{Проверка доказательства.} 
Помимо проверки найденных значений, аргумент исключает остальные варианты. 
\end{learnbox}

\chapter{Эталонные ответы к вопросам по ходу чтения}
Все вопросы на проверку понимания, переносы из разобранных примеров и мини-пробы из эвристик имеют ниже проверяемые эталоны. Их назначение — проверка аргумента после самостоятельной попытки. Самооценка допускает индивидуальный ответ; для неё приведён образец структуры.

\section{Вопросы после теоретических глав}

\noindent\muid{MB-W01-L1-TH-01}\par \textbf{Вопрос.} Объясните своими словами, какая часть проверки $1+2=3$, $2+3=5$, $3+4=7$ является наблюдением, а какая формула даёт доказательство для любого $n$.\par 
\textbf{Образец решения.} Проверки $n=1,2,3$ относятся лишь к трём числам. Равенство $n+(n+1)=2n+1$ получается объединением слагаемых при произвольном $n$, поэтому даёт общее доказательство.\medskip

\noindent\muid{MB-W01-L1-TH-02}\par \textbf{Вопрос.} Для каждого выражения $x^2=9$, «все квадраты неотрицательны», $2+2=4$ определите, является ли оно высказыванием в записанном виде.\par 
\textbf{Образец решения.} Открытая формула $x^2=9$ имеет свободную переменную. Предложение «все квадраты неотрицательны» содержит универсальность по подразумеваемой области; $2+2=4$ — истинное числовое высказывание.\medskip

\noindent\muid{MB-W01-L1-TH-03}\par \textbf{Вопрос.} Сравните $\forall n\in\mathbb N\ \exists m\in\mathbb N\colon m>n$ и $\exists m\in\mathbb N\ \forall n\in\mathbb N\colon m>n$.\par 
\textbf{Образец решения.} Первое утверждение истинно: для каждого $n$ берём $m=n+1$. Второе ложно: для любого фиксированного $m$ число $n=m+1$ нарушает $m>n$.\medskip

\noindent\muid{MB-W01-L1-TH-04}\par \textbf{Вопрос.} Приведите пример неверного универсального утверждения и объясните, почему ваш объект действительно удовлетворяет всем условиям.\par 
\textbf{Образец решения.} Например, «каждый вещественный $x$ имеет $x^2>0$» опровергается допустимым $x=0$: квадрат равен нулю. Важно назвать область и проверить нарушение строгого знака.\medskip

\noindent\muid{MB-W01-L1-TH-05}\par \textbf{Вопрос.} Как отличить полный перебор от недостаточной проверки отдельных случаев?\par 
\textbf{Образец решения.} В полном переборе исчерпывающе проверено всё конечное множество возможностей. Три подстановки при переменной из бесконечного множества представляют только выборку.\medskip

\noindent\muid{MB-W01-L1-TH-06}\par \textbf{Вопрос.} В каком месте цепочки $a=b\Rightarrow a+c=b+c\Rightarrow 2(a+c)=2(b+c)$ используется правило сохранения равенства?\par 
\textbf{Образец решения.} Переход $a=b\Rightarrow a+c=b+c$ основан на прибавлении одинакового $c$ к обеим частям; следующий переход — на умножении двух равных чисел на $2$.\medskip

\noindent\muid{MB-W01-L1-TH-07}\par \textbf{Вопрос.} Почему из $a(x-2)=0$ нельзя без уточнения сделать вывод $x=2$?\par 
\textbf{Образец решения.} Из $a(x-2)=0$ можно заключить $x=2$ только при $a\ne0$; при $a=0$ любое $x$ подходит.\medskip

\noindent\muid{MB-W01-L1-TH-08}\par \textbf{Вопрос.} Запишите утверждение, обратное к «если $x>7$, то $x>0$», и проверьте его на числе $x=1$.\par 
\textbf{Образец решения.} Обратное: «если $x>0$, то $x>7$». Оно ложно при $x=1$: предпосылка верна, заключение нарушено.\medskip

\noindent\muid{MB-W01-L1-TH-09}\par \textbf{Вопрос.} Объясните, в чём состоит разница между доказательством кратности для всех $n$ и проверкой только $n=1$.\par 
\textbf{Образец решения.} Формула $12n+18=6(2n+3)$ при $n\in\mathbb Z$ предъявляет целый множитель $2n+3$ и потому доказывает кратность $6$ универсально. Проверка $n=1$ этого не делает.\medskip

\noindent\muid{MB-W01-L1-TH-10}\par \textbf{Вопрос.} Почему рисунок с двумя похожими углами сам по себе не доказывает их равенство?\par 
\textbf{Образец решения.} Масштаб рисунка зависит от способа изображения и может быть неточным. Равенство углов нужно вывести из условия, определения или теоремы.\medskip

\noindent\muid{MB-W01-L1-TH-11}\par \textbf{Вопрос.} Найдите незаконный переход в записи $(x-1)(x+2)=0\Rightarrow x-1=0$ и объясните, какой случай потерян.\par 
\textbf{Образец решения.} Из $(x-1)(x+2)=0$ следует $x=1$ или $x=-2$; запись с одним множителем теряет корень $-2$.\medskip

\noindent\muid{MB-W01-L1-TH-12}\par \textbf{Вопрос.} Составьте мини-чек-лист из пяти вопросов к чужому доказательству; объясните выбор каждого.\par 
\textbf{Образец решения.} Образец: указана ли область? правильно ли расставлены кванторы? обоснован ли каждый переход? не выполнено ли деление на нуль? соответствует ли итог требуемому тезису?\medskip

\noindent\muid{MB-W01-L1-TH-13}\par \textbf{Вопрос.} Запишите два навыка, которые уже получаются, и два, которые требуют повторения.\par 
\textbf{Образец решения.} Индивидуальный образец самооценки: «умею читать квантор и проверять область; тренирую построение строгой общей цепочки и доказательство единственности». Содержание самооценки зависит от ученика.\medskip

\section{Переносы после 30 разобранных примеров}

\noindent\muid{MB-W01-L1-TR-01}\par \textbf{Вопрос.} Составьте аналогичное ложное числовое высказывание.\par 
\textbf{Образец решения.} Например, $7+5=13$ — ложное числовое высказывание, поскольку сумма равна $12$.\medskip

\noindent\muid{MB-W01-L1-TR-02}\par \textbf{Вопрос.} Добавьте квантор, чтобы получить истинное утверждение.\par 
\textbf{Образец решения.} Истинное высказывание: $\exists x\in\mathbb R\colon 2x+3=7$. Свидетель $x=2$.\medskip

\noindent\muid{MB-W01-L1-TR-03}\par \textbf{Вопрос.} Проверьте, остаётся ли оно истинным при $n\in\mathbb Z$.\par 
\textbf{Образец решения.} На целых неверно: при $n=-2$ получаем $3n+2=-4<0$.\medskip

\noindent\muid{MB-W01-L1-TR-04}\par \textbf{Вопрос.} Где в доказательстве используется целочисленность $k$?\par 
\textbf{Образец решения.} Целочисленность $k$ гарантирует, что $2k$ и $2(k+1)$ действительно кратны двум по определению.\medskip

\noindent\muid{MB-W01-L1-TR-05}\par \textbf{Вопрос.} Почему важно назвать случай равенства?\par 
\textbf{Образец решения.} При $x=0$ квадрат равен нулю, поэтому строгий знак $>0$ не выполняется.\medskip

\noindent\muid{MB-W01-L1-TR-06}\par \textbf{Вопрос.} Что изменится, если требуется $n\in\mathbb Z$?\par 
\textbf{Образец решения.} В целых числах существуют два свидетеля: $n=7$ и $n=-7$; в натуральных подходит только $7$.\medskip

\noindent\muid{MB-W01-L1-TR-07}\par \textbf{Вопрос.} Запишите аналог для $n-1<n$.\par 
\textbf{Образец решения.} Из $1>0$ прибавлением $n-1$ к обеим частям получаем $n>n-1$, то есть $n-1<n$.\medskip

\noindent\muid{MB-W01-L1-TR-08}\par \textbf{Вопрос.} Каков ответ при ограничении $x\in\mathbb N$?\par 
\textbf{Образец решения.} При $x\in\mathbb N$ уравнение $x^2=4$ имеет только решение $2$; $-2$ исключён областью.\medskip

\noindent\muid{MB-W01-L1-TR-09}\par \textbf{Вопрос.} Укажите все случаи равенства внутри $\mathbb N$.\par 
\textbf{Образец решения.} $n^2=n$ даёт $n(n-1)=0$, а среди натуральных выполняется только $n=1$.\medskip

\noindent\muid{MB-W01-L1-TR-10}\par \textbf{Вопрос.} Как изменится ответ для $2x=8$?\par 
\textbf{Образец решения.} Для $2x=8$ есть целое решение $x=4$, так что существование доказано.\medskip

\noindent\muid{MB-W01-L1-TR-11}\par \textbf{Вопрос.} Переформулируйте утверждение, исключив найденное исключение.\par 
\textbf{Образец решения.} Верная формулировка: «каждое простое число, отличное от $2$, нечётно». Если простое чётно, оно делится на $2$ и потому равно $2$.\medskip

\noindent\muid{MB-W01-L1-TR-12}\par \textbf{Вопрос.} Придумайте другую верную цепочку из двух неравенств.\par 
\textbf{Образец решения.} Например, $x>10$ вместе с $10>5$ даёт $x>5$.\medskip

\noindent\muid{MB-W01-L1-TR-13}\par \textbf{Вопрос.} Как выглядит аналогичное равенство для множителя $-3$?\par 
\textbf{Образец решения.} $-3(a+b)=-3a-3b$ по распределительному закону.\medskip

\noindent\muid{MB-W01-L1-TR-14}\par \textbf{Вопрос.} Сформулируйте похожее следствие при умножении на $c$.\par 
\textbf{Образец решения.} Если $a=b$, то умножение на любое вещественное $c$ даёт $ac=bc$.\medskip

\noindent\muid{MB-W01-L1-TR-15}\par \textbf{Вопрос.} Проверьте случай $n=-2$.\par 
\textbf{Образец решения.} При $n=-2$ имеем $n(n+1)=2$, а $n^2=4$, поэтому строгое $2>4$ ложно.\medskip

\noindent\muid{MB-W01-L1-TR-16}\par \textbf{Вопрос.} Почему при $n=\tfrac14$ такой вывод потребовал бы иной проверки?\par 
\textbf{Образец решения.} При $n=\tfrac14$ выражение $12n+18=21$, что не кратно $6$; множитель $2n+3=\tfrac72$ не целый.\medskip

\noindent\muid{MB-W01-L1-TR-17}\par \textbf{Вопрос.} Поясните отдельно случай $c<0$.\par 
\textbf{Образец решения.} Даже если $c<0$, прибавление одного и того же $c$ сохраняет порядок. Например, $3>1\Rightarrow3-10>1-10$.\medskip

\noindent\muid{MB-W01-L1-TR-18}\par \textbf{Вопрос.} Найдите допустимые $a,b$, для которых выполняется равенство.\par 
\textbf{Образец решения.} Например, $a=1,b=0$ удовлетворяют условиям и дают $a+b=a$.\medskip

\noindent\muid{MB-W01-L1-TR-19}\par \textbf{Вопрос.} Верно ли обратное утверждение?\par 
\textbf{Образец решения.} Обратное верно: если $x=0$, то $x^2=0$. В этом примере утверждения равносильны.\medskip

\noindent\muid{MB-W01-L1-TR-20}\par \textbf{Вопрос.} Выведите аналогичную формулу для четырёх последовательных чисел.\par 
\textbf{Образец решения.} Сумма $n+(n+1)+(n+2)+(n+3)=4n+6$.\medskip

\noindent\muid{MB-W01-L1-TR-21}\par \textbf{Вопрос.} Сравните области, в которых утверждение верно.\par 
\textbf{Образец решения.} На целых произведение $n(n+1)$ всегда неотрицательно; на вещественных контрпример $n=-\tfrac12$ показывает отрицательное значение.\medskip

\noindent\muid{MB-W01-L1-TR-22}\par \textbf{Вопрос.} Докажите также обратное направление.\par 
\textbf{Образец решения.} Если $-1<x<1$, то $|x|<1$, поэтому $x^2<1$. Это обратная стрелка и вместе с исходной образует эквивалентность.\medskip

\noindent\muid{MB-W01-L1-TR-23}\par \textbf{Вопрос.} Испытайте $a=-1$, $b=-2$.\par 
\textbf{Образец решения.} При $a=-1,b=-2$ имеем $a>b$, но $a^2=1<4=b^2$; условие $b\ge0$ существенно.\medskip

\noindent\muid{MB-W01-L1-TR-24}\par \textbf{Вопрос.} Какая строка меняется для строгих знаков?\par 
\textbf{Образец решения.} Если $a>b$ и $b>c$, то $a>c$ по транзитивности строгого порядка.\medskip

\noindent\muid{MB-W01-L1-TR-25}\par \textbf{Вопрос.} Почему нельзя применить тот же ход к $a=-1$?\par 
\textbf{Образец решения.} $\sqrt a$ как вещественное выражение не определено при $a=-1$; рассуждение выходит за заданную область.\medskip

\noindent\muid{MB-W01-L1-TR-26}\par \textbf{Вопрос.} Как изменится список для кубов?\par 
\textbf{Образец решения.} При делении куба целого числа на $4$ возможны остатки $0$, $1$ и $3$: чётные кубы дают $0$, нечётные при остатках $1$ и $3$ сохраняют соответствующие остатки.\medskip

\noindent\muid{MB-W01-L1-TR-27}\par \textbf{Вопрос.} Где точно применено требование $x\ne0$?\par 
\textbf{Образец решения.} Условие $x\ne0$ делает деление $xy=0$ на $x$ допустимым; при $x=0$ заключение «$x=0$ или $y=0$» уже доказано.\medskip

\noindent\muid{MB-W01-L1-TR-28}\par \textbf{Вопрос.} Составьте собственный пример на неравенствах.\par 
\textbf{Образец решения.} Например, из $x>3$ следует $x^2>9$ при $x>3$, а контрапозиция утверждает: из $x^2\le9$ следует $x\le3$. Обратная стрелка $x^2>9\Rightarrow x>3$ ложна при $x=-4$.\medskip

\noindent\muid{MB-W01-L1-TR-29}\par \textbf{Вопрос.} Какова самая короткая форма записи этого доказательства?\par 
\textbf{Образец решения.} Среди трёх последовательных чисел одно кратно $3$ и одно чётное; следовательно, произведение делится на взаимно простые $3$ и $2$, то есть на $6$.\medskip

\noindent\muid{MB-W01-L1-TR-30}\par \textbf{Вопрос.} Сколько первых подстановок может скрывать ошибочное обобщение?\par 
\textbf{Образец решения.} Формула даёт простые значения для $n=1,\ldots,39$; уже при $n=40$ она равна $41^2$, поэтому множество первых удачных проверок не заменяет доказательства.\medskip

\section{Проверки эвристик}

\noindent\muid{MB-W01-L1-HR-01}\par \textbf{Вопрос.} Докажите $(n+1)-n=1$ для каждого целого $n$.\par 
\textbf{Образец решения.} $(n+1)-n=1$ по сокращению противоположных слагаемых при произвольном целом $n$.\medskip

\noindent\muid{MB-W01-L1-HR-02}\par \textbf{Вопрос.} Докажите существование целого корня $x^2=36$.\par 
\textbf{Образец решения.} Положим $x=6\in\mathbb Z$: $x^2=36$, существование доказано.\medskip

\noindent\muid{MB-W01-L1-HR-03}\par \textbf{Вопрос.} Опровергните «каждое вещественное число имеет ненулевой квадрат».\par 
\textbf{Образец решения.} $x=0$ допустим и даёт $x^2=0$, значит универсальная формулировка с «ненулевым» ложна.\medskip

\noindent\muid{MB-W01-L1-HR-04}\par \textbf{Вопрос.} Найдите все решения $x+4=9$.\par 
\textbf{Образец решения.} Из $x+4=9$ следует $x=5$; обратная подстановка верна, другие решения невозможны после вычитания $4$.\medskip

\noindent\muid{MB-W01-L1-HR-05}\par \textbf{Вопрос.} Истинно ли «если $x>6$, то $x>2$»?\par 
\textbf{Образец решения.} Верно: $x>6$ и $6>2$ дают $x>2$.\medskip

\noindent\muid{MB-W01-L1-HR-06}\par \textbf{Вопрос.} Проверьте $(a-2)^2=a^2-4a+4$.\par 
\textbf{Образец решения.} $(a-2)^2=a^2-4a+4$ при всех вещественных $a$ по формуле квадрата разности.\medskip

\noindent\muid{MB-W01-L1-HR-07}\par \textbf{Вопрос.} Докажите $x^2+2x+1\ge0$.\par 
\textbf{Образец решения.} $x^2+2x+1=(x+1)^2\ge0$; равенство при $x=-1$.\medskip

\noindent\muid{MB-W01-L1-HR-08}\par \textbf{Вопрос.} Докажите $|x|\ge0$ при $x<0$, $x=0$, $x>0$.\par 
\textbf{Образец решения.} $|x|=-x>0$ при $x<0$, $|0|=0$, $|x|=x>0$ при $x>0$; случаи исчерпывающие.\medskip

\noindent\muid{MB-W01-L1-HR-09}\par \textbf{Вопрос.} Придумайте трапецию, которая не является прямоугольником.\par 
\textbf{Образец решения.} Трапеция с вершинами $(-2,0),(2,0),(1,2),(-1,2)$ имеет одну пару параллельных сторон с разной длиной и потому не является прямоугольником.\medskip

\noindent\muid{MB-W01-L1-HR-10}\par \textbf{Вопрос.} Почему из $x^2=4$ нельзя заключить только $x=2$?\par 
\textbf{Образец решения.} Равенство $x^2=4$ имеет два решения $x=2$ и $x=-2$; потерян второй корень.\medskip

\noindent\muid{MB-W01-L1-HR-11}\par \textbf{Вопрос.} Преобразуйте $(n+2)+(n+3)$.\par 
\textbf{Образец решения.} $(n+2)+(n+3)=2n+5$ по объединению подобных слагаемых.\medskip

\section{Промежуточные контрольные вопросы}

\noindent\muid{MB-W01-L1-CP-01}\par \textbf{Вопрос.} Назовите контрпример к «все вещественные $x$ имеют $x^2>0$» и свидетель к «существует целое $x$ с $x^2=9$». \par 
\textbf{Образец решения.} Контрпример к $x^2>0$ — $x=0$; свидетель к $x^2=9$ над целыми — $x=3$ (также $-3$).\medskip

\noindent\muid{MB-W01-L1-CP-02}\par \textbf{Вопрос.} Почему проверка $n=1,2,3$ не доказывает свойство для всех натуральных $n$?\par 
\textbf{Образец решения.} Проверены лишь три допустимых значения из бесконечной области. Универсальное доказательство должно быть применимо к произвольному натуральному $n$.\medskip

\noindent\muid{MB-W01-L1-CP-03}\par \textbf{Вопрос.} Почему $x^2=4\Rightarrow x=2$ неверно без дополнительного условия?\par 
\textbf{Образец решения.} $x^2=4$ означает $(x-2)(x+2)=0$, то есть $x=2$ или $x=-2$; один вариант без дополнительного $x\ge0$ не оправдан.\medskip

\noindent\muid{MB-W01-L1-CP-04}\par \textbf{Вопрос.} При каких условиях из $ax=ay$ следует $x=y$?\par 
\textbf{Образец решения.} Нужно $a\ne0$. Только тогда деление обеих частей $ax=ay$ на $a$ допустимо и даёт $x=y$.\medskip

\noindent\muid{MB-W01-L1-CP-05}\par \textbf{Вопрос.} Проверьте, что полное решение содержит и проверку, и объяснение отсутствия других вариантов.\par 
\textbf{Образец решения.} Для «всех решений» доказываем необходимость полученного списка и проверяем достаточность каждого кандидата. Для «докажите» каждый ключевой переход должен иметь основание.\medskip

\chapter{Таблица ответов и итоговая самооценка}
\begin{keybox}Таблица может использоваться независимо от полного решения. Для задач на доказательство указан тезис и его ключевое основание; для задач с отрицательным ответом приведён допустимый контрпример. Все 90 строк соответствуют практике один к одному.\end{keybox}
\small
\begin{longtable}{p{.30\textwidth}p{.62\textwidth}}\toprule Идентификатор & Ответ / ключевая идея\\\midrule\endfirsthead\toprule Идентификатор & Ответ / ключевая идея\\\midrule\endhead

\muid{MB-W01-L1-PR-A-001} & Истина: $14-6=8$.\\

\muid{MB-W01-L1-PR-A-002} & Ложь: $8\cdot7=56\ne54$.\\

\muid{MB-W01-L1-PR-A-003} & Открытая формула; нужен выбор $x$ или квантор.\\

\muid{MB-W01-L1-PR-A-004} & Истина для каждого $x\in\mathbb R$.\\

\muid{MB-W01-L1-PR-A-005} & Свидетель $n=5$.\\

\muid{MB-W01-L1-PR-A-006} & Нет; единственное числовое решение $n=-1$.\\

\muid{MB-W01-L1-PR-A-007} & Нельзя: $-3\notin\mathbb N$.\\

\muid{MB-W01-L1-PR-A-008} & Истина: $x^2\ge0$ при всех $x\in\mathbb R$.\\

\muid{MB-W01-L1-PR-A-009} & Ложь; контрпример $x=0$.\\

\muid{MB-W01-L1-PR-A-010} & Истина; разность равна $1>0$.\\

\muid{MB-W01-L1-PR-A-011} & Истина; например $x=4$.\\

\muid{MB-W01-L1-PR-A-012} & Истина на указанном пятиэлементном множестве.\\

\muid{MB-W01-L1-PR-A-013} & Обратное ложно; $x=1$.\\

\muid{MB-W01-L1-PR-A-014} & Нет; это инструкция к решению.\\

\muid{MB-W01-L1-PR-A-015} & Истина: разность сторон равна $1$.\\

\muid{MB-W01-L1-PR-A-016} & Аргумент недостаточен: проверены два значения.\\

\muid{MB-W01-L1-PR-A-017} & Истина; $x=0$.\\

\muid{MB-W01-L1-PR-A-018} & Да; $n=5$.\\

\muid{MB-W01-L1-PR-A-019} & Истина по свойству нуля при сложении.\\

\muid{MB-W01-L1-PR-A-020} & Ложь; $n=0$.\\

\muid{MB-W01-L1-PR-A-021} & Истина: обе части тождественны.\\

\muid{MB-W01-L1-PR-A-022} & Ложь; при $(1,1)$ имеем $4\ne2$.\\

\muid{MB-W01-L1-PR-A-023} & Дано $n\in\mathbb N$; основание $n>0$; вывод $n+2>2$.\\

\muid{MB-W01-L1-PR-A-024} & Истина; равенство при $x=0$.\\

\muid{MB-W01-L1-PR-A-025} & Ложь; контрпример $4$.\\

\muid{MB-W01-L1-PR-A-026} & Истина; $k=-5$.\\

\muid{MB-W01-L1-PR-A-027} & Истина; $x\ge4\ge1$.\\

\muid{MB-W01-L1-PR-A-028} & Ложь; $n=4$.\\

\muid{MB-W01-L1-PR-A-029} & «Каждое вещественное число имеет неотрицательный квадрат».\\

\muid{MB-W01-L1-PR-A-030} & $x\in\{-1,1\}$.\\

\muid{MB-W01-L1-PR-B-031} & Тождество; результат всегда $10$.\\

\muid{MB-W01-L1-PR-B-032} & Истинно в $\mathbb N,\mathbb Z$; ложно в $\mathbb R$.\\

\muid{MB-W01-L1-PR-B-033} & Верно; $2x>16>12$.\\

\muid{MB-W01-L1-PR-B-034} & Ложь; $x=-2$.\\

\muid{MB-W01-L1-PR-B-035} & Истина: $15k+25=5(3k+5)$.\\

\muid{MB-W01-L1-PR-B-036} & Ошибка: сокращение на нуль.\\

\muid{MB-W01-L1-PR-B-037} & Тождество: $12x$.\\

\muid{MB-W01-L1-PR-B-038} & Истина; $3n=7(3k)$.\\

\muid{MB-W01-L1-PR-B-039} & Ложь; $(a,b)=(-3,2)$.\\

\muid{MB-W01-L1-PR-B-040} & Существуют $m=0$ и $m=5$.\\

\muid{MB-W01-L1-PR-B-041} & Истина; равенство только при $x=0$.\\

\muid{MB-W01-L1-PR-B-042} & Тождество; $4k+6$.\\

\muid{MB-W01-L1-PR-B-043} & Истина; $a+b\ge5$.\\

\muid{MB-W01-L1-PR-B-044} & Ложь; контрпример $2\cdot2=4$.\\

\muid{MB-W01-L1-PR-B-045} & Обратное ложно; $(1,-1)$.\\

\muid{MB-W01-L1-PR-B-046} & Истина: $-2a<-2b$.\\

\muid{MB-W01-L1-PR-B-047} & $x\in\{-1,4\}$.\\

\muid{MB-W01-L1-PR-B-048} & Истина; $m=n+2$.\\

\muid{MB-W01-L1-PR-B-049} & Ложь; для любого $m$ подходит $n=m+1$.\\

\muid{MB-W01-L1-PR-B-050} & Ложь; $n=1$.\\

\muid{MB-W01-L1-PR-B-051} & Истина; два случая исчерпывают $\mathbb Z$.\\

\muid{MB-W01-L1-PR-B-052} & Да; делим на ненулевое $a$.\\

\muid{MB-W01-L1-PR-B-053} & Истина; $(x-1)^2\ge0$.\\

\muid{MB-W01-L1-PR-B-054} & Ложь; $1+1=2>1\cdot1$.\\

\muid{MB-W01-L1-PR-B-055} & Истина; $a\ge b\ge0$.\\

\muid{MB-W01-L1-PR-B-056} & Истина; $x^4=(x^2)^2\ge0$.\\

\muid{MB-W01-L1-PR-B-057} & Истина при $t\ge1$.\\

\muid{MB-W01-L1-PR-B-058} & Единственное решение $x=-4$.\\

\muid{MB-W01-L1-PR-B-059} & Допустимо как ссылка на известный закон; кругово как его доказательство.\\

\muid{MB-W01-L1-PR-B-060} & Первая импликация истинна; вторая ложна при $x=0$.\\

\muid{MB-W01-L1-PR-C-061} & Верно на $n\le0$ или $n\ge2$; исключение $n=1$.\\

\muid{MB-W01-L1-PR-C-062} & Минимум $1$ при $x=-3$.\\

\muid{MB-W01-L1-PR-C-063} & Ошибка: незаконная обратная импликация.\\

\muid{MB-W01-L1-PR-C-064} & Три случая образуют полное разбиение.\\

\muid{MB-W01-L1-PR-C-065} & $a=b$ или $a=-b$.\\

\muid{MB-W01-L1-PR-C-066} & Контрпример $(3,-3)$; исправление $a=\pm b$.\\

\muid{MB-W01-L1-PR-C-067} & Истина; для данного $a$ берём $b=-a$.\\

\muid{MB-W01-L1-PR-C-068} & Ложь: $a=0$ требует $b=0$, а $a=1$ — $b=-1$.\\

\muid{MB-W01-L1-PR-C-069} & Истина; $(u-v)(u+v)>0$.\\

\muid{MB-W01-L1-PR-C-070} & $|x|\le3\Leftrightarrow-3\le x\le3$.\\

\muid{MB-W01-L1-PR-C-071} & Нужно $x\ge0$; без него $x=\pm3$.\\

\muid{MB-W01-L1-PR-C-072} & $a\ge(a+b)/2\ge b$.\\

\muid{MB-W01-L1-PR-C-073} & $x=4$, единственное решение.\\

\muid{MB-W01-L1-PR-C-074} & Невозможно: $x^2+y^2\ge0$.\\

\muid{MB-W01-L1-PR-C-075} & Ложь; $(-1,-2)$.\\

\muid{MB-W01-L1-PR-C-076} & $n\in\{-2,-1,0,1,2\}$.\\

\muid{MB-W01-L1-PR-C-077} & Истина: $1/a<1/b$.\\

\muid{MB-W01-L1-PR-C-078} & Истина: подстановка равных величин сохраняет куб.\\

\muid{MB-W01-L1-PR-C-079} & Ошибка в строгости: верно $\ge0$, равенство при $x=-1$.\\

\muid{MB-W01-L1-PR-C-080} & $0\le ab\le6$.\\

\muid{MB-W01-L1-PR-D-081} & Равносильно $x\in[1,3]$.\\

\muid{MB-W01-L1-PR-D-082} & Первое истинно ($y=1-x$); второе ложно.\\

\muid{MB-W01-L1-PR-D-083} & $a=b=c=0$.\\

\muid{MB-W01-L1-PR-D-084} & Нет; при $n=17$ значение $17\cdot19$ составное.\\

\muid{MB-W01-L1-PR-D-085} & Ложь; симметричная трапеция с диагоналями $\sqrt{13}$.\\

\muid{MB-W01-L1-PR-D-086} & Невозможно: из $x>2$ следует $x^2>4$.\\

\muid{MB-W01-L1-PR-D-087} & Ложно на $\mathbb N$; истинно при $x\in\mathbb Z$, $x\le0$ или $x\ge2$.\\

\muid{MB-W01-L1-PR-D-088} & Нужно $a\ne0$; при $a=0$, $b=1$, $c=2$ вывод неверен.\\

\muid{MB-W01-L1-PR-D-089} & $ab\le9$; равенство при $a=b=3$.\\

\muid{MB-W01-L1-PR-D-090} & Единственная пара $(x,y)=(2,2)$.\\

\bottomrule\end{longtable}\normalsize
\section{Личная итоговая карта}
\begin{longtable}{p{.62\textwidth}p{.24\textwidth}}\toprule Компетенция & Самооценка\\\midrule\endhead
Различаю высказывание и открытую формулу & умею / тренирую\\
Указываю область переменной и кванторы & умею / тренирую\\
Не путаю пример и доказательство для всех & умею / тренирую\\
Предъявляю корректный контрпример & умею / тренирую\\
Строю цепь обоснованных переходов & умею / тренирую\\
Проверяю обратное направление и крайние случаи & умею / тренирую\\
Пишу доказательство без скрытого деления на нуль & умею / тренирую\\
Могу восстановить решение задачи уровня B без подсказки & умею / тренирую\\\bottomrule\end{longtable}
\section{Итоговый мини-тур и оценивание}
Рекомендуемая работа на 90 минут: B-037 (10 минут), B-052 (15 минут), C-073 (20 минут), C-078 (20 минут) и D-090 (25 минут). Оценивайте каждое решение по рубрике 0–4 балла, максимум 20. Критерий допуска к следующему уроку: не менее 14 баллов, при этом в двух задачах должны присутствовать полные строгие доказательства и ни в одной не должно быть неотмеченного запрещённого деления.

\chapter{Источники, происхождение и юридические пояснения}
Все условия примеров и самостоятельных задач, а также решения и рисунки этого издания подготовлены специально для Math Blast. Внешние материалы использованы как ориентиры для подбора методики и последующего чтения; текст и рисунки этих изданий не перепечатывались. Применимость условий чужих ресурсов к публичному изданию следует проверять независимо.
\begin{enumerate}

\item \textbf{План Math Blast}. ArtemLevin / math-blast, index.html, DATA.weeks[0].\par \url{https://github.com/ArtemLevin/math-blast/blob/main/index.html}\par Основной источник последовательности занятий; тема и домашняя работа недели 1.

\item \textbf{А. Я. Канель-Белов, А. К. Ковальджи}. «Как решают нестандартные задачи», под ред. В. О. Бугаенко, М.: МЦНМО, 2008, 96 с..\par \url{https://math.ru/lib/460}\par Для дальнейшего чтения: разделы о методах, доказательствах и оформлении решений.

\item \textbf{В. А. Успенский}. «Простейшие примеры математических доказательств», М.: МЦНМО, 2009, 56 с..\par \url{https://math.ru/lib/mmmf/34}\par Дополнительная перспектива на понятие доказательства и аксиоматический подход.

\item \textbf{МЦНМО / Problems.ru}. База задач и правила использования материалов.\par \url{https://problems.ru/about_system.php}\par Источники задач для самостоятельного подбора. Условия и решения из базы в пособие не переносились.

\item \textbf{Art of Problem Solving}. How to Write a Math Solution.\par \url{https://artofproblemsolving.com/blog/articles/how-to-write-a-solution}\par Советы по структуре олимпиадного письменного решения; англоязычный ресурс.

\item \textbf{Всероссийская олимпиада школьников}. Официальный раздел «Математика».\par \url{https://vserosolimp.edsoo.ru/matematika}\par Регламенты, задания и критерии; актуальность конкретного года сверять отдельно.

\end{enumerate}
\section{Редакционная трассировка}
План: \texttt{ArtemLevin/math-blast}, \texttt{index.html}, \texttt{DATA.weeks[0].lesson1}. Фиксированный commit плана: \texttt{277841af9f279459929c1f5137e92ec96d5d1e3e}. ID издания: \texttt{MB-W01-L1}; ID занятия: 1. Каждая задача, подсказка, решение и строка ответа привязаны к стабильному идентификатору \texttt{MB-W01-L1-PR-...}. Проверка статуса издания выполняется по manifest.json и отчёту QA.\par
\textit{Подготовлено как кандидат для математического и полиграфического ревью.}\par
\end{document}
